% Options for packages loaded elsewhere
\PassOptionsToPackage{unicode}{hyperref}
\PassOptionsToPackage{hyphens}{url}
\documentclass[
  11pt,
]{report}
\usepackage{xcolor}
\usepackage[margin=1in]{geometry}
\usepackage{amsmath,amssymb}
\setcounter{secnumdepth}{-\maxdimen} % remove section numbering
\usepackage{iftex}
\ifPDFTeX
  \usepackage[T1]{fontenc}
  \usepackage[utf8]{inputenc}
  \usepackage{textcomp} % provide euro and other symbols
\else % if luatex or xetex
  \usepackage{unicode-math} % this also loads fontspec
  \defaultfontfeatures{Scale=MatchLowercase}
  \defaultfontfeatures[\rmfamily]{Ligatures=TeX,Scale=1}
\fi
\usepackage{lmodern}
\ifPDFTeX\else
  % xetex/luatex font selection
\fi
% Use upquote if available, for straight quotes in verbatim environments
\IfFileExists{upquote.sty}{\usepackage{upquote}}{}
\IfFileExists{microtype.sty}{% use microtype if available
  \usepackage[]{microtype}
  \UseMicrotypeSet[protrusion]{basicmath} % disable protrusion for tt fonts
}{}
\makeatletter
\@ifundefined{KOMAClassName}{% if non-KOMA class
  \IfFileExists{parskip.sty}{%
    \usepackage{parskip}
  }{% else
    \setlength{\parindent}{0pt}
    \setlength{\parskip}{6pt plus 2pt minus 1pt}}
}{% if KOMA class
  \KOMAoptions{parskip=half}}
\makeatother
\usepackage{color}
\usepackage{fancyvrb}
\newcommand{\VerbBar}{|}
\newcommand{\VERB}{\Verb[commandchars=\\\{\}]}
\DefineVerbatimEnvironment{Highlighting}{Verbatim}{commandchars=\\\{\}}
% Add ',fontsize=\small' for more characters per line
\newenvironment{Shaded}{}{}
\newcommand{\AlertTok}[1]{\textcolor[rgb]{1.00,0.00,0.00}{\textbf{#1}}}
\newcommand{\AnnotationTok}[1]{\textcolor[rgb]{0.38,0.63,0.69}{\textbf{\textit{#1}}}}
\newcommand{\AttributeTok}[1]{\textcolor[rgb]{0.49,0.56,0.16}{#1}}
\newcommand{\BaseNTok}[1]{\textcolor[rgb]{0.25,0.63,0.44}{#1}}
\newcommand{\BuiltInTok}[1]{\textcolor[rgb]{0.00,0.50,0.00}{#1}}
\newcommand{\CharTok}[1]{\textcolor[rgb]{0.25,0.44,0.63}{#1}}
\newcommand{\CommentTok}[1]{\textcolor[rgb]{0.38,0.63,0.69}{\textit{#1}}}
\newcommand{\CommentVarTok}[1]{\textcolor[rgb]{0.38,0.63,0.69}{\textbf{\textit{#1}}}}
\newcommand{\ConstantTok}[1]{\textcolor[rgb]{0.53,0.00,0.00}{#1}}
\newcommand{\ControlFlowTok}[1]{\textcolor[rgb]{0.00,0.44,0.13}{\textbf{#1}}}
\newcommand{\DataTypeTok}[1]{\textcolor[rgb]{0.56,0.13,0.00}{#1}}
\newcommand{\DecValTok}[1]{\textcolor[rgb]{0.25,0.63,0.44}{#1}}
\newcommand{\DocumentationTok}[1]{\textcolor[rgb]{0.73,0.13,0.13}{\textit{#1}}}
\newcommand{\ErrorTok}[1]{\textcolor[rgb]{1.00,0.00,0.00}{\textbf{#1}}}
\newcommand{\ExtensionTok}[1]{#1}
\newcommand{\FloatTok}[1]{\textcolor[rgb]{0.25,0.63,0.44}{#1}}
\newcommand{\FunctionTok}[1]{\textcolor[rgb]{0.02,0.16,0.49}{#1}}
\newcommand{\ImportTok}[1]{\textcolor[rgb]{0.00,0.50,0.00}{\textbf{#1}}}
\newcommand{\InformationTok}[1]{\textcolor[rgb]{0.38,0.63,0.69}{\textbf{\textit{#1}}}}
\newcommand{\KeywordTok}[1]{\textcolor[rgb]{0.00,0.44,0.13}{\textbf{#1}}}
\newcommand{\NormalTok}[1]{#1}
\newcommand{\OperatorTok}[1]{\textcolor[rgb]{0.40,0.40,0.40}{#1}}
\newcommand{\OtherTok}[1]{\textcolor[rgb]{0.00,0.44,0.13}{#1}}
\newcommand{\PreprocessorTok}[1]{\textcolor[rgb]{0.74,0.48,0.00}{#1}}
\newcommand{\RegionMarkerTok}[1]{#1}
\newcommand{\SpecialCharTok}[1]{\textcolor[rgb]{0.25,0.44,0.63}{#1}}
\newcommand{\SpecialStringTok}[1]{\textcolor[rgb]{0.73,0.40,0.53}{#1}}
\newcommand{\StringTok}[1]{\textcolor[rgb]{0.25,0.44,0.63}{#1}}
\newcommand{\VariableTok}[1]{\textcolor[rgb]{0.10,0.09,0.49}{#1}}
\newcommand{\VerbatimStringTok}[1]{\textcolor[rgb]{0.25,0.44,0.63}{#1}}
\newcommand{\WarningTok}[1]{\textcolor[rgb]{0.38,0.63,0.69}{\textbf{\textit{#1}}}}
\usepackage{longtable,booktabs,array}
\newcounter{none} % for unnumbered tables
\usepackage{calc} % for calculating minipage widths
% Correct order of tables after \paragraph or \subparagraph
\usepackage{etoolbox}
\makeatletter
\patchcmd\longtable{\par}{\if@noskipsec\mbox{}\fi\par}{}{}
\makeatother
% Allow footnotes in longtable head/foot
\IfFileExists{footnotehyper.sty}{\usepackage{footnotehyper}}{\usepackage{footnote}}
\makesavenoteenv{longtable}
\setlength{\emergencystretch}{3em} % prevent overfull lines
\providecommand{\tightlist}{%
  \setlength{\itemsep}{0pt}\setlength{\parskip}{0pt}}
\usepackage{bookmark}
\IfFileExists{xurl.sty}{\usepackage{xurl}}{} % add URL line breaks if available
\urlstyle{same}
\hypersetup{
  hidelinks,
  pdfcreator={LaTeX via pandoc}}

\author{}
\date{}

\begin{document}

\chapter{VQE Potential Energy Surface \& 3D Molecular Orbital Renderer:
Full Research \& Implementation
Report}\label{vqe-potential-energy-surface-3d-molecular-orbital-renderer-full-research-implementation-report}

\begin{center}\rule{0.5\linewidth}{0.5pt}\end{center}

\section{Part 1 --- Theoretical
Foundations}\label{part-1-theoretical-foundations}

\subsection{1. Quantum Chemistry from First
Principles}\label{quantum-chemistry-from-first-principles}

\subsubsection{1.1 The Molecular Hamiltonian Under the Born-Oppenheimer
Approximation}\label{the-molecular-hamiltonian-under-the-born-oppenheimer-approximation}

The full non-relativistic Hamiltonian for a molecule with \(N_e\)
electrons and \(N_n\) nuclei is (in atomic units,
\(\hbar = m_e = e = 4\pi\epsilon_0 = 1\)):

\[\hat{H}_{\text{mol}} = \hat{T}_n + \hat{T}_e + \hat{V}_{nn} + \hat{V}_{en} + \hat{V}_{ee}\]

where:

\[\hat{T}_n = -\sum_{A=1}^{N_n} \frac{1}{2M_A}\nabla_A^2 \qquad \text{(nuclear kinetic energy)}\]

\[\hat{T}_e = -\sum_{i=1}^{N_e} \frac{1}{2}\nabla_i^2 \qquad \text{(electronic kinetic energy)}\]

\[\hat{V}_{nn} = \sum_{A<B} \frac{Z_A Z_B}{|\mathbf{R}_A - \mathbf{R}_B|} \qquad \text{(nuclear-nuclear repulsion)}\]

\[\hat{V}_{en} = -\sum_{i=1}^{N_e}\sum_{A=1}^{N_n} \frac{Z_A}{|\mathbf{r}_i - \mathbf{R}_A|} \qquad \text{(electron-nuclear attraction)}\]

\[\hat{V}_{ee} = \sum_{i<j} \frac{1}{|\mathbf{r}_i - \mathbf{r}_j|} \qquad \text{(electron-electron repulsion)}\]

Here \(M_A\) and \(Z_A\) are the mass and atomic number of nucleus
\(A\), \(\mathbf{R}_A\) are nuclear coordinates, and \(\mathbf{r}_i\)
are electron coordinates.

The \textbf{Born-Oppenheimer (BO) approximation} exploits the mass ratio
\(m_e / M_A \sim 10^{-3}\)--\(10^{-5}\): nuclei are vastly heavier than
electrons, so electrons respond instantaneously to nuclear motion. Under
this separation, the total wavefunction factorises:

\[\Psi(\mathbf{r}, \mathbf{R}) \approx \psi_{\text{el}}(\mathbf{r}; \mathbf{R}) \cdot \chi_{\text{nuc}}(\mathbf{R})\]

where \(\psi_{\text{el}}\) depends parametrically on nuclear positions
\(\mathbf{R}\). The \textbf{electronic Hamiltonian} at fixed nuclear
geometry is:

\[\hat{H}_{\text{el}}(\mathbf{R}) = \hat{T}_e + \hat{V}_{en} + \hat{V}_{ee} + \hat{V}_{nn}\]

(the nuclear repulsion \(\hat{V}_{nn}\) is a constant at fixed
\(\mathbf{R}\)). Solving the electronic Schrödinger equation:

\[\hat{H}_{\text{el}}|\psi_k(\mathbf{R})\rangle = E_k(\mathbf{R})|\psi_k(\mathbf{R})\rangle\]

yields the electronic energy \(E_k(\mathbf{R})\) as a function of
nuclear coordinates --- this is the \textbf{potential energy surface
(PES)}.

\subsubsection{1.2 The Electronic Structure Problem: Exponential
Complexity}\label{the-electronic-structure-problem-exponential-complexity}

The electronic wavefunction
\(\psi_{\text{el}}(\mathbf{r}_1, \sigma_1, \ldots, \mathbf{r}_{N_e}, \sigma_{N_e})\)
is a function of \(3N_e\) spatial coordinates and \(N_e\) spin
variables. For fermions, it must be antisymmetric under particle
exchange.

Given a basis of \(M\) spin-orbitals \(\{\phi_p\}\), the Hilbert space
of \(N_e\) electrons occupying \(M\) orbitals has dimension
\(\binom{M}{N_e}\). For \(M = 100\) orbitals and \(N_e = 20\) electrons,
\(\binom{100}{20} \approx 5.4 \times 10^{20}\) --- this is the Full
Configuration Interaction (FCI) dimension. Storing a single FCI vector
at double precision would require \(\sim 4 \times 10^9\) TB. This
exponential scaling is why classical exact diagonalisation is
intractable beyond \(\sim 20\) orbitals and is the fundamental
motivation for quantum computing approaches.

\subsubsection{1.3 Second-Quantized
Hamiltonian}\label{second-quantized-hamiltonian}

In second quantisation, we introduce creation \(\hat{a}_p^\dagger\) and
annihilation \(\hat{a}_p\) operators for each spin-orbital \(\phi_p\),
satisfying the fermionic anti-commutation relations:

\[\{\hat{a}_p, \hat{a}_q^\dagger\} = \delta_{pq}, \qquad \{\hat{a}_p, \hat{a}_q\} = 0, \qquad \{\hat{a}_p^\dagger, \hat{a}_q^\dagger\} = 0\]

The electronic Hamiltonian becomes:

\[\hat{H}_{\text{el}} = \sum_{p,q} h_{pq}\, \hat{a}_p^\dagger \hat{a}_q + \frac{1}{2}\sum_{p,q,r,s} h_{pqrs}\, \hat{a}_p^\dagger \hat{a}_q^\dagger \hat{a}_s \hat{a}_r + E_{\text{nuc}}\]

where:

\begin{itemize}
\item
  \textbf{One-electron integrals} (kinetic energy + electron-nuclear
  attraction):
  \[h_{pq} = \int \phi_p^*(\mathbf{x})\left(-\frac{1}{2}\nabla^2 - \sum_A \frac{Z_A}{|\mathbf{r} - \mathbf{R}_A|}\right)\phi_q(\mathbf{x})\,d\mathbf{x}\]
\item
  \textbf{Two-electron integrals} (electron-electron repulsion, in
  chemists' notation):
  \[h_{pqrs} = \iint \frac{\phi_p^*(\mathbf{x}_1)\phi_r(\mathbf{x}_1)\phi_q^*(\mathbf{x}_2)\phi_s(\mathbf{x}_2)}{|\mathbf{r}_1 - \mathbf{r}_2|}\,d\mathbf{x}_1\,d\mathbf{x}_2\]
\end{itemize}

Here \(\mathbf{x} = (\mathbf{r}, \sigma)\) denotes combined spatial and
spin coordinates. The two-electron integrals in \textbf{physicists'
notation} (antisymmetrised) are:

\[\langle pq || rs \rangle = h_{pqrs} - h_{pqsr}\]

\(E_{\text{nuc}} = \hat{V}_{nn}\) is the constant nuclear repulsion
energy at the given geometry.

The number of one-electron integrals scales as \(O(M^2)\) and
two-electron integrals as \(O(M^4)\). These integrals are computed
classically by quantum chemistry packages (PySCF, Psi4, etc.) and serve
as input to the quantum algorithm.

\subsubsection{1.4 Molecular Orbitals, Basis Sets, and the Fock
Matrix}\label{molecular-orbitals-basis-sets-and-the-fock-matrix}

\textbf{Molecular orbitals (MOs)} are expressed as linear combinations
of atomic orbitals (LCAO):

\[\phi_p(\mathbf{r}) = \sum_{\mu=1}^{K} C_{\mu p}\, \chi_\mu(\mathbf{r})\]

where \(\{\chi_\mu\}\) is the \textbf{atomic orbital (AO) basis set} and
\(C_{\mu p}\) are the MO coefficients obtained by solving the
Hartree-Fock equations.

\textbf{Basis sets} approximate the true atomic orbitals with
Gaussian-type orbitals (GTOs):

\begin{itemize}
\item
  \textbf{STO-3G (Minimal basis):} Each Slater-type orbital (STO) is
  approximated by 3 Gaussian primitives. For H₂: 2 basis functions (one
  per atom). For H₂O: 7 basis functions (1s, 2s, 2p×3 on O; 1s on each
  H). This is the smallest possible basis --- qualitatively correct but
  quantitatively poor.
\item
  \textbf{6-31G* (Split-valence + polarisation):} Valence orbitals are
  split into two functions (3+1 Gaussians), core orbitals use 6
  Gaussians, and polarisation d-functions are added on heavy atoms. For
  H₂O: 19 basis functions. Significantly more accurate, especially for
  geometries and relative energies.
\end{itemize}

The \textbf{Fock matrix} \(\mathbf{F}\) in the AO basis is:

\[F_{\mu\nu} = h_{\mu\nu}^{\text{core}} + \sum_{\lambda\sigma} P_{\lambda\sigma}\left[(\mu\nu|\lambda\sigma) - \frac{1}{2}(\mu\lambda|\nu\sigma)\right]\]

where \(h^{\text{core}}\) is the one-electron (core) Hamiltonian matrix,
\(P_{\lambda\sigma}\) is the density matrix, and
\((\mu\nu|\lambda\sigma)\) are two-electron integrals in the AO basis.
The Hartree-Fock equations (Roothaan-Hall equations) are the generalised
eigenvalue problem:

\[\mathbf{F}\mathbf{C} = \mathbf{S}\mathbf{C}\boldsymbol{\epsilon}\]

where \(\mathbf{S}\) is the overlap matrix
\(S_{\mu\nu} = \langle\chi_\mu|\chi_\nu\rangle\) and
\(\boldsymbol{\epsilon}\) contains the orbital energies. This is solved
self-consistently (SCF): the Fock matrix depends on the density matrix
\(\mathbf{P} = \mathbf{C}_{\text{occ}}\mathbf{C}_{\text{occ}}^T\), which
depends on \(\mathbf{C}\), which comes from diagonalising
\(\mathbf{F}\).

\subsubsection{1.5 Hartree-Fock Theory and Its Failure at Bond
Dissociation}\label{hartree-fock-theory-and-its-failure-at-bond-dissociation}

\textbf{Restricted Hartree-Fock (RHF)} constrains the wavefunction to be
a single Slater determinant:

\[|\Phi_{\text{HF}}\rangle = \hat{a}_{1}^\dagger \hat{a}_{2}^\dagger \cdots \hat{a}_{N_e}^\dagger |0\rangle\]

This captures the mean-field electron-electron interaction (each
electron moves in the average potential of all others) but misses
\textbf{correlation energy} --- the energy difference between the exact
solution and the HF solution:

\[E_{\text{corr}} = E_{\text{exact}} - E_{\text{HF}} < 0\]

HF works well near equilibrium geometries where the single-determinant
picture is qualitatively correct. However, it fails catastrophically at
bond dissociation because:

\begin{enumerate}
\def\labelenumi{\arabic{enumi}.}
\tightlist
\item
  \textbf{Static (strong) correlation:} At large internuclear
  separation, the correct wavefunction becomes a superposition of
  multiple determinants. For H₂ dissociation, the exact wavefunction at
  \(R \to \infty\) is:
\end{enumerate}

\[|\Psi\rangle \to \frac{1}{\sqrt{2}}(|\sigma_g\bar{\sigma}_g\rangle - |\sigma_u\bar{\sigma}_u\rangle)\]

(a superposition of doubly-occupied bonding and anti-bonding
configurations). RHF forces a single determinant
\(|\sigma_g\bar{\sigma}_g\rangle\), which incorrectly mixes in ionic
terms (\(\text{H}^+\text{H}^-\)) at dissociation, predicting an energy
\(\sim 0.25\) Hartree too high.

\begin{enumerate}
\def\labelenumi{\arabic{enumi}.}
\setcounter{enumi}{1}
\tightlist
\item
  \textbf{Dynamic correlation:} Short-range electron-electron cusp
  conditions are never captured by a single determinant. This
  contributes a correlation energy typically \(\sim 1\%\) of total
  energy but \(\sim 100\%\) of chemical bond energies.
\end{enumerate}

Post-HF methods (MP2, CISD, CCSD(T), FCI) recover correlation energy at
increasing computational cost. VQE on a quantum computer can in
principle reach FCI accuracy with polynomial quantum resources.

\begin{center}\rule{0.5\linewidth}{0.5pt}\end{center}

\subsection{2. Fermion-to-Qubit
Mappings}\label{fermion-to-qubit-mappings}

\subsubsection{2.1 Jordan-Wigner
Transformation}\label{jordan-wigner-transformation}

The Jordan-Wigner (JW) transformation maps fermionic operators to qubit
(Pauli) operators by encoding the occupation of each spin-orbital in a
qubit and using \(Z\)-strings to enforce anti-commutation:

\[\hat{a}_p^\dagger \mapsto \frac{1}{2}(\hat{X}_p - i\hat{Y}_p) \otimes \prod_{q=0}^{p-1} \hat{Z}_q = \frac{1}{2}\hat{\sigma}_p^- \bigotimes_{q<p} \hat{Z}_q\]

\[\hat{a}_p \mapsto \frac{1}{2}(\hat{X}_p + i\hat{Y}_p) \otimes \prod_{q=0}^{p-1} \hat{Z}_q = \frac{1}{2}\hat{\sigma}_p^+ \bigotimes_{q<p} \hat{Z}_q\]

where \(\hat{\sigma}_p^{\pm} = (\hat{X}_p \mp i\hat{Y}_p)/2\) are the
qubit raising/lowering operators and the \(Z\)-string
\(\prod_{q<p}\hat{Z}_q\) provides the fermionic sign (parity of all
lower-indexed orbitals).

\textbf{Verification of anti-commutation:} Consider
\(\hat{a}_p^\dagger \hat{a}_q + \hat{a}_q \hat{a}_p^\dagger\) for
\(p \neq q\). The \(Z\)-strings on the overlapping qubits between \(p\)
and \(q\) introduce a sign flip, yielding:

\[\{\hat{a}_p, \hat{a}_q^\dagger\}_{\text{JW}} = \delta_{pq}\hat{I}\]

as required.

\textbf{Number operator:}
\(\hat{n}_p = \hat{a}_p^\dagger \hat{a}_p \mapsto \frac{1}{2}(\hat{I} - \hat{Z}_p)\),
which has eigenvalue 0 on \(|0\rangle\) (unoccupied) and 1 on
\(|1\rangle\) (occupied).

\textbf{One-body terms under JW:}

\[\hat{a}_p^\dagger \hat{a}_q \mapsto \frac{1}{4}\bigl[(\hat{X}_p\hat{X}_q + \hat{Y}_p\hat{Y}_q) + i(\hat{X}_p\hat{Y}_q - \hat{Y}_p\hat{X}_q)\bigr]\prod_{k=\min(p,q)+1}^{\max(p,q)-1}\hat{Z}_k\]

For \(p = q\):
\(\hat{a}_p^\dagger\hat{a}_p = \frac{1}{2}(\hat{I} - \hat{Z}_p)\).

The \(Z\)-string between indices \(p\) and \(q\) makes the operator
weight (number of Pauli terms) scale as \(O(N)\) for distant orbitals.
This is the principal cost of JW: non-local operators.

\textbf{Two-body terms under JW:} Each
\(\hat{a}_p^\dagger\hat{a}_q^\dagger\hat{a}_s\hat{a}_r\) term generates
a sum of Pauli strings of weight up to \(O(N)\). The total number of
Pauli terms in the qubit Hamiltonian scales as \(O(M^4)\).

\subsubsection{2.2 Parity Mapping}\label{parity-mapping}

The \textbf{parity mapping} (Bravyi-Kitaev parity basis) encodes the
parity of the occupation number rather than the occupation itself. Qubit
\(p\) stores the parity \(\bigoplus_{q \leq p} n_q\) (sum modulo 2 of
occupations up to orbital \(p\)).

The creation operator becomes:

\[\hat{a}_p^\dagger \mapsto \frac{1}{2}\hat{X}_p \bigotimes_{q \in U(p)} \hat{X}_q \cdot \bigotimes_{k \in P(p)} \hat{Z}_k - \frac{i}{2}\hat{Y}_p \bigotimes_{q \in U(p)} \hat{X}_q\]

where \(U(p)\) and \(P(p)\) are the ``update'' and ``parity'' sets
determined by the binary tree structure.

\textbf{Qubit reduction advantage:} In the parity basis, the last qubit
stores the total parity of the electron number, and the second-to-last
stores the total spin parity. For systems with fixed electron number and
fixed spin (which is always the case in chemistry), these two qubits are
redundant and can be removed, reducing the qubit count by 2:

\[M \text{ spin-orbitals} \to M - 2 \text{ qubits (parity mapping with tapering)}\]

This is significant for small molecules: H₂ (STO-3G) goes from 4 qubits
(JW) to 2 qubits (parity + tapering).

\subsubsection{2.3 Worked Example: H₂ in
STO-3G}\label{worked-example-hux2082-in-sto-3g}

\textbf{Setup:} H₂ has 2 electrons in a minimal basis of 2 spatial
orbitals (\(\sigma_g\), \(\sigma_u\)) giving 4 spin-orbitals:
\(\sigma_{g\alpha}\), \(\sigma_{g\beta}\), \(\sigma_{u\alpha}\),
\(\sigma_{u\beta}\). This maps to 4 qubits under JW.

\textbf{Electronic integrals (at equilibrium \(R = 0.735\) Å):}

Using PySCF to compute the integrals and PennyLane's \texttt{qml.qchem}
to perform the JW mapping, the qubit Hamiltonian for H₂ (STO-3G) is:

\[\hat{H} = g_0 \hat{I} + g_1 \hat{Z}_0 + g_2 \hat{Z}_1 + g_3 \hat{Z}_2 + g_4 \hat{Z}_3 + g_5 \hat{Z}_0\hat{Z}_1 + g_6 \hat{Z}_0\hat{Z}_2 + g_7 \hat{Z}_0\hat{Z}_3 + g_8 \hat{Z}_1\hat{Z}_2 + g_9 \hat{Z}_1\hat{Z}_3 + g_{10}\hat{Z}_2\hat{Z}_3 + g_{11}\hat{X}_0\hat{X}_1\hat{Y}_2\hat{Y}_3 + g_{12}\hat{X}_0\hat{Y}_1\hat{Y}_2\hat{X}_3 + g_{13}\hat{Y}_0\hat{X}_1\hat{X}_2\hat{Y}_3 + g_{14}\hat{Y}_0\hat{Y}_1\hat{X}_2\hat{X}_3\]

with numerical coefficients (at \(R = 0.735\) Å):

{\def\LTcaptype{none} % do not increment counter
\begin{longtable}[]{@{}lc@{}}
\toprule\noalign{}
Term & Coefficient (Hartree) \\
\midrule\noalign{}
\endhead
\bottomrule\noalign{}
\endlastfoot
\(\hat{I}\) & \(-0.0988\) \\
\(\hat{Z}_0\) & \(+0.1711\) \\
\(\hat{Z}_1\) & \(+0.1711\) \\
\(\hat{Z}_2\) & \(-0.2232\) \\
\(\hat{Z}_3\) & \(-0.2232\) \\
\(\hat{Z}_0\hat{Z}_1\) & \(+0.1686\) \\
\(\hat{Z}_0\hat{Z}_2\) & \(+0.1205\) \\
\(\hat{Z}_0\hat{Z}_3\) & \(+0.1659\) \\
\(\hat{Z}_1\hat{Z}_2\) & \(+0.1659\) \\
\(\hat{Z}_1\hat{Z}_3\) & \(+0.1205\) \\
\(\hat{Z}_2\hat{Z}_3\) & \(+0.1743\) \\
\(\hat{X}_0\hat{X}_1\hat{Y}_2\hat{Y}_3\) & \(-0.0454\) \\
\(\hat{X}_0\hat{Y}_1\hat{Y}_2\hat{X}_3\) & \(+0.0454\) \\
\(\hat{Y}_0\hat{X}_1\hat{X}_2\hat{Y}_3\) & \(+0.0454\) \\
\(\hat{Y}_0\hat{Y}_1\hat{X}_2\hat{X}_3\) & \(-0.0454\) \\
\end{longtable}
}

The 15-term qubit Hamiltonian acts on \(2^4 = 16\)-dimensional Hilbert
space. The HF reference state is \(|1100\rangle\) (first two
spin-orbitals occupied). The FCI ground-state energy at equilibrium is
\(E_{\text{FCI}} = -1.1373\) Hartree.

With \textbf{parity mapping + 2-qubit tapering}, this reduces to a
2-qubit Hamiltonian:

\[\hat{H}_{\text{tapered}} = a_0\hat{I} + a_1\hat{Z}_0 + a_2\hat{Z}_1 + a_3\hat{Z}_0\hat{Z}_1 + a_4\hat{X}_0\hat{X}_1\]

with only 5 terms on 2 qubits --- dramatically simpler to simulate.

\subsubsection{2.4 Qubit Count and Circuit Depth
Comparison}\label{qubit-count-and-circuit-depth-comparison}

{\def\LTcaptype{none} % do not increment counter
\begin{longtable}[]{@{}
  >{\raggedright\arraybackslash}p{(\linewidth - 14\tabcolsep) * \real{0.1266}}
  >{\raggedright\arraybackslash}p{(\linewidth - 14\tabcolsep) * \real{0.0886}}
  >{\centering\arraybackslash}p{(\linewidth - 14\tabcolsep) * \real{0.1519}}
  >{\centering\arraybackslash}p{(\linewidth - 14\tabcolsep) * \real{0.1139}}
  >{\centering\arraybackslash}p{(\linewidth - 14\tabcolsep) * \real{0.1139}}
  >{\centering\arraybackslash}p{(\linewidth - 14\tabcolsep) * \real{0.1392}}
  >{\centering\arraybackslash}p{(\linewidth - 14\tabcolsep) * \real{0.1266}}
  >{\centering\arraybackslash}p{(\linewidth - 14\tabcolsep) * \real{0.1392}}@{}}
\toprule\noalign{}
\begin{minipage}[b]{\linewidth}\raggedright
Molecule
\end{minipage} & \begin{minipage}[b]{\linewidth}\raggedright
Basis
\end{minipage} & \begin{minipage}[b]{\linewidth}\centering
Spatial Orbs
\end{minipage} & \begin{minipage}[b]{\linewidth}\centering
Spin-Orbs (M)
\end{minipage} & \begin{minipage}[b]{\linewidth}\centering
JW Qubits
\end{minipage} & \begin{minipage}[b]{\linewidth}\centering
Parity (tapered)
\end{minipage} & \begin{minipage}[b]{\linewidth}\centering
UCCSD Params
\end{minipage} & \begin{minipage}[b]{\linewidth}\centering
CNOT Depth (UCCSD)
\end{minipage} \\
\midrule\noalign{}
\endhead
\bottomrule\noalign{}
\endlastfoot
H₂ & STO-3G & 2 & 4 & 4 & 2 & 3 & \textasciitilde16 \\
LiH & STO-3G & 6 & 12 & 12 & 10 & \textasciitilde52 &
\textasciitilde1200 \\
H₂O & STO-3G & 7 & 14 & 14 & 12 & \textasciitilde80 &
\textasciitilde3000 \\
H₂O & 6-31G* & 19 & 38 & 38 & 36 & \textasciitilde2000 &
\textasciitilde80000 \\
\end{longtable}
}

The CNOT depth scales roughly as \(O(M^3 N_e)\) for UCCSD, making large
molecules intractable on NISQ hardware without active space reduction.

\begin{center}\rule{0.5\linewidth}{0.5pt}\end{center}

\subsection{3. Variational Quantum Eigensolver
(VQE)}\label{variational-quantum-eigensolver-vqe}

\subsubsection{3.1 The Variational
Principle}\label{the-variational-principle}

The variational principle states that for any normalised trial state
\(|\psi(\boldsymbol{\theta})\rangle\):

\[E(\boldsymbol{\theta}) = \langle\psi(\boldsymbol{\theta})|\hat{H}|\psi(\boldsymbol{\theta})\rangle \geq E_0\]

where \(E_0\) is the exact ground-state energy. This follows directly
from expanding \(|\psi(\boldsymbol{\theta})\rangle\) in the eigenbasis
\(\{|E_k\rangle\}\) of \(\hat{H}\):

\[E(\boldsymbol{\theta}) = \sum_k |c_k|^2 E_k \geq E_0 \sum_k |c_k|^2 = E_0\]

since \(E_k \geq E_0\) for all \(k\) and \(\sum_k|c_k|^2 = 1\).

This means any parameterised quantum circuit that prepares
\(|\psi(\boldsymbol{\theta})\rangle\) gives an \textbf{upper bound} on
the ground-state energy. By minimising \(E(\boldsymbol{\theta})\) over
\(\boldsymbol{\theta}\), we approach \(E_0\) from above. The tightness
of the bound depends on the expressibility of the ansatz --- whether the
parameterised family can represent (or closely approximate) the true
ground state.

\subsubsection{3.2 The UCCSD Ansatz}\label{the-uccsd-ansatz}

\textbf{Classical Coupled Cluster (CC):} The coupled-cluster
wavefunction is:

\[|\Psi_{\text{CC}}\rangle = e^{\hat{T}}|\Phi_{\text{HF}}\rangle\]

where \(\hat{T} = \hat{T}_1 + \hat{T}_2 + \cdots\) is the cluster
operator. For CCSD (Singles and Doubles):

\[\hat{T}_1 = \sum_{i \in \text{occ}} \sum_{a \in \text{virt}} t_i^a\, \hat{a}_a^\dagger \hat{a}_i\]

\[\hat{T}_2 = \sum_{i<j \in \text{occ}} \sum_{a<b \in \text{virt}} t_{ij}^{ab}\, \hat{a}_a^\dagger \hat{a}_b^\dagger \hat{a}_j \hat{a}_i\]

where \(i, j\) label occupied orbitals, \(a, b\) label virtual
(unoccupied) orbitals, and \(t_i^a\), \(t_{ij}^{ab}\) are the cluster
amplitudes (variational parameters).

\textbf{Problem for quantum computing:} The classical CC operator
\(e^{\hat{T}}\) is not unitary (\(\hat{T}\) is not anti-Hermitian), so
it cannot be directly implemented as a quantum circuit.

\textbf{Unitary Coupled Cluster (UCC):} Replace \(\hat{T}\) with the
anti-Hermitian operator \(\hat{T} - \hat{T}^\dagger\):

\[|\Psi_{\text{UCC}}\rangle = e^{\hat{T} - \hat{T}^\dagger}|\Phi_{\text{HF}}\rangle\]

This is manifestly unitary: \(\hat{U} = e^{\hat{T} - \hat{T}^\dagger}\)
satisfies \(\hat{U}^\dagger\hat{U} = \hat{I}\) since
\((\hat{T} - \hat{T}^\dagger)^\dagger = -(\hat{T} - \hat{T}^\dagger)\).

For UCCSD:

\[\hat{T} - \hat{T}^\dagger = \sum_{i,a} t_i^a(\hat{a}_a^\dagger\hat{a}_i - \hat{a}_i^\dagger\hat{a}_a) + \sum_{i<j,a<b} t_{ij}^{ab}(\hat{a}_a^\dagger\hat{a}_b^\dagger\hat{a}_j\hat{a}_i - \hat{a}_i^\dagger\hat{a}_j^\dagger\hat{a}_b\hat{a}_a)\]

The number of parameters is
\(n_{\text{singles}} = N_{\text{occ}} \times N_{\text{virt}}\) plus
\(n_{\text{doubles}} = \binom{N_{\text{occ}}}{2}\binom{N_{\text{virt}}}{2}\).

\subsubsection{3.3 Trotterisation of UCCSD into
Gates}\label{trotterisation-of-uccsd-into-gates}

The UCC unitary \(e^{\hat{T}-\hat{T}^\dagger}\) cannot be decomposed
exactly into single/two-qubit gates because the individual excitation
operators do not commute. We use the first-order Suzuki-Trotter
decomposition:

\[e^{\hat{T}-\hat{T}^\dagger} \approx \prod_{\mu} e^{\theta_\mu \hat{\tau}_\mu}\]

where \(\hat{\tau}_\mu\) are individual excitation generators (single or
double), \(\theta_\mu\) are the corresponding amplitudes, and the
product is over all excitations.

\textbf{Single excitation gate:} For the excitation \(i \to a\):

\[e^{\theta(a_a^\dagger a_i - a_i^\dagger a_a)}\]

Under JW, this maps to a rotation involving \(X\), \(Y\) Pauli strings
with \(Z\)-strings between orbital indices \(i\) and \(a\). PennyLane
provides \texttt{qml.SingleExcitation(theta,\ wires={[}i,\ a{]})} which
implements this as a Givens rotation:

\[\text{SingleExcitation}(\theta) = \begin{pmatrix} 1 & 0 & 0 & 0 \\ 0 & \cos(\theta/2) & -\sin(\theta/2) & 0 \\ 0 & \sin(\theta/2) & \cos(\theta/2) & 0 \\ 0 & 0 & 0 & 1 \end{pmatrix}\]

acting on the \(\{|01\rangle, |10\rangle\}\) subspace of wires
\([i, a]\).

\textbf{Double excitation gate:} For \(i, j \to a, b\):

\[e^{\theta(a_a^\dagger a_b^\dagger a_j a_i - a_i^\dagger a_j^\dagger a_b a_a)}\]

PennyLane provides
\texttt{qml.DoubleExcitation(theta,\ wires={[}i,\ j,\ a,\ b{]})} which
implements a 4-qubit rotation in the \(\{|1100\rangle, |0011\rangle\}\)
subspace. This gate decomposes into \(\sim 16\) CNOT gates.

\textbf{Trotterisation error:} The first-order Trotter formula has error
\(O(\sum_{\mu < \nu} \|[\hat{\tau}_\mu, \hat{\tau}_\nu]\| \cdot |\theta_\mu||\theta_\nu|)\).
In practice, for VQE, the Trotter error is absorbed into the variational
parameters --- the optimizer compensates for the discretisation, so the
ground-state energy can still be reached if the Trotterised ansatz is
expressive enough.

\subsubsection{3.4 Hardware-Efficient
Ansatz}\label{hardware-efficient-ansatz}

An alternative to UCCSD is the \textbf{hardware-efficient ansatz (HEA)},
which consists of layers of native single-qubit rotations and entangling
gates that are directly executable on a given hardware topology:

\[|\psi(\boldsymbol{\theta})\rangle = \prod_{l=1}^{L}\left[\hat{U}_{\text{ent}} \cdot \prod_{i=1}^{N_q} R_Y(\theta_i^{(l)}) R_Z(\phi_i^{(l)})\right] |\Phi_{\text{HF}}\rangle\]

where \(\hat{U}_{\text{ent}}\) is a layer of CNOT (or CZ) gates
following the hardware connectivity.

\textbf{Advantages over UCCSD:} - Much shallower circuits (fewer CNOTs
per layer). - No chemistry-specific structure needed --- can be used for
any Hamiltonian. - Circuit depth is a tunable hyperparameter (number of
layers \(L\)).

\textbf{Disadvantages:} - No physical motivation: the circuit does not
encode any knowledge of the electronic structure problem. - Suffers from
barren plateaus more readily as \(N_q\) and \(L\) increase. - Does not
preserve electron number or spin symmetry unless explicitly constrained.
- Convergence to the ground state is not guaranteed (the landscape may
have many local minima).

For small molecules (H₂, LiH), HEA with 1--3 layers can reach chemical
accuracy. For larger molecules, UCCSD (or adaptive variants like
ADAPT-VQE) is generally superior.

\subsubsection{3.5 The Full VQE Loop}\label{the-full-vqe-loop}

\begin{enumerate}
\def\labelenumi{\arabic{enumi}.}
\item
  \textbf{State preparation:} Prepare
  \(|\psi(\boldsymbol{\theta})\rangle = \hat{U}(\boldsymbol{\theta})|\Phi_{\text{HF}}\rangle\)
  on the quantum computer.
\item
  \textbf{Measurement:} Measure
  \(\langle\psi(\boldsymbol{\theta})|\hat{H}|\psi(\boldsymbol{\theta})\rangle\)
  by decomposing \(\hat{H}\) into Pauli terms:
  \[\hat{H} = \sum_k c_k \hat{P}_k\] Measure each
  \(\langle\hat{P}_k\rangle\) from repeated preparations + measurements.
  Group commuting Pauli terms to reduce the number of distinct
  measurement circuits.
\item
  \textbf{Classical optimization:} Update \(\boldsymbol{\theta}\) using
  a classical optimizer (gradient-based or gradient-free) to minimise
  \(E(\boldsymbol{\theta})\).
\item
  \textbf{Convergence:} Repeat until
  \(|E(\boldsymbol{\theta}^{(k)}) - E(\boldsymbol{\theta}^{(k-1)})| < \varepsilon\)
  or maximum iterations reached.
\end{enumerate}

The hybrid quantum-classical nature of VQE keeps quantum circuit depth
shallow (only state preparation + measurement) while offloading the
optimisation to classical hardware. The bottleneck is the number of
circuit evaluations needed for sufficient shot statistics and optimiser
convergence.

\begin{center}\rule{0.5\linewidth}{0.5pt}\end{center}

\subsection{4. Potential Energy Surface
(PES)}\label{potential-energy-surface-pes}

\subsubsection{4.1 Definition}\label{definition}

The \textbf{potential energy surface} (PES) is the electronic energy
\(E_0(\mathbf{R})\) as a function of nuclear coordinates \(\mathbf{R}\):

\[\text{PES}: \mathbb{R}^{3N_n - 6} \to \mathbb{R}\]

(subtracting 3 translational and 3 rotational degrees of freedom, or
\(3N_n - 5\) for linear molecules).

For a diatomic molecule like H₂, the PES is a one-dimensional curve
\(E(R)\) where \(R\) is the internuclear distance. For polyatomics, the
PES is a high-dimensional hypersurface.

\subsubsection{4.2 What the PES Reveals}\label{what-the-pes-reveals}

\begin{itemize}
\item
  \textbf{Equilibrium geometry \(R_e\):} The minimum of \(E(R)\), where
  \(dE/dR = 0\) and \(d^2E/dR^2 > 0\). For H₂, \(R_e \approx 0.74\) Å.
\item
  \textbf{Dissociation energy \(D_e\):} The energy difference between
  the minimum and the asymptotic separated-atom limit:
  \(D_e = E(R \to \infty) - E(R_e)\). For H₂, \(D_e \approx 4.75\) eV
  (exact: \(4.7467\) eV).
\item
  \textbf{Bond stretching frequency \(\omega_e\):} Related to the
  curvature at the minimum: \(\omega_e = \sqrt{k/\mu}\) where
  \(k = d^2E/dR^2|_{R_e}\) and \(\mu\) is the reduced mass.
\item
  \textbf{Transition states:} Saddle points on the PES (for polyatomics)
  correspond to transition states connecting reactants and products.
\end{itemize}

\subsubsection{4.3 Accuracy Hierarchy}\label{accuracy-hierarchy}

The hierarchy of computational chemistry methods, in increasing accuracy
and cost:

{\def\LTcaptype{none} % do not increment counter
\begin{longtable}[]{@{}
  >{\raggedright\arraybackslash}p{(\linewidth - 6\tabcolsep) * \real{0.1951}}
  >{\centering\arraybackslash}p{(\linewidth - 6\tabcolsep) * \real{0.2195}}
  >{\centering\arraybackslash}p{(\linewidth - 6\tabcolsep) * \real{0.2927}}
  >{\centering\arraybackslash}p{(\linewidth - 6\tabcolsep) * \real{0.2927}}@{}}
\toprule\noalign{}
\begin{minipage}[b]{\linewidth}\raggedright
Method
\end{minipage} & \begin{minipage}[b]{\linewidth}\centering
Scaling
\end{minipage} & \begin{minipage}[b]{\linewidth}\centering
Correlation
\end{minipage} & \begin{minipage}[b]{\linewidth}\centering
PES Quality
\end{minipage} \\
\midrule\noalign{}
\endhead
\bottomrule\noalign{}
\endlastfoot
HF & \(O(M^4)\) & None & Fails at dissociation \\
MP2 & \(O(M^5)\) & Perturbative & Fair near equilibrium, fails at
dissociation \\
CISD & \(O(M^6)\) & Variational, not size-consistent & Good near
equilibrium \\
CCSD(T) & \(O(M^7)\) & ``Gold standard'' of single-ref methods &
Excellent near equilibrium \\
FCI & Exponential & Exact (within basis) & Exact reference \\
VQE (UCCSD) & Quantum poly & Variational, size-consistent & Approaches
FCI \\
\end{longtable}
}

\textbf{Chemical accuracy} is defined as \(\Delta E \leq 1.6\) mHartree
\(\approx 1\) kcal/mol \(\approx 43\) meV. Achieving chemical accuracy
across the entire PES (including dissociation) requires a method that
correctly handles strong correlation --- this is where VQE with a
multi-reference-capable ansatz can outperform CCSD(T).

\subsubsection{4.4 H₂ Dissociation Curve: The
Benchmark}\label{hux2082-dissociation-curve-the-benchmark}

The H₂ molecule is the standard test case for VQE. The exact
(FCI/STO-3G) dissociation curve has:

\begin{itemize}
\tightlist
\item
  \(R_e = 0.735\) Å, \(E(R_e) = -1.1373\) Hartree
\item
  \(E(R \to \infty) = -0.9997\) Hartree (two separated hydrogen atoms)
\item
  \(D_e = 0.1376\) Hartree \(= 3.745\) eV (in STO-3G; the basis set
  limit value is higher)
\end{itemize}

\textbf{RHF failure:} At \(R = 3.0\) Å, the RHF energy is \(\sim -0.90\)
Hartree while the exact energy is \(\sim -1.00\) Hartree --- an error of
\(\sim 0.1\) Hartree (\(\sim 63\) kcal/mol), completely unacceptable.

\textbf{VQE target:} Reproduce the FCI/STO-3G curve to within
\textbf{chemical accuracy} (1.6 mHartree) at all bond lengths
\(R \in [0.3, 3.0]\) Å. With the UCCSD ansatz on 4 qubits (or 2 qubits
with parity + tapering), this is achievable in noiseless simulation and
is the primary validation target.

\begin{center}\rule{0.5\linewidth}{0.5pt}\end{center}

\subsection{5. Molecular Orbitals and Wavefunction
Visualization}\label{molecular-orbitals-and-wavefunction-visualization}

\subsubsection{5.1 LCAO Expansion}\label{lcao-expansion}

Each molecular orbital is expanded in the AO basis:

\[\phi_p(\mathbf{r}) = \sum_{\mu=1}^{K} C_{\mu p}\, \chi_\mu(\mathbf{r})\]

where \(\chi_\mu\) are Gaussian-type basis functions centred on nuclei:

\[\chi_\mu(\mathbf{r}) = N_\mu\, (x - A_x)^{l_x}(y - A_y)^{l_y}(z - A_z)^{l_z}\, e^{-\alpha_\mu |\mathbf{r} - \mathbf{A}|^2}\]

with angular momentum quantum numbers \((l_x, l_y, l_z)\), exponent
\(\alpha_\mu\), centre \(\mathbf{A}\), and normalisation \(N_\mu\). The
MO coefficients \(C_{\mu p}\) come from the converged HF calculation.

\subsubsection{\texorpdfstring{5.2 Evaluating \(\psi(\mathbf{r})\) on a
3D
Grid}{5.2 Evaluating \textbackslash psi(\textbackslash mathbf\{r\}) on a 3D Grid}}\label{evaluating-psimathbfr-on-a-3d-grid}

To visualise an MO, evaluate \(\phi_p(\mathbf{r})\) on a uniform 3D
Cartesian grid:

\begin{enumerate}
\def\labelenumi{\arabic{enumi}.}
\item
  Define the grid:
  \(\mathbf{r}_{ijk} = (x_{\min} + i\Delta x, \, y_{\min} + j\Delta y, \, z_{\min} + k\Delta z)\)
  for \(i = 0, \ldots, N_x-1\), etc. Typical grid:
  \(\Delta x = \Delta y = \Delta z = 0.1\) Bohr, extent \(\pm 5\) Bohr
  from molecular centre, giving \(\sim 100^3 = 10^6\) grid points.
\item
  For each grid point, compute each AO value:
  \[\chi_\mu(\mathbf{r}_{ijk}) = N_\mu \prod_d (r_d - A_d)^{l_d} \exp(-\alpha_\mu|\mathbf{r}_{ijk} - \mathbf{A}|^2)\]
\item
  Contract with MO coefficients:
  \[\phi_p(\mathbf{r}_{ijk}) = \sum_\mu C_{\mu p}\, \chi_\mu(\mathbf{r}_{ijk})\]
\end{enumerate}

This produces a \textbf{volumetric scalar field}
\(V_p[i,j,k] = \phi_p(\mathbf{r}_{ijk})\) of shape \((N_x, N_y, N_z)\).

\textbf{Optimisation:} For GTOs, the exponential decay makes most basis
functions negligible beyond a few Bohr from their centre. Use a cutoff
radius per AO (e.g., where \(\chi_\mu < 10^{-8}\)) and only evaluate
within that radius. The evaluation is vectorised by forming a
displacement tensor \(\Delta_{g\mu d} = r_{gd} - A_{\mu d}\) of shape
\((N_{\text{grid}}, K, 3)\), computing squared distances
\(|\Delta_{g\mu}|^2 = \sum_d \Delta_{g\mu d}^2\), and evaluating all
Gaussians in a single broadcasted exponential
\(G_{g\mu} = \exp(-\alpha_\mu |\Delta_{g\mu}|^2)\). The angular
polynomial prefactors are applied element-wise, and the final MO values
are obtained by contracting with MO coefficients:
\(\phi_p(\mathbf{r}_g) = \sum_\mu C_{\mu p}\, [\text{angular}]_{g\mu}\, G_{g\mu}\).

\emph{(Full implementation: \texttt{report2\_scripts/orbital\_grid.py})}

\subsubsection{5.3 Electron Density and the One-Particle Reduced Density
Matrix}\label{electron-density-and-the-one-particle-reduced-density-matrix}

The \textbf{electron density} is the sum over occupied MOs:

\[\rho(\mathbf{r}) = \sum_{i \in \text{occ}} n_i\, |\phi_i(\mathbf{r})|^2\]

where \(n_i\) is the occupation number (2 for RHF doubly-occupied
orbitals). More generally, for a correlated wavefunction with
one-particle reduced density matrix (1-RDM) \(\gamma_{pq}\):

\[\rho(\mathbf{r}) = \sum_{p,q} \gamma_{pq}\, \phi_p^*(\mathbf{r})\phi_q(\mathbf{r})\]

where
\(\gamma_{pq} = \langle\Psi|\hat{a}_p^\dagger\hat{a}_q|\Psi\rangle\).
For VQE, the 1-RDM can be measured on the quantum computer by measuring
\(\langle\hat{a}_p^\dagger\hat{a}_q\rangle\) in the Pauli basis (each
\(\hat{a}_p^\dagger\hat{a}_q\) maps to a sum of Pauli strings under JW).

\subsubsection{5.4 Isosurface Extraction}\label{isosurface-extraction}

An \textbf{isosurface} is the set of points where a scalar field equals
a constant value:

\[S_c = \{\mathbf{r} : \phi_p(\mathbf{r}) = c\}\]

For MO visualisation, the conventional isovalues are \(c = \pm 0.05\)
a.u. (atomic units). The positive lobe (\(c = +0.05\)) shows where the
wavefunction is positive; the negative lobe (\(c = -0.05\)) shows where
it's negative. These lobes are typically coloured blue/red respectively.

\textbf{Why \(\pm 0.05\) a.u.?} This isovalue encloses approximately
90\% of the electron density associated with the orbital, providing a
visually useful representation of the orbital shape without being too
diffuse (smaller \(|c|\)) or too compact (larger \(|c|\)).

The \textbf{marching cubes algorithm} (Lorensen \& Cline, 1987) extracts
a triangle mesh approximating the isosurface from the volumetric data.
For each voxel (cube formed by 8 neighbouring grid points):

\begin{enumerate}
\def\labelenumi{\arabic{enumi}.}
\tightlist
\item
  Classify each corner as inside (\(\phi > c\)) or outside
  (\(\phi < c\)) the isosurface.
\item
  The 8 binary corner classifications give \(2^8 = 256\) possible
  configurations, reduced to 15 unique topological cases by symmetry.
\item
  Look up the corresponding triangle configuration in a precomputed
  table (the \textbf{marching cubes lookup table}).
\item
  Interpolate vertex positions along cube edges where the isosurface
  crosses (linear interpolation between corner values).
\item
  Compute face normals from the gradient of the scalar field at each
  vertex (or from the triangle vertices directly).
\end{enumerate}

\begin{center}\rule{0.5\linewidth}{0.5pt}\end{center}

\section{Part 2 --- Software
Architecture}\label{part-2-software-architecture}

\subsection{Project Structure}\label{project-structure}

\begin{verbatim}
vqe_explorer/
├── main.py                  # Orchestration, CLI, entry point
├── molecule.py              # Molecule definition, geometry, basis set
├── classical_hf.py          # PySCF interface: HF, integrals, MO coefficients
├── hamiltonian_builder.py   # Second-quantized → qubit Hamiltonian
├── ansatz.py                # UCCSD and hardware-efficient ansätze
├── vqe_solver.py            # VQE optimization loop
├── pes_scanner.py           # PES sweep over bond lengths
├── orbital_grid.py          # MO evaluation on 3D grid
├── isosurface.py            # Marching cubes isosurface extraction
├── renderer.py              # OpenGL 3D renderer
├── hardware_runner.py       # IBM / IonQ hardware submission
├── benchmarker.py           # Method comparison (HF, CISD, CCSD(T), FCI, VQE)
├── requirements.txt
├── config.yaml              # Default configuration
└── tests/
    ├── test_molecule.py
    ├── test_hamiltonian.py
    ├── test_vqe.py
    ├── test_pes.py
    ├── test_orbital.py
    ├── test_isosurface.py
    └── test_renderer.py
\end{verbatim}

\subsection{Data Flow Diagram}\label{data-flow-diagram}

\begin{verbatim}
                        ┌──────────────┐
                        │ config.yaml  │
                        └──────┬───────┘
                               │
                        ┌──────▼───────┐
                        │   main.py    │   CLI / orchestration
                        └──────┬───────┘
                               │
              ┌────────────────┼────────────────────┐
              │                │                     │
     ┌────────▼───────┐ ┌─────▼──────────┐ ┌───────▼──────────┐
     │  molecule.py   │ │ pes_scanner.py │ │  benchmarker.py  │
     │ (geometry,     │ │ (bond length   │ │ (HF/CISD/CCSD(T) │
     │  basis set)    │ │  loop)         │ │  /FCI comparison) │
     └────────┬───────┘ └─────┬──────────┘ └───────┬──────────┘
              │               │                     │
     ┌────────▼───────┐      │                     │
     │ classical_hf.py│◄─────┘─────────────────────┘
     │ (PySCF: HF,    │
     │  integrals, MO)│
     └────────┬───────┘
              │
     ┌────────▼──────────────┐
     │ hamiltonian_builder.py│
     │ (JW/Parity mapping,  │
     │  QuTiP cross-check)  │
     └────────┬──────────────┘
              │
    ┌─────────┼──────────────────┐
    │         │                  │
┌───▼───┐ ┌──▼──────────┐ ┌─────▼─────────────┐
│ansatz │ │ vqe_solver  │ │ hardware_runner.py│
│ .py   │ │   .py       │ │ (IBM/IonQ device) │
└───┬───┘ └──┬──────────┘ └─────┬─────────────┘
    │        │                   │
    └────────┼───────────────────┘
             │
    ┌────────▼───────┐     ┌───────────────┐
    │ orbital_grid.py│────►│isosurface.py  │
    │ (3D MO grid)   │     │(marching cubes)│
    └────────────────┘     └───────┬───────┘
                                   │
                           ┌───────▼───────┐
                           │  renderer.py  │
                           │ (OpenGL:      │
                           │  isosurface + │
                           │  PES curve +  │
                           │  mol. geom.)  │
                           └───────────────┘
\end{verbatim}

\subsection{Module Specifications}\label{module-specifications}

\subsubsection{\texorpdfstring{1.
\texttt{molecule.py}}{1. molecule.py}}\label{molecule.py}

\textbf{Purpose:} Define molecular systems with geometry, basis set, and
charge/multiplicity.

This module provides two core data classes and a set of factory/utility
functions:

\begin{itemize}
\tightlist
\item
  \textbf{\texttt{Atom}} --- Holds an element symbol (e.g.,
  \texttt{\textquotesingle{}H\textquotesingle{}},
  \texttt{\textquotesingle{}O\textquotesingle{}}) and a Cartesian
  position vector \((x, y, z)\) in Ångströms.
\item
  \textbf{\texttt{MoleculeSpec}} --- Aggregates a list of \texttt{Atom}
  objects together with basis set name (default \texttt{sto-3g}), net
  charge (default 0), spin multiplicity \(2S+1\) (default 1 = singlet),
  and a human-readable name for logging.
\end{itemize}

{\def\LTcaptype{none} % do not increment counter
\begin{longtable}[]{@{}
  >{\raggedright\arraybackslash}p{(\linewidth - 6\tabcolsep) * \real{0.2083}}
  >{\raggedright\arraybackslash}p{(\linewidth - 6\tabcolsep) * \real{0.2708}}
  >{\raggedright\arraybackslash}p{(\linewidth - 6\tabcolsep) * \real{0.3333}}
  >{\raggedright\arraybackslash}p{(\linewidth - 6\tabcolsep) * \real{0.1875}}@{}}
\toprule\noalign{}
\begin{minipage}[b]{\linewidth}\raggedright
Function
\end{minipage} & \begin{minipage}[b]{\linewidth}\raggedright
Description
\end{minipage} & \begin{minipage}[b]{\linewidth}\raggedright
Key Parameters
\end{minipage} & \begin{minipage}[b]{\linewidth}\raggedright
Returns
\end{minipage} \\
\midrule\noalign{}
\endhead
\bottomrule\noalign{}
\endlastfoot
\texttt{create\_h2(bond\_length)} & Places two H atoms symmetrically
along the \(z\)-axis at \(\pm R/2\). & \(R\) in Å (default 0.735) &
\texttt{MoleculeSpec} for H₂ \\
\texttt{create\_lih(bond\_length)} & Places Li at the origin and H at
distance \(R\) along \(z\). & \(R\) in Å (default 1.595) &
\texttt{MoleculeSpec} for LiH \\
\texttt{create\_h2o(oh\_length,\ angle\_deg)} & Constructs a bent H₂O
geometry with O at the origin and two H atoms placed at bond length
\(R_{\text{OH}}\) and opening angle \(\alpha\). & \(R_{\text{OH}}\) in Å
(default 0.957), \(\alpha\) in degrees (default 104.5) &
\texttt{MoleculeSpec} for H₂O \\
\texttt{set\_bond\_length(mol,\ a,\ b,\ new\_length)} & Returns a deep
copy of \texttt{mol} with the distance between atoms \texttt{a} and
\texttt{b} set to \texttt{new\_length}, keeping the bond midpoint fixed.
The direction vector
\(\hat{\mathbf{d}} = (\mathbf{r}_b - \mathbf{r}_a)/\|\mathbf{r}_b - \mathbf{r}_a\|\)
is preserved while both atoms are displaced symmetrically. & Atom
indices, new \(R\) in Å & New \texttt{MoleculeSpec} \\
\end{longtable}
}

\begin{Shaded}
\begin{Highlighting}[]
\ImportTok{from}\NormalTok{ \_\_future\_\_ }\ImportTok{import}\NormalTok{ annotations}

\ImportTok{import}\NormalTok{ copy}
\ImportTok{from}\NormalTok{ dataclasses }\ImportTok{import}\NormalTok{ dataclass}

\ImportTok{import}\NormalTok{ numpy }\ImportTok{as}\NormalTok{ np}


\AttributeTok{@dataclass}
\KeywordTok{class}\NormalTok{ Atom:}
    \CommentTok{"""A single atom with element symbol and Cartesian position (Å)."""}
\NormalTok{    symbol: }\BuiltInTok{str}
\NormalTok{    position: np.ndarray  }\CommentTok{\# shape: (3,)}


\AttributeTok{@dataclass}
\KeywordTok{class}\NormalTok{ MoleculeSpec:}
    \CommentTok{"""Full specification of a molecular system."""}
\NormalTok{    atoms: }\BuiltInTok{list}\NormalTok{[Atom]}
\NormalTok{    basis: }\BuiltInTok{str} \OperatorTok{=} \StringTok{"sto{-}3g"}
\NormalTok{    charge: }\BuiltInTok{int} \OperatorTok{=} \DecValTok{0}
\NormalTok{    multiplicity: }\BuiltInTok{int} \OperatorTok{=} \DecValTok{1}
\NormalTok{    name: }\BuiltInTok{str} \OperatorTok{=} \StringTok{""}


\KeywordTok{def}\NormalTok{ create\_h2(bond\_length: }\BuiltInTok{float} \OperatorTok{=} \FloatTok{0.735}\NormalTok{) }\OperatorTok{{-}\textgreater{}}\NormalTok{ MoleculeSpec:}
    \CommentTok{"""Place two H atoms symmetrically along z at ±R/2."""}
\NormalTok{    ...}


\KeywordTok{def}\NormalTok{ create\_lih(bond\_length: }\BuiltInTok{float} \OperatorTok{=} \FloatTok{1.595}\NormalTok{) }\OperatorTok{{-}\textgreater{}}\NormalTok{ MoleculeSpec:}
    \CommentTok{"""Place Li at origin and H at distance R along z."""}
\NormalTok{    ...}


\KeywordTok{def}\NormalTok{ create\_h2o(}
\NormalTok{    oh\_length: }\BuiltInTok{float} \OperatorTok{=} \FloatTok{0.957}\NormalTok{,}
\NormalTok{    angle\_deg: }\BuiltInTok{float} \OperatorTok{=} \FloatTok{104.5}\NormalTok{,}
\NormalTok{) }\OperatorTok{{-}\textgreater{}}\NormalTok{ MoleculeSpec:}
    \CommentTok{"""Construct bent H₂O: O at origin, two H at R\_OH and angle α."""}
\NormalTok{    ...}


\KeywordTok{def}\NormalTok{ set\_bond\_length(}
\NormalTok{    mol: MoleculeSpec,}
\NormalTok{    atom\_idx\_a: }\BuiltInTok{int}\NormalTok{,}
\NormalTok{    atom\_idx\_b: }\BuiltInTok{int}\NormalTok{,}
\NormalTok{    new\_length: }\BuiltInTok{float}\NormalTok{,}
\NormalTok{) }\OperatorTok{{-}\textgreater{}}\NormalTok{ MoleculeSpec:}
    \CommentTok{"""Return a deep copy with bond a–b set to new\_length (Å).}

\CommentTok{    Preserves bond midpoint and direction unit vector.}
\CommentTok{    """}
\NormalTok{    ...}
\end{Highlighting}
\end{Shaded}

\emph{(Full implementation: \texttt{report2\_scripts/molecule.py})}

\subsubsection{\texorpdfstring{2.
\texttt{classical\_hf.py}}{2. classical\_hf.py}}\label{classical_hf.py}

\textbf{Purpose:} Interface to PySCF for Hartree-Fock calculations,
integral extraction, and MO coefficients.

This module wraps PySCF to provide all the classical quantum-chemistry
quantities needed downstream. It exposes a result container and four
functions:

\begin{itemize}
\tightlist
\item
  \textbf{\texttt{HFResult}} --- Data class bundling everything produced
  by a converged RHF calculation:

  \begin{itemize}
  \tightlist
  \item
    Total HF energy \(E_{\text{HF}}\) and nuclear repulsion energy
    \(E_{\text{nuc}}\) (both in Hartree).
  \item
    MO coefficient matrix \(\mathbf{C}\) of shape \((K, M)\) where \(K\)
    = number of AO basis functions and \(M\) = number of MOs.
  \item
    Orbital energies \(\{\epsilon_p\}\) of shape \((M,)\).
  \item
    One-electron integrals \(h_{pq}\) in the MO basis, shape \((M, M)\).
  \item
    Two-electron integrals \(h_{pqrs}\) in the MO basis, shape
    \((M, M, M, M)\).
  \item
    AO overlap matrix \(\mathbf{S}\) of shape \((K, K)\), electron count
    \(N_e\), orbital count \(M\), and AO labels.
  \end{itemize}
\end{itemize}

{\def\LTcaptype{none} % do not increment counter
\begin{longtable}[]{@{}
  >{\raggedright\arraybackslash}p{(\linewidth - 4\tabcolsep) * \real{0.2703}}
  >{\raggedright\arraybackslash}p{(\linewidth - 4\tabcolsep) * \real{0.4865}}
  >{\raggedright\arraybackslash}p{(\linewidth - 4\tabcolsep) * \real{0.2432}}@{}}
\toprule\noalign{}
\begin{minipage}[b]{\linewidth}\raggedright
Function
\end{minipage} & \begin{minipage}[b]{\linewidth}\raggedright
What it computes
\end{minipage} & \begin{minipage}[b]{\linewidth}\raggedright
Returns
\end{minipage} \\
\midrule\noalign{}
\endhead
\bottomrule\noalign{}
\endlastfoot
\texttt{run\_hartree\_fock(mol\_spec)} & Builds a PySCF \texttt{Mole}
object from the \texttt{MoleculeSpec}, runs RHF SCF, then extracts the
converged Fock matrix \(\mathbf{F}\), MO coefficients \(\mathbf{C}\),
orbital energies \(\{\epsilon_p\}\), and transforms the one- and
two-electron integrals to the MO basis via
\(h_{pq} = \mathbf{C}^T \mathbf{h}^{\text{core}} \mathbf{C}\) and the
four-index AO-to-MO transformation for \(h_{pqrs}\). &
\texttt{HFResult} \\
\texttt{get\_fci\_energy(mol\_spec)} & Runs RHF followed by PySCF's Full
CI solver to obtain the exact ground-state energy within the given
basis. & \(E_{\text{FCI}}\) (float, Hartree) \\
\texttt{get\_cisd\_energy(mol\_spec)} & Runs RHF + Configuration
Interaction Singles and Doubles. & \(E_{\text{CISD}}\) (float,
Hartree) \\
\texttt{get\_ccsd\_t\_energy(mol\_spec)} & Runs RHF + CCSD, then adds
the perturbative triples correction (T). & \(E_{\text{CCSD(T)}}\)
(float, Hartree) \\
\end{longtable}
}

\begin{Shaded}
\begin{Highlighting}[]
\ImportTok{from}\NormalTok{ \_\_future\_\_ }\ImportTok{import}\NormalTok{ annotations}

\ImportTok{from}\NormalTok{ dataclasses }\ImportTok{import}\NormalTok{ dataclass}

\ImportTok{import}\NormalTok{ numpy }\ImportTok{as}\NormalTok{ np}
\ImportTok{from}\NormalTok{ pyscf }\ImportTok{import}\NormalTok{ ao2mo, gto, scf, fci, ci, cc}


\AttributeTok{@dataclass}
\KeywordTok{class}\NormalTok{ HFResult:}
    \CommentTok{"""Results from a converged RHF calculation."""}
\NormalTok{    energy\_hf: }\BuiltInTok{float}                        \CommentTok{\# E\_HF (Hartree)}
\NormalTok{    energy\_nuc: }\BuiltInTok{float}                       \CommentTok{\# nuclear repulsion (Hartree)}
\NormalTok{    mo\_coefficients: np.ndarray             }\CommentTok{\# shape: (K, M) — AO{-}to{-}MO matrix}
\NormalTok{    mo\_energies: np.ndarray                 }\CommentTok{\# shape: (M,)}
\NormalTok{    one\_electron\_integrals: np.ndarray      }\CommentTok{\# h\_pq in MO basis, shape: (M, M)}
\NormalTok{    two\_electron\_integrals: np.ndarray      }\CommentTok{\# h\_pqrs in MO basis, shape: (M, M, M, M)}
\NormalTok{    overlap\_matrix: np.ndarray              }\CommentTok{\# S\_μν in AO basis, shape: (K, K)}
\NormalTok{    n\_electrons: }\BuiltInTok{int}
\NormalTok{    n\_orbitals: }\BuiltInTok{int}
\NormalTok{    ao\_labels: }\BuiltInTok{list}\NormalTok{[}\BuiltInTok{str}\NormalTok{]}


\KeywordTok{def}\NormalTok{ run\_hartree\_fock(mol\_spec: MoleculeSpec) }\OperatorTok{{-}\textgreater{}}\NormalTok{ HFResult:}
    \CommentTok{"""Build PySCF Mole, run RHF, extract MO{-}basis integrals.}

\CommentTok{    Transforms one{-}electron integrals via C\^{}T h\^{}core C and}
\CommentTok{    two{-}electron integrals via ao2mo.full.}
\CommentTok{    """}
\NormalTok{    ...}


\KeywordTok{def}\NormalTok{ get\_fci\_energy(mol\_spec: MoleculeSpec) }\OperatorTok{{-}\textgreater{}} \BuiltInTok{float}\NormalTok{:}
    \CommentTok{"""Run RHF + Full CI; return E\_FCI (Hartree)."""}
\NormalTok{    ...}


\KeywordTok{def}\NormalTok{ get\_cisd\_energy(mol\_spec: MoleculeSpec) }\OperatorTok{{-}\textgreater{}} \BuiltInTok{float}\NormalTok{:}
    \CommentTok{"""Run RHF + CISD; return E\_CISD (Hartree)."""}
\NormalTok{    ...}


\KeywordTok{def}\NormalTok{ get\_ccsd\_t\_energy(mol\_spec: MoleculeSpec) }\OperatorTok{{-}\textgreater{}} \BuiltInTok{float}\NormalTok{:}
    \CommentTok{"""Run RHF + CCSD(T); return E\_CCSD(T) (Hartree)."""}
\NormalTok{    ...}
\end{Highlighting}
\end{Shaded}

\emph{(Full implementation: \texttt{report2\_scripts/classical\_hf.py})}

\subsubsection{\texorpdfstring{3.
\texttt{hamiltonian\_builder.py}}{3. hamiltonian\_builder.py}}\label{hamiltonian_builder.py}

\textbf{Purpose:} Build the qubit Hamiltonian from molecular integrals
using Jordan-Wigner or Parity mapping. Cross-validate with QuTiP exact
diagonalisation.

This module converts classical molecular integrals into a qubit-operator
representation suitable for VQE, and provides an independent validation
path.

\begin{itemize}
\tightlist
\item
  \textbf{\texttt{QubitHamiltonian}} --- Data class holding the
  PennyLane Hamiltonian object, qubit count \(N_q\), number of Pauli
  terms, the mapping name
  (\texttt{\textquotesingle{}jordan\_wigner\textquotesingle{}} or
  \texttt{\textquotesingle{}parity\textquotesingle{}}), and an optional
  exact ground-state energy for validation.
\end{itemize}

{\def\LTcaptype{none} % do not increment counter
\begin{longtable}[]{@{}
  >{\raggedright\arraybackslash}p{(\linewidth - 2\tabcolsep) * \real{0.4348}}
  >{\raggedright\arraybackslash}p{(\linewidth - 2\tabcolsep) * \real{0.5652}}@{}}
\toprule\noalign{}
\begin{minipage}[b]{\linewidth}\raggedright
Function
\end{minipage} & \begin{minipage}[b]{\linewidth}\raggedright
Description
\end{minipage} \\
\midrule\noalign{}
\endhead
\bottomrule\noalign{}
\endlastfoot
\texttt{build\_hamiltonian(h\_pq,\ h\_pqrs,\ n\_electrons,\ E\_nuc,\ mapping,\ ...)}
& Assembles the second-quantised fermionic Hamiltonian
\(\hat{H} = \sum_{pq}h_{pq}\hat{a}_p^\dagger\hat{a}_q + \tfrac{1}{2}\sum_{pqrs}h_{pqrs}\hat{a}_p^\dagger\hat{a}_q^\dagger\hat{a}_s\hat{a}_r + E_{\text{nuc}}\)
from the supplied MO-basis integrals, then applies the chosen
fermion-to-qubit mapping (Jordan-Wigner or parity) via PennyLane.
Optionally restricts to an active space of the specified number of
active electrons and orbitals. Returns a \texttt{QubitHamiltonian}. \\
\texttt{validate\_with\_qutip(qubit\_hamiltonian)} & Converts the
PennyLane Hamiltonian to a QuTiP \texttt{Qobj}, performs exact
diagonalisation, and returns the ground-state energy \(E_0\). This
provides an independent check that the fermion-to-qubit mapping is
correct. \\
\texttt{pennylane\_to\_qutip(hamiltonian,\ n\_qubits)} & Maps each Pauli
string \(c_k \hat{P}_k\) in the PennyLane Hamiltonian to a tensor
product of QuTiP Pauli matrices (\(\hat{\sigma}_x\), \(\hat{\sigma}_y\),
\(\hat{\sigma}_z\), \(\hat{I}\)), accumulating into a single QuTiP
\texttt{Qobj} of dimension \((2^{N_q}, 2^{N_q})\). \\
\end{longtable}
}

\begin{Shaded}
\begin{Highlighting}[]
\ImportTok{from}\NormalTok{ \_\_future\_\_ }\ImportTok{import}\NormalTok{ annotations}

\ImportTok{from}\NormalTok{ dataclasses }\ImportTok{import}\NormalTok{ dataclass}

\ImportTok{import}\NormalTok{ numpy }\ImportTok{as}\NormalTok{ np}
\ImportTok{import}\NormalTok{ pennylane }\ImportTok{as}\NormalTok{ qml}
\ImportTok{import}\NormalTok{ qutip}


\AttributeTok{@dataclass}
\KeywordTok{class}\NormalTok{ QubitHamiltonian:}
    \CommentTok{"""Container for a qubit Hamiltonian and metadata."""}
\NormalTok{    hamiltonian: qml.Hamiltonian}
\NormalTok{    n\_qubits: }\BuiltInTok{int}
\NormalTok{    n\_paulis: }\BuiltInTok{int}
\NormalTok{    mapping: }\BuiltInTok{str}
\NormalTok{    fci\_energy: }\BuiltInTok{float} \OperatorTok{|} \VariableTok{None} \OperatorTok{=} \VariableTok{None}


\KeywordTok{def}\NormalTok{ build\_hamiltonian(}
\NormalTok{    one\_electron\_integrals: np.ndarray,}
\NormalTok{    two\_electron\_integrals: np.ndarray,}
\NormalTok{    n\_electrons: }\BuiltInTok{int}\NormalTok{,}
\NormalTok{    nuclear\_repulsion: }\BuiltInTok{float}\NormalTok{,}
\NormalTok{    mapping: }\BuiltInTok{str} \OperatorTok{=} \StringTok{"jordan\_wigner"}\NormalTok{,}
\NormalTok{    active\_electrons: }\BuiltInTok{int} \OperatorTok{|} \VariableTok{None} \OperatorTok{=} \VariableTok{None}\NormalTok{,}
\NormalTok{    active\_orbitals: }\BuiltInTok{int} \OperatorTok{|} \VariableTok{None} \OperatorTok{=} \VariableTok{None}\NormalTok{,}
\NormalTok{) }\OperatorTok{{-}\textgreater{}}\NormalTok{ QubitHamiltonian:}
    \CommentTok{"""Assemble fermionic Hamiltonian from MO integrals and map to qubits.}

\CommentTok{    Supports Jordan{-}Wigner and parity mappings via PennyLane.}
\CommentTok{    Optionally restricts to an active space.}
\CommentTok{    """}
\NormalTok{    ...}


\KeywordTok{def}\NormalTok{ validate\_with\_qutip(qubit\_hamiltonian: QubitHamiltonian) }\OperatorTok{{-}\textgreater{}} \BuiltInTok{float}\NormalTok{:}
    \CommentTok{"""Convert to QuTiP Qobj, diagonalise, return ground{-}state energy E₀."""}
\NormalTok{    ...}


\KeywordTok{def}\NormalTok{ pennylane\_to\_qutip(}
\NormalTok{    hamiltonian: qml.Hamiltonian,}
\NormalTok{    n\_qubits: }\BuiltInTok{int}\NormalTok{,}
\NormalTok{) }\OperatorTok{{-}\textgreater{}}\NormalTok{ qutip.Qobj:}
    \CommentTok{"""Map PennyLane Pauli Hamiltonian to a QuTiP Qobj of dim (2\^{}N, 2\^{}N)."""}
\NormalTok{    ...}
\end{Highlighting}
\end{Shaded}

\emph{(Full implementation:
\texttt{report2\_scripts/hamiltonian\_builder.py})}

\subsubsection{\texorpdfstring{4.
\texttt{ansatz.py}}{4. ansatz.py}}\label{ansatz.py}

\textbf{Purpose:} Construct UCCSD and hardware-efficient ansätze using
PennyLane.

This module provides two ansatz constructors, each returning a callable
circuit function and a parameter count.

\begin{quote}
\textbf{\texttt{build\_uccsd\_ansatz(n\_qubits,\ n\_electrons,\ hf\_state=None)}}
constructs the Unitary Coupled Cluster Singles and Doubles ansatz. It
first enumerates all spin-allowed single excitations \(\{(i \to a)\}\)
and double excitations \(\{(i,j \to a,b)\}\) using PennyLane's
\texttt{qchem.excitations}. The UCCSD unitary is:
\[|\Psi(\boldsymbol{\theta})\rangle = e^{\hat{T}(\boldsymbol{\theta}) - \hat{T}^\dagger(\boldsymbol{\theta})}|\Phi_{\text{HF}}\rangle\]
where \(\hat{T} = \hat{T}_1 + \hat{T}_2\) with
\(\hat{T}_1 = \sum_{i,a}\theta_i^a \hat{a}_a^\dagger \hat{a}_i\) and
\(\hat{T}_2 = \sum_{i<j,a<b}\theta_{ij}^{ab}\hat{a}_a^\dagger\hat{a}_b^\dagger\hat{a}_j\hat{a}_i\).
PennyLane's \texttt{qml.UCCSD} (or individual \texttt{SingleExcitation}
/ \texttt{DoubleExcitation} gates) compiles this into a gate sequence
via first-order Trotterisation. The HF reference state fills the first
\(N_e\) spin-orbitals (or uses a user-supplied binary vector). Returns
\texttt{(ansatz\_fn,\ n\_params,\ excitation\_list)}.
\end{quote}

\begin{quote}
\textbf{\texttt{build\_hardware\_efficient\_ansatz(n\_qubits,\ n\_layers=2,\ entangler=\textquotesingle{}CNOT\textquotesingle{})}}
builds a hardware-efficient ansatz consisting of \(L\) layers, each
comprising \(R_Y(\theta)\) and \(R_Z(\phi)\) rotations on every qubit
followed by a ladder of CNOT (or CZ) entangling gates on
nearest-neighbour pairs. The total parameter count is \(2 N_q L\).
Returns \texttt{(ansatz\_fn,\ n\_params)}.
\end{quote}

\begin{Shaded}
\begin{Highlighting}[]
\ImportTok{from}\NormalTok{ \_\_future\_\_ }\ImportTok{import}\NormalTok{ annotations}

\ImportTok{from}\NormalTok{ typing }\ImportTok{import}\NormalTok{ Callable}

\ImportTok{import}\NormalTok{ numpy }\ImportTok{as}\NormalTok{ np}
\ImportTok{import}\NormalTok{ pennylane }\ImportTok{as}\NormalTok{ qml}


\KeywordTok{def}\NormalTok{ build\_uccsd\_ansatz(}
\NormalTok{    n\_qubits: }\BuiltInTok{int}\NormalTok{,}
\NormalTok{    n\_electrons: }\BuiltInTok{int}\NormalTok{,}
\NormalTok{    hf\_state: np.ndarray }\OperatorTok{|} \VariableTok{None} \OperatorTok{=} \VariableTok{None}\NormalTok{,}
\NormalTok{) }\OperatorTok{{-}\textgreater{}} \BuiltInTok{tuple}\NormalTok{[Callable, }\BuiltInTok{int}\NormalTok{, }\BuiltInTok{list}\NormalTok{]:}
    \CommentTok{"""Construct UCCSD ansatz from PennyLane qchem excitations.}

\CommentTok{    Returns (ansatz\_fn, n\_params, excitation\_list).}
\CommentTok{    ansatz\_fn accepts a parameter array of shape (n\_params,).}
\CommentTok{    """}
\NormalTok{    ...}


\KeywordTok{def}\NormalTok{ build\_hardware\_efficient\_ansatz(}
\NormalTok{    n\_qubits: }\BuiltInTok{int}\NormalTok{,}
\NormalTok{    n\_layers: }\BuiltInTok{int} \OperatorTok{=} \DecValTok{2}\NormalTok{,}
\NormalTok{    entangler: }\BuiltInTok{str} \OperatorTok{=} \StringTok{"CNOT"}\NormalTok{,}
\NormalTok{) }\OperatorTok{{-}\textgreater{}} \BuiltInTok{tuple}\NormalTok{[Callable, }\BuiltInTok{int}\NormalTok{]:}
    \CommentTok{"""Build L{-}layer hardware{-}efficient ansatz (RY, RZ + entangler).}

\CommentTok{    Total parameters: 2 * n\_qubits * n\_layers.}
\CommentTok{    Returns (ansatz\_fn, n\_params).}
\CommentTok{    """}
\NormalTok{    ...}
\end{Highlighting}
\end{Shaded}

\emph{(Full implementation: \texttt{report2\_scripts/ansatz.py})}

\subsubsection{\texorpdfstring{5.
\texttt{vqe\_solver.py}}{5. vqe\_solver.py}}\label{vqe_solver.py}

\textbf{Purpose:} VQE optimisation loop with multiple optimiser
backends.

\begin{itemize}
\tightlist
\item
  \textbf{\texttt{VQEResult}} --- Data class storing the best energy
  \(E^*\) (Hartree), optimal parameters \(\boldsymbol{\theta}^*\), full
  trajectory \((\boldsymbol{\theta}^{(k)}, E^{(k)})\) at each iteration,
  total iteration count, convergence flag, and gradient norms.
\end{itemize}

\begin{quote}
\textbf{\texttt{run\_vqe(hamiltonian,\ ansatz\_fn,\ n\_params,\ n\_qubits,\ ...)}}
executes the VQE optimisation loop. A PennyLane QNode wraps the ansatz
and evaluates
\(E(\boldsymbol{\theta}) = \langle\psi(\boldsymbol{\theta})|\hat{H}|\psi(\boldsymbol{\theta})\rangle\).
The solver supports four optimiser backends:

\begin{itemize}
\tightlist
\item
  \textbf{Adam / GradientDescent}: PennyLane built-in optimisers using
  automatic differentiation (backprop or parameter-shift).
\item
  \textbf{COBYLA}: SciPy's derivative-free constrained optimiser.
\item
  \textbf{SPSA} (Simultaneous Perturbation Stochastic Approximation):
  estimates the gradient from two function evaluations per step via
  random perturbation
  \(\hat{g}_k = \frac{f(\boldsymbol{\theta}+c_k\boldsymbol{\Delta}) - f(\boldsymbol{\theta}-c_k\boldsymbol{\Delta})}{2c_k\boldsymbol{\Delta}}\)
  with decaying step sizes \(a_k = a_0/(k+A)^\alpha\) and
  \(c_k = c_0/k^\gamma\).
\end{itemize}

Convergence is declared when \(|E^{(k)} - E^{(k-1)}| < \varepsilon\).
Initial parameters default to zero with a small random perturbation to
break symmetry. An optional callback is invoked each iteration with
\texttt{(iteration,\ params,\ energy)}.
\end{quote}

\begin{Shaded}
\begin{Highlighting}[]
\ImportTok{from}\NormalTok{ \_\_future\_\_ }\ImportTok{import}\NormalTok{ annotations}

\ImportTok{from}\NormalTok{ dataclasses }\ImportTok{import}\NormalTok{ dataclass}
\ImportTok{from}\NormalTok{ typing }\ImportTok{import}\NormalTok{ Callable}

\ImportTok{import}\NormalTok{ numpy }\ImportTok{as}\NormalTok{ np}
\ImportTok{import}\NormalTok{ pennylane }\ImportTok{as}\NormalTok{ qml}


\AttributeTok{@dataclass}
\KeywordTok{class}\NormalTok{ VQEResult:}
    \CommentTok{"""Container for VQE optimisation results."""}
\NormalTok{    optimal\_energy: }\BuiltInTok{float}
\NormalTok{    optimal\_params: np.ndarray}
\NormalTok{    trajectory: }\BuiltInTok{list}\NormalTok{[}\BuiltInTok{tuple}\NormalTok{[np.ndarray, }\BuiltInTok{float}\NormalTok{]]}
\NormalTok{    n\_iterations: }\BuiltInTok{int}
\NormalTok{    converged: }\BuiltInTok{bool}
\NormalTok{    gradient\_norms: }\BuiltInTok{list}\NormalTok{[}\BuiltInTok{float}\NormalTok{]}


\KeywordTok{def}\NormalTok{ run\_vqe(}
\NormalTok{    hamiltonian: qml.Hamiltonian,}
\NormalTok{    ansatz\_fn: Callable,}
\NormalTok{    n\_params: }\BuiltInTok{int}\NormalTok{,}
\NormalTok{    n\_qubits: }\BuiltInTok{int}\NormalTok{,}
\NormalTok{    init\_params: np.ndarray }\OperatorTok{|} \VariableTok{None} \OperatorTok{=} \VariableTok{None}\NormalTok{,}
\NormalTok{    optimizer: }\BuiltInTok{str} \OperatorTok{=} \StringTok{"Adam"}\NormalTok{,}
\NormalTok{    learning\_rate: }\BuiltInTok{float} \OperatorTok{=} \FloatTok{0.05}\NormalTok{,}
\NormalTok{    maxiter: }\BuiltInTok{int} \OperatorTok{=} \DecValTok{200}\NormalTok{,}
\NormalTok{    tol: }\BuiltInTok{float} \OperatorTok{=} \FloatTok{1e{-}6}\NormalTok{,}
\NormalTok{    device\_name: }\BuiltInTok{str} \OperatorTok{=} \StringTok{"default.qubit"}\NormalTok{,}
\NormalTok{    diff\_method: }\BuiltInTok{str} \OperatorTok{=} \StringTok{"backprop"}\NormalTok{,}
\NormalTok{    callback: Callable }\OperatorTok{|} \VariableTok{None} \OperatorTok{=} \VariableTok{None}\NormalTok{,}
\NormalTok{) }\OperatorTok{{-}\textgreater{}}\NormalTok{ VQEResult:}
    \CommentTok{"""Run the VQE optimisation loop.}

\CommentTok{    Supports Adam, GradientDescent, COBYLA, and SPSA backends.}
\CommentTok{    Convergence declared when |ΔE| \textless{} tol.}
\CommentTok{    """}
\NormalTok{    ...}
\end{Highlighting}
\end{Shaded}

\emph{(Full implementation: \texttt{report2\_scripts/vqe\_solver.py})}

\subsubsection{\texorpdfstring{6.
\texttt{pes\_scanner.py}}{6. pes\_scanner.py}}\label{pes_scanner.py}

\textbf{Purpose:} Sweep over bond lengths, run VQE at each geometry,
collect PES data.

\begin{itemize}
\tightlist
\item
  \textbf{\texttt{PESData}} --- Data class holding arrays of bond
  lengths (Å), VQE energies, HF energies, FCI energies, and optionally
  CISD and CCSD(T) energies, plus the list of optimal VQE parameter
  vectors at each geometry.
\end{itemize}

\begin{quote}
\textbf{\texttt{scan\_pes(molecule\_factory,\ bond\_lengths,\ ...)}}
iterates over an array of bond lengths \(\{R_1, \ldots, R_N\}\). At each
\(R_k\) it:

\begin{enumerate}
\def\labelenumi{\arabic{enumi}.}
\tightlist
\item
  Calls \texttt{molecule\_factory(R\_k)} to build the geometry.
\item
  Runs Hartree-Fock via \texttt{run\_hartree\_fock} to extract
  integrals.
\item
  Builds the qubit Hamiltonian via \texttt{build\_hamiltonian}.
\item
  Constructs the chosen ansatz (UCCSD or hardware-efficient).
\item
  Runs VQE, optionally \textbf{warm-starting} from the previous
  geometry's optimal parameters \(\boldsymbol{\theta}^*_{k-1}\) to
  accelerate convergence along the scan.
\item
  Optionally computes HF, CISD, CCSD(T), and FCI energies via
  \texttt{classical\_hf} for comparison.
\end{enumerate}

The warm-start strategy exploits the fact that molecular orbitals change
smoothly with geometry, so \(\boldsymbol{\theta}^*_{k-1}\) is a good
initial guess for \(\boldsymbol{\theta}^*_k\).
\end{quote}

\begin{Shaded}
\begin{Highlighting}[]
\ImportTok{from}\NormalTok{ \_\_future\_\_ }\ImportTok{import}\NormalTok{ annotations}

\ImportTok{from}\NormalTok{ dataclasses }\ImportTok{import}\NormalTok{ dataclass}
\ImportTok{from}\NormalTok{ typing }\ImportTok{import}\NormalTok{ Callable}

\ImportTok{import}\NormalTok{ numpy }\ImportTok{as}\NormalTok{ np}


\AttributeTok{@dataclass}
\KeywordTok{class}\NormalTok{ PESData:}
    \CommentTok{"""Container for PES scan results."""}
\NormalTok{    bond\_lengths: np.ndarray}
\NormalTok{    energies\_vqe: np.ndarray}
\NormalTok{    energies\_hf: np.ndarray}
\NormalTok{    energies\_fci: np.ndarray}
\NormalTok{    energies\_cisd: np.ndarray }\OperatorTok{|} \VariableTok{None} \OperatorTok{=} \VariableTok{None}
\NormalTok{    energies\_ccsd\_t: np.ndarray }\OperatorTok{|} \VariableTok{None} \OperatorTok{=} \VariableTok{None}
\NormalTok{    vqe\_params: }\BuiltInTok{list}\NormalTok{[np.ndarray] }\OperatorTok{|} \VariableTok{None} \OperatorTok{=} \VariableTok{None}


\KeywordTok{def}\NormalTok{ scan\_pes(}
\NormalTok{    molecule\_factory: Callable[[}\BuiltInTok{float}\NormalTok{], MoleculeSpec],}
\NormalTok{    bond\_lengths: np.ndarray,}
\NormalTok{    basis: }\BuiltInTok{str} \OperatorTok{=} \StringTok{"sto{-}3g"}\NormalTok{,}
\NormalTok{    mapping: }\BuiltInTok{str} \OperatorTok{=} \StringTok{"jordan\_wigner"}\NormalTok{,}
\NormalTok{    ansatz\_type: }\BuiltInTok{str} \OperatorTok{=} \StringTok{"uccsd"}\NormalTok{,}
\NormalTok{    optimizer: }\BuiltInTok{str} \OperatorTok{=} \StringTok{"Adam"}\NormalTok{,}
\NormalTok{    maxiter: }\BuiltInTok{int} \OperatorTok{=} \DecValTok{200}\NormalTok{,}
\NormalTok{    use\_prev\_params: }\BuiltInTok{bool} \OperatorTok{=} \VariableTok{True}\NormalTok{,}
\NormalTok{    compute\_classical: }\BuiltInTok{bool} \OperatorTok{=} \VariableTok{True}\NormalTok{,}
\NormalTok{    callback: Callable }\OperatorTok{|} \VariableTok{None} \OperatorTok{=} \VariableTok{None}\NormalTok{,}
\NormalTok{) }\OperatorTok{{-}\textgreater{}}\NormalTok{ PESData:}
    \CommentTok{"""Sweep bond lengths, run VQE at each geometry, collect PES data.}

\CommentTok{    Warm{-}starts from previous geometry\textquotesingle{}s optimal parameters when}
\CommentTok{    use\_prev\_params is True and parameter count is unchanged.}
\CommentTok{    """}
\NormalTok{    ...}
\end{Highlighting}
\end{Shaded}

\emph{(Full implementation: \texttt{report2\_scripts/pes\_scanner.py})}

\subsubsection{\texorpdfstring{7.
\texttt{orbital\_grid.py}}{7. orbital\_grid.py}}\label{orbital_grid.py}

\textbf{Purpose:} Evaluate molecular orbital wavefunctions on a 3D
Cartesian grid.

\begin{itemize}
\tightlist
\item
  \textbf{\texttt{OrbitalVolumeData}} --- Data class holding the 3D
  scalar field array of shape \((N_x, N_y, N_z)\), grid origin, grid
  spacing (all in Bohr), grid shape, and the MO index.
\end{itemize}

\begin{quote}
\textbf{\texttt{evaluate\_orbital\_on\_grid(mo\_coefficients,\ orbital\_index,\ atom\_coords,\ atom\_symbols,\ basis\_name,\ grid\_extent=5.0,\ grid\_spacing=0.1)}}
computes the MO wavefunction \(\phi_p(\mathbf{r})\) on a uniform
Cartesian grid. The grid extends \(\pm\)\texttt{grid\_extent} Bohr from
the molecular centre in each direction with spacing
\texttt{grid\_spacing}. At each grid point \(\mathbf{r}_g\), PySCF's
\texttt{eval\_gto} evaluates all AO basis functions
\(\{\chi_\mu(\mathbf{r}_g)\}\) efficiently (exploiting Gaussian
sparsity), then contracts with the \(p\)-th column of the MO coefficient
matrix:
\[\phi_p(\mathbf{r}_g) = \sum_{\mu} C_{\mu p}\, \chi_\mu(\mathbf{r}_g)\]
The result is reshaped into a 3D volume of shape \((N_x, N_y, N_z)\).
\end{quote}

\begin{quote}
\textbf{\texttt{evaluate\_electron\_density(mo\_coefficients,\ n\_occupied,\ ...)}}
computes the total RHF electron density on the same grid type by summing
over occupied spatial orbitals:
\[\rho(\mathbf{r}) = \sum_{i=0}^{N_{\text{occ}}-1} 2\,|\phi_i(\mathbf{r})|^2\]
The factor of 2 accounts for double occupation in RHF. Returns an
\texttt{OrbitalVolumeData} with \texttt{orbital\_index\ =\ -1} to
indicate a density rather than a single orbital.
\end{quote}

\begin{Shaded}
\begin{Highlighting}[]
\ImportTok{from}\NormalTok{ \_\_future\_\_ }\ImportTok{import}\NormalTok{ annotations}

\ImportTok{from}\NormalTok{ dataclasses }\ImportTok{import}\NormalTok{ dataclass}

\ImportTok{import}\NormalTok{ numpy }\ImportTok{as}\NormalTok{ np}
\ImportTok{from}\NormalTok{ pyscf }\ImportTok{import}\NormalTok{ gto}


\AttributeTok{@dataclass}
\KeywordTok{class}\NormalTok{ OrbitalVolumeData:}
    \CommentTok{"""Volumetric data for a molecular orbital on a 3D grid."""}
\NormalTok{    values: np.ndarray           }\CommentTok{\# shape: (N\_x, N\_y, N\_z)}
\NormalTok{    grid\_origin: }\BuiltInTok{tuple}\NormalTok{[}\BuiltInTok{float}\NormalTok{, }\BuiltInTok{float}\NormalTok{, }\BuiltInTok{float}\NormalTok{]    }\CommentTok{\# in Bohr}
\NormalTok{    grid\_spacing: }\BuiltInTok{tuple}\NormalTok{[}\BuiltInTok{float}\NormalTok{, }\BuiltInTok{float}\NormalTok{, }\BuiltInTok{float}\NormalTok{]   }\CommentTok{\# in Bohr}
\NormalTok{    grid\_shape: }\BuiltInTok{tuple}\NormalTok{[}\BuiltInTok{int}\NormalTok{, }\BuiltInTok{int}\NormalTok{, }\BuiltInTok{int}\NormalTok{]}
\NormalTok{    orbital\_index: }\BuiltInTok{int}


\KeywordTok{def}\NormalTok{ evaluate\_orbital\_on\_grid(}
\NormalTok{    mo\_coefficients: np.ndarray,}
\NormalTok{    orbital\_index: }\BuiltInTok{int}\NormalTok{,}
\NormalTok{    atom\_coords: np.ndarray,}
\NormalTok{    atom\_symbols: }\BuiltInTok{list}\NormalTok{[}\BuiltInTok{str}\NormalTok{],}
\NormalTok{    basis\_name: }\BuiltInTok{str}\NormalTok{,}
\NormalTok{    grid\_extent: }\BuiltInTok{float} \OperatorTok{=} \FloatTok{5.0}\NormalTok{,}
\NormalTok{    grid\_spacing: }\BuiltInTok{float} \OperatorTok{=} \FloatTok{0.1}\NormalTok{,}
\NormalTok{) }\OperatorTok{{-}\textgreater{}}\NormalTok{ OrbitalVolumeData:}
    \CommentTok{"""Evaluate MO φ\_p(r) on a uniform 3D Cartesian grid.}

\CommentTok{    Uses PySCF eval\_gto for efficient AO evaluation, then}
\CommentTok{    contracts with MO coefficient column p.}
\CommentTok{    """}
\NormalTok{    ...}


\KeywordTok{def}\NormalTok{ evaluate\_electron\_density(}
\NormalTok{    mo\_coefficients: np.ndarray,}
\NormalTok{    n\_occupied: }\BuiltInTok{int}\NormalTok{,}
\NormalTok{    atom\_coords: np.ndarray,}
\NormalTok{    atom\_symbols: }\BuiltInTok{list}\NormalTok{[}\BuiltInTok{str}\NormalTok{],}
\NormalTok{    basis\_name: }\BuiltInTok{str}\NormalTok{,}
\NormalTok{    grid\_extent: }\BuiltInTok{float} \OperatorTok{=} \FloatTok{5.0}\NormalTok{,}
\NormalTok{    grid\_spacing: }\BuiltInTok{float} \OperatorTok{=} \FloatTok{0.1}\NormalTok{,}
\NormalTok{) }\OperatorTok{{-}\textgreater{}}\NormalTok{ OrbitalVolumeData:}
    \CommentTok{"""Compute total RHF electron density ρ(r) = Σ 2|φ\_i(r)|².}

\CommentTok{    Returns OrbitalVolumeData with orbital\_index = {-}1.}
\CommentTok{    """}
\NormalTok{    ...}
\end{Highlighting}
\end{Shaded}

\emph{(Full implementation: \texttt{report2\_scripts/orbital\_grid.py})}

\subsubsection{\texorpdfstring{8.
\texttt{isosurface.py}}{8. isosurface.py}}\label{isosurface.py}

\textbf{Purpose:} Marching cubes isosurface extraction from volumetric
data.

\begin{itemize}
\tightlist
\item
  \textbf{\texttt{IsosurfaceMesh}} --- Data class holding a triangle
  mesh: vertex positions \((N_v, 3)\), vertex normals \((N_v, 3)\), face
  index triples \((N_f, 3)\), and a per-vertex sign array (\(+1\) for
  positive lobe, \(-1\) for negative lobe) used for colouring.
\end{itemize}

\begin{quote}
\textbf{\texttt{marching\_cubes(volume,\ isovalue,\ grid\_origin,\ grid\_spacing)}}
extracts the isosurface \(S_c = \{\mathbf{r} : \phi(\mathbf{r}) = c\}\)
from a 3D scalar field using scikit-image's implementation of the
Lorensen \& Cline (1987) algorithm. Each voxel's 8 corners are
classified as inside/outside the isosurface, yielding an 8-bit index
into a lookup table of 256 triangulation cases (15 unique topologies).
Edge-crossing positions are linearly interpolated, and vertex normals
are computed from the scalar field gradient. Vertices are translated to
world coordinates by adding the grid origin.
\end{quote}

\begin{quote}
\textbf{\texttt{extract\_dual\_isosurface(volume,\ isovalue,\ grid\_origin,\ grid\_spacing)}}
calls \texttt{marching\_cubes} twice --- once at \(+c\) and once at
\(-c\) --- then concatenates the two meshes. Vertices from the positive
isosurface are tagged with sign \(= +1\); negative with sign \(= -1\).
Face indices of the negative mesh are offset by the vertex count of the
positive mesh before concatenation.
\end{quote}

\begin{Shaded}
\begin{Highlighting}[]
\ImportTok{from}\NormalTok{ \_\_future\_\_ }\ImportTok{import}\NormalTok{ annotations}

\ImportTok{from}\NormalTok{ dataclasses }\ImportTok{import}\NormalTok{ dataclass}

\ImportTok{import}\NormalTok{ numpy }\ImportTok{as}\NormalTok{ np}
\ImportTok{from}\NormalTok{ skimage.measure }\ImportTok{import}\NormalTok{ marching\_cubes }\ImportTok{as}\NormalTok{ sk\_marching\_cubes}


\AttributeTok{@dataclass}
\KeywordTok{class}\NormalTok{ IsosurfaceMesh:}
    \CommentTok{"""Triangle mesh from isosurface extraction."""}
\NormalTok{    vertices: np.ndarray    }\CommentTok{\# shape: (N\_v, 3)}
\NormalTok{    normals: np.ndarray     }\CommentTok{\# shape: (N\_v, 3)}
\NormalTok{    faces: np.ndarray       }\CommentTok{\# shape: (N\_f, 3), uint32}
\NormalTok{    signs: np.ndarray       }\CommentTok{\# shape: (N\_v,), +1 or {-}1}


\KeywordTok{def}\NormalTok{ marching\_cubes(}
\NormalTok{    volume: np.ndarray,}
\NormalTok{    isovalue: }\BuiltInTok{float}\NormalTok{,}
\NormalTok{    grid\_origin: }\BuiltInTok{tuple}\NormalTok{[}\BuiltInTok{float}\NormalTok{, }\BuiltInTok{float}\NormalTok{, }\BuiltInTok{float}\NormalTok{],}
\NormalTok{    grid\_spacing: }\BuiltInTok{tuple}\NormalTok{[}\BuiltInTok{float}\NormalTok{, }\BuiltInTok{float}\NormalTok{, }\BuiltInTok{float}\NormalTok{],}
\NormalTok{) }\OperatorTok{{-}\textgreater{}}\NormalTok{ IsosurfaceMesh:}
    \CommentTok{"""Extract isosurface \{r : φ(r) = c\} via scikit{-}image marching cubes.}

\CommentTok{    Vertices are translated to world coordinates by adding grid\_origin.}
\CommentTok{    """}
\NormalTok{    ...}


\KeywordTok{def}\NormalTok{ extract\_dual\_isosurface(}
\NormalTok{    volume: np.ndarray,}
\NormalTok{    isovalue: }\BuiltInTok{float}\NormalTok{,}
\NormalTok{    grid\_origin: }\BuiltInTok{tuple}\NormalTok{[}\BuiltInTok{float}\NormalTok{, }\BuiltInTok{float}\NormalTok{, }\BuiltInTok{float}\NormalTok{],}
\NormalTok{    grid\_spacing: }\BuiltInTok{tuple}\NormalTok{[}\BuiltInTok{float}\NormalTok{, }\BuiltInTok{float}\NormalTok{, }\BuiltInTok{float}\NormalTok{],}
\NormalTok{) }\OperatorTok{{-}\textgreater{}}\NormalTok{ IsosurfaceMesh:}
    \CommentTok{"""Extract +c and {-}c isosurfaces; concatenate with sign tags.}

\CommentTok{    Positive{-}lobe vertices: sign = +1. Negative: sign = {-}1.}
\CommentTok{    Face indices of the negative mesh are offset by |V\_pos|.}
\CommentTok{    """}
\NormalTok{    ...}
\end{Highlighting}
\end{Shaded}

\emph{(Full implementation: \texttt{report2\_scripts/isosurface.py})}

\subsubsection{\texorpdfstring{9.
\texttt{renderer.py}}{9. renderer.py}}\label{renderer.py}

\textbf{Purpose:} OpenGL renderer for orbital isosurfaces, PES curves,
and molecular geometry.

This module provides the full 3D rendering pipeline using GLFW + OpenGL
3.3 Core Profile.

{\def\LTcaptype{none} % do not increment counter
\begin{longtable}[]{@{}
  >{\raggedright\arraybackslash}p{(\linewidth - 2\tabcolsep) * \real{0.5806}}
  >{\raggedright\arraybackslash}p{(\linewidth - 2\tabcolsep) * \real{0.4194}}@{}}
\toprule\noalign{}
\begin{minipage}[b]{\linewidth}\raggedright
Function / Class
\end{minipage} & \begin{minipage}[b]{\linewidth}\raggedright
Description
\end{minipage} \\
\midrule\noalign{}
\endhead
\bottomrule\noalign{}
\endlastfoot
\texttt{init\_opengl\_context(width,\ height,\ title)} & Initialises a
GLFW window with an OpenGL 3.3 core-profile context, enables depth
testing, 4× MSAA, and alpha blending. Returns the GLFW window handle. \\
\texttt{create\_orbital\_vao(mesh)} & Uploads an \texttt{IsosurfaceMesh}
to the GPU. Vertex layout:
\texttt{{[}x,\ y,\ z,\ nx,\ ny,\ nz,\ sign{]}} (7 float32 values, 28
bytes per vertex). Position at attribute location 0, normal at location
1, sign at location 2. Uses \texttt{GL\_DYNAMIC\_DRAW} since the mesh
changes with bond length. Returns
\texttt{(VAO,\ VBO,\ EBO,\ n\_indices)}. \\
\texttt{update\_orbital\_vao(vbo,\ ebo,\ mesh)} & Orphans the old buffer
and uploads new mesh data to the same VBO/EBO --- a double-buffering
strategy that avoids GPU sync stalls. Returns the new index count. \\
\texttt{create\_pes\_curve\_vao(bond\_lengths,\ energies\_dict)} &
Creates a VAO for rendering multiple PES curves as coloured line strips.
Each method (HF, CISD, VQE, FCI, etc.) maps to a distinct RGBA colour.
Vertices store \texttt{{[}R,\ E,\ 0,\ r,\ g,\ b,\ a{]}}. Returns
\texttt{(VAO,\ VBO)}. \\
\texttt{draw\_molecule(atom\_positions,\ atom\_symbols,\ bonds,\ shader\_program)}
& Renders atoms as spheres and bonds as cylinders using instanced
drawing of unit meshes. Atom colours follow the CPK convention (H =
white, O = red, Li = purple, etc.) and radii are scaled van der Waals
radii (\(\times 0.3\)). \\
\texttt{ArcballCamera} & Camera class supporting arcball rotation
(left-drag), zoom (scroll), and heading/pitch angles. Computes view
(\(\mathbf{V} = \text{lookAt}\)) and perspective (\(\mathbf{P}\))
matrices for the shader pipeline. \\
\end{longtable}
}

\begin{Shaded}
\begin{Highlighting}[]
\ImportTok{from}\NormalTok{ \_\_future\_\_ }\ImportTok{import}\NormalTok{ annotations}

\ImportTok{import}\NormalTok{ ctypes}

\ImportTok{import}\NormalTok{ glfw}
\ImportTok{import}\NormalTok{ numpy }\ImportTok{as}\NormalTok{ np}
\ImportTok{from}\NormalTok{ OpenGL.GL }\ImportTok{import} \OperatorTok{*}


\KeywordTok{def}\NormalTok{ init\_opengl\_context(}
\NormalTok{    width: }\BuiltInTok{int} \OperatorTok{=} \DecValTok{1400}\NormalTok{,}
\NormalTok{    height: }\BuiltInTok{int} \OperatorTok{=} \DecValTok{900}\NormalTok{,}
\NormalTok{    title: }\BuiltInTok{str} \OperatorTok{=} \StringTok{"VQE Molecular Explorer"}\NormalTok{,}
\NormalTok{) }\OperatorTok{{-}\textgreater{}} \BuiltInTok{object}\NormalTok{:}
    \CommentTok{"""Init GLFW window with OpenGL 3.3 core profile, depth, MSAA, blending."""}
\NormalTok{    ...}


\KeywordTok{def}\NormalTok{ create\_orbital\_vao(}
\NormalTok{    mesh: IsosurfaceMesh,}
\NormalTok{) }\OperatorTok{{-}\textgreater{}} \BuiltInTok{tuple}\NormalTok{[}\BuiltInTok{int}\NormalTok{, }\BuiltInTok{int}\NormalTok{, }\BuiltInTok{int}\NormalTok{, }\BuiltInTok{int}\NormalTok{]:}
    \CommentTok{"""Upload IsosurfaceMesh to GPU.}

\CommentTok{    Vertex layout: [x, y, z, nx, ny, nz, sign] — 7 floats, 28 bytes.}
\CommentTok{    Returns (VAO, VBO, EBO, n\_indices).}
\CommentTok{    """}
\NormalTok{    ...}


\KeywordTok{def}\NormalTok{ update\_orbital\_vao(}
\NormalTok{    vbo: }\BuiltInTok{int}\NormalTok{,}
\NormalTok{    ebo: }\BuiltInTok{int}\NormalTok{,}
\NormalTok{    mesh: IsosurfaceMesh,}
\NormalTok{) }\OperatorTok{{-}\textgreater{}} \BuiltInTok{int}\NormalTok{:}
    \CommentTok{"""Orphan{-}and{-}reupload new mesh data to existing VBO/EBO.}

\CommentTok{    Returns new index count.}
\CommentTok{    """}
\NormalTok{    ...}


\KeywordTok{def}\NormalTok{ create\_pes\_curve\_vao(}
\NormalTok{    bond\_lengths: np.ndarray,}
\NormalTok{    energies\_dict: }\BuiltInTok{dict}\NormalTok{[}\BuiltInTok{str}\NormalTok{, np.ndarray],}
\NormalTok{) }\OperatorTok{{-}\textgreater{}} \BuiltInTok{tuple}\NormalTok{[}\BuiltInTok{int}\NormalTok{, }\BuiltInTok{int}\NormalTok{]:}
    \CommentTok{"""Create VAO for PES curves as coloured line strips.}

\CommentTok{    Vertex layout: [R, E, 0, r, g, b, a] — 7 floats per vertex.}
\CommentTok{    Returns (VAO, VBO).}
\CommentTok{    """}
\NormalTok{    ...}


\KeywordTok{def}\NormalTok{ draw\_molecule(}
\NormalTok{    atom\_positions: np.ndarray,}
\NormalTok{    atom\_symbols: }\BuiltInTok{list}\NormalTok{[}\BuiltInTok{str}\NormalTok{],}
\NormalTok{    bonds: }\BuiltInTok{list}\NormalTok{[}\BuiltInTok{tuple}\NormalTok{[}\BuiltInTok{int}\NormalTok{, }\BuiltInTok{int}\NormalTok{]],}
\NormalTok{    shader\_program: }\BuiltInTok{int}\NormalTok{,}
\NormalTok{) }\OperatorTok{{-}\textgreater{}} \VariableTok{None}\NormalTok{:}
    \CommentTok{"""Render atoms as spheres (CPK colours, 0.3× vdW radii) and bonds as cylinders."""}
\NormalTok{    ...}


\KeywordTok{class}\NormalTok{ ArcballCamera:}
    \CommentTok{"""Arcball camera: left{-}drag rotates, scroll zooms.}

\CommentTok{    Provides view (lookAt) and perspective (45°, near=0.1, far=200)}
\CommentTok{    matrices for the shader pipeline.}
\CommentTok{    """}

    \KeywordTok{def} \FunctionTok{\_\_init\_\_}\NormalTok{(}
        \VariableTok{self}\NormalTok{,}
\NormalTok{        position: np.ndarray }\OperatorTok{|} \VariableTok{None} \OperatorTok{=} \VariableTok{None}\NormalTok{,}
\NormalTok{        target: np.ndarray }\OperatorTok{|} \VariableTok{None} \OperatorTok{=} \VariableTok{None}\NormalTok{,}
\NormalTok{    ) }\OperatorTok{{-}\textgreater{}} \VariableTok{None}\NormalTok{: ...}

    \KeywordTok{def}\NormalTok{ on\_mouse\_drag(}\VariableTok{self}\NormalTok{, dx: }\BuiltInTok{float}\NormalTok{, dy: }\BuiltInTok{float}\NormalTok{) }\OperatorTok{{-}\textgreater{}} \VariableTok{None}\NormalTok{: ...}

    \KeywordTok{def}\NormalTok{ on\_scroll(}\VariableTok{self}\NormalTok{, delta: }\BuiltInTok{float}\NormalTok{) }\OperatorTok{{-}\textgreater{}} \VariableTok{None}\NormalTok{: ...}

    \KeywordTok{def}\NormalTok{ get\_view\_matrix(}\VariableTok{self}\NormalTok{) }\OperatorTok{{-}\textgreater{}}\NormalTok{ np.ndarray: ...}

    \KeywordTok{def}\NormalTok{ get\_projection\_matrix(}\VariableTok{self}\NormalTok{, aspect: }\BuiltInTok{float}\NormalTok{, fov: }\BuiltInTok{float} \OperatorTok{=} \FloatTok{45.0}\NormalTok{) }\OperatorTok{{-}\textgreater{}}\NormalTok{ np.ndarray: ...}
\end{Highlighting}
\end{Shaded}

\emph{(Full implementation: \texttt{report2\_scripts/renderer.py})}

\subsubsection{\texorpdfstring{10.
\texttt{hardware\_runner.py}}{10. hardware\_runner.py}}\label{hardware_runner.py}

\textbf{Purpose:} Submit VQE circuits to IBM Quantum or IonQ via
PennyLane device plugins.

\begin{itemize}
\tightlist
\item
  \textbf{\texttt{HardwareVQEResult}} --- Data class storing the
  measured energy, raw bitstring counts, shot count, hardware backend
  name, and the number of commuting Pauli groups measured.
\end{itemize}

{\def\LTcaptype{none} % do not increment counter
\begin{longtable}[]{@{}
  >{\raggedright\arraybackslash}p{(\linewidth - 2\tabcolsep) * \real{0.4348}}
  >{\raggedright\arraybackslash}p{(\linewidth - 2\tabcolsep) * \real{0.5652}}@{}}
\toprule\noalign{}
\begin{minipage}[b]{\linewidth}\raggedright
Function
\end{minipage} & \begin{minipage}[b]{\linewidth}\raggedright
Description
\end{minipage} \\
\midrule\noalign{}
\endhead
\bottomrule\noalign{}
\endlastfoot
\texttt{create\_ibm\_device(n\_qubits,\ backend\_name,\ shots,\ ibm\_token)}
& Creates a PennyLane device backed by the PennyLane-Qiskit plugin.
Authenticates with an IBM Quantum API token (from argument or
\texttt{IBMQ\_TOKEN} environment variable). Default backend:
\texttt{ibm\_brisbane}, default shots: 8192. \\
\texttt{create\_ionq\_device(n\_qubits,\ backend,\ shots,\ api\_key)} &
Creates a PennyLane device backed by the PennyLane-IonQ plugin. Default
backend: \texttt{ionq\_harmony} (simulator); switch to \texttt{ionq.qpu}
for real hardware. Authenticates via argument or \texttt{IONQ\_API\_KEY}
environment variable. \\
\texttt{estimate\_shot\_budget(hamiltonian,\ target\_precision)} &
Estimates the total number of shots \(S\) required to achieve standard
deviation \(\leq \epsilon\) on the energy estimate, using the variance
bound \(\text{Var}[\hat{E}] \leq \sum_k |c_k|^2 / S\) (assuming
independent Pauli measurements). This gives
\(S \geq \sum_k |c_k|^2 / \epsilon^2\). For H₂ with
\(\sum|c_k|^2 \approx 0.3\) and \(\epsilon = 1.6 \times 10^{-3}\) Ha,
the bound yields \(S \approx 117{,}000\). The result is clamped to
\([1000, 10^7]\). \\
\end{longtable}
}

\begin{Shaded}
\begin{Highlighting}[]
\ImportTok{from}\NormalTok{ \_\_future\_\_ }\ImportTok{import}\NormalTok{ annotations}

\ImportTok{import}\NormalTok{ os}
\ImportTok{from}\NormalTok{ dataclasses }\ImportTok{import}\NormalTok{ dataclass}

\ImportTok{import}\NormalTok{ numpy }\ImportTok{as}\NormalTok{ np}
\ImportTok{import}\NormalTok{ pennylane }\ImportTok{as}\NormalTok{ qml}


\AttributeTok{@dataclass}
\KeywordTok{class}\NormalTok{ HardwareVQEResult:}
    \CommentTok{"""Results from a hardware VQE execution."""}
\NormalTok{    energy: }\BuiltInTok{float}
\NormalTok{    raw\_counts: }\BuiltInTok{dict}\NormalTok{[}\BuiltInTok{str}\NormalTok{, }\BuiltInTok{int}\NormalTok{]}
\NormalTok{    shots: }\BuiltInTok{int}
\NormalTok{    backend\_name: }\BuiltInTok{str}
\NormalTok{    n\_pauli\_groups: }\BuiltInTok{int}


\KeywordTok{def}\NormalTok{ create\_ibm\_device(}
\NormalTok{    n\_qubits: }\BuiltInTok{int}\NormalTok{,}
\NormalTok{    backend\_name: }\BuiltInTok{str} \OperatorTok{=} \StringTok{"ibm\_brisbane"}\NormalTok{,}
\NormalTok{    shots: }\BuiltInTok{int} \OperatorTok{=} \DecValTok{8192}\NormalTok{,}
\NormalTok{    ibm\_token: }\BuiltInTok{str} \OperatorTok{|} \VariableTok{None} \OperatorTok{=} \VariableTok{None}\NormalTok{,}
\NormalTok{) }\OperatorTok{{-}\textgreater{}}\NormalTok{ qml.Device:}
    \CommentTok{"""Create PennyLane device via qiskit.ibmq plugin.}

\CommentTok{    Authenticates with ibm\_token or IBMQ\_TOKEN env var.}
\CommentTok{    """}
\NormalTok{    ...}


\KeywordTok{def}\NormalTok{ create\_ionq\_device(}
\NormalTok{    n\_qubits: }\BuiltInTok{int}\NormalTok{,}
\NormalTok{    backend: }\BuiltInTok{str} \OperatorTok{=} \StringTok{"ionq\_harmony"}\NormalTok{,}
\NormalTok{    shots: }\BuiltInTok{int} \OperatorTok{=} \DecValTok{1024}\NormalTok{,}
\NormalTok{    api\_key: }\BuiltInTok{str} \OperatorTok{|} \VariableTok{None} \OperatorTok{=} \VariableTok{None}\NormalTok{,}
\NormalTok{) }\OperatorTok{{-}\textgreater{}}\NormalTok{ qml.Device:}
    \CommentTok{"""Create PennyLane device via ionq.simulator plugin.}

\CommentTok{    Authenticates with api\_key or IONQ\_API\_KEY env var.}
\CommentTok{    """}
\NormalTok{    ...}


\KeywordTok{def}\NormalTok{ estimate\_shot\_budget(}
\NormalTok{    hamiltonian: qml.Hamiltonian,}
\NormalTok{    target\_precision: }\BuiltInTok{float} \OperatorTok{=} \FloatTok{1.6e{-}3}\NormalTok{,}
\NormalTok{) }\OperatorTok{{-}\textgreater{}} \BuiltInTok{int}\NormalTok{:}
    \CommentTok{"""Estimate shots S ≥ Σ|c\_k|² / ε² for energy precision ε.}

\CommentTok{    Result clamped to [1000, 10⁷].}
\CommentTok{    """}
\NormalTok{    ...}
\end{Highlighting}
\end{Shaded}

\emph{(Full implementation:
\texttt{report2\_scripts/hardware\_runner.py})}

\subsubsection{\texorpdfstring{11.
\texttt{benchmarker.py}}{11. benchmarker.py}}\label{benchmarker.py}

\textbf{Purpose:} Compare VQE results against classical methods across
the PES.

\begin{itemize}
\tightlist
\item
  \textbf{\texttt{BenchmarkResult}} --- Data class holding scanned bond
  lengths, list of method names, a dictionary of energy arrays keyed by
  method, correlation energy recovery ratios
  \(r_{\text{corr}} = (E_{\text{method}} - E_{\text{HF}}) / (E_{\text{FCI}} - E_{\text{HF}})\),
  and absolute errors \(|E_{\text{method}} - E_{\text{FCI}}|\) for each
  method at every bond length.
\end{itemize}

\begin{quote}
\textbf{\texttt{run\_benchmark(pes\_data)}} takes a \texttt{PESData}
object (from \texttt{scan\_pes}) and computes, for every available
method, the correlation energy recovery ratio and the absolute error
relative to FCI. Division-by-zero is guarded at geometries where
\(E_{\text{FCI}} \approx E_{\text{HF}}\).
\end{quote}

\begin{quote}
\textbf{\texttt{print\_benchmark\_table(result,\ bond\_lengths\_subset=None)}}
formats the benchmark into a human-readable table showing \(R\), each
method's energy, and the VQE-vs-FCI error. If no subset is specified,
five evenly spaced bond lengths are selected automatically.
\end{quote}

\begin{Shaded}
\begin{Highlighting}[]
\ImportTok{from}\NormalTok{ \_\_future\_\_ }\ImportTok{import}\NormalTok{ annotations}

\ImportTok{from}\NormalTok{ dataclasses }\ImportTok{import}\NormalTok{ dataclass}

\ImportTok{import}\NormalTok{ numpy }\ImportTok{as}\NormalTok{ np}


\AttributeTok{@dataclass}
\KeywordTok{class}\NormalTok{ BenchmarkResult:}
    \CommentTok{"""Results from a multi{-}method PES benchmark."""}
\NormalTok{    bond\_lengths: np.ndarray}
\NormalTok{    methods: }\BuiltInTok{list}\NormalTok{[}\BuiltInTok{str}\NormalTok{]}
\NormalTok{    energies: }\BuiltInTok{dict}\NormalTok{[}\BuiltInTok{str}\NormalTok{, np.ndarray]}
\NormalTok{    correlation\_recovery: }\BuiltInTok{dict}\NormalTok{[}\BuiltInTok{str}\NormalTok{, np.ndarray]}
\NormalTok{    errors\_vs\_fci: }\BuiltInTok{dict}\NormalTok{[}\BuiltInTok{str}\NormalTok{, np.ndarray]}


\KeywordTok{def}\NormalTok{ run\_benchmark(pes\_data: PESData) }\OperatorTok{{-}\textgreater{}}\NormalTok{ BenchmarkResult:}
    \CommentTok{"""Compute correlation recovery ratio and |E {-} E\_FCI| for each method.}

\CommentTok{    Division{-}by{-}zero guarded where |E\_FCI {-} E\_HF| \textless{} 1e{-}10.}
\CommentTok{    """}
\NormalTok{    ...}


\KeywordTok{def}\NormalTok{ print\_benchmark\_table(}
\NormalTok{    result: BenchmarkResult,}
\NormalTok{    bond\_lengths\_subset: np.ndarray }\OperatorTok{|} \VariableTok{None} \OperatorTok{=} \VariableTok{None}\NormalTok{,}
\NormalTok{) }\OperatorTok{{-}\textgreater{}} \BuiltInTok{str}\NormalTok{:}
    \CommentTok{"""Format benchmark as a human{-}readable table.}

\CommentTok{    Selects 5 evenly spaced bond lengths if no subset given.}
\CommentTok{    """}
\NormalTok{    ...}
\end{Highlighting}
\end{Shaded}

\emph{(Full implementation: \texttt{report2\_scripts/benchmarker.py})}

\subsubsection{\texorpdfstring{12.
\texttt{main.py}}{12. main.py}}\label{main.py}

\textbf{Purpose:} CLI entry point and workflow orchestration.

\begin{quote}
\textbf{\texttt{parse\_args()}} defines the CLI interface via
\texttt{argparse}. Supported flags include \texttt{-\/-molecule}
(\texttt{h2}, \texttt{lih}, \texttt{h2o}), \texttt{-\/-basis} (e.g.,
\texttt{sto-3g}, \texttt{6-31g*}), \texttt{-\/-bond-range}
(\texttt{start,stop,n\_points}), \texttt{-\/-ansatz} (\texttt{uccsd} or
\texttt{hardware\_efficient}), \texttt{-\/-optimizer} (\texttt{Adam},
\texttt{GradientDescent}, \texttt{COBYLA}, \texttt{SPSA}),
\texttt{-\/-maxiter}, \texttt{-\/-hardware} / \texttt{-\/-backend} for
real-device execution, \texttt{-\/-render} to enable the OpenGL viewer,
\texttt{-\/-benchmark} to compute all classical reference methods,
\texttt{-\/-orbital} for the MO index to visualise, and
\texttt{-\/-config} for a YAML configuration file.
\end{quote}

\begin{quote}
\textbf{\texttt{run\_pipeline(args)}} orchestrates the full workflow:

\begin{enumerate}
\def\labelenumi{\arabic{enumi}.}
\tightlist
\item
  Parses the bond-range string into a \texttt{linspace} array.
\item
  Selects the appropriate molecule factory function.
\item
  Calls \texttt{scan\_pes} to sweep bond lengths with VQE.
\item
  If \texttt{-\/-benchmark} is set, calls \texttt{run\_benchmark} and
  \texttt{print\_benchmark\_table}.
\item
  If \texttt{-\/-render} is set, runs HF at equilibrium, evaluates the
  selected MO on a 3D grid, extracts the dual isosurface, and launches
  the OpenGL render loop with arcball camera interaction.
\item
  Prints a summary: equilibrium VQE energy, bond length, FCI energy, and
  absolute error.
\end{enumerate}
\end{quote}

\begin{Shaded}
\begin{Highlighting}[]
\ImportTok{from}\NormalTok{ \_\_future\_\_ }\ImportTok{import}\NormalTok{ annotations}

\ImportTok{import}\NormalTok{ argparse}

\ImportTok{import}\NormalTok{ numpy }\ImportTok{as}\NormalTok{ np}


\KeywordTok{def}\NormalTok{ parse\_args() }\OperatorTok{{-}\textgreater{}}\NormalTok{ argparse.Namespace:}
    \CommentTok{"""Define CLI: {-}{-}molecule, {-}{-}basis, {-}{-}bond{-}range, {-}{-}ansatz, {-}{-}optimizer,}
\CommentTok{    {-}{-}maxiter, {-}{-}hardware, {-}{-}backend, {-}{-}render, {-}{-}benchmark, {-}{-}orbital, {-}{-}config.}
\CommentTok{    """}
\NormalTok{    ...}


\KeywordTok{def}\NormalTok{ run\_pipeline(args: argparse.Namespace) }\OperatorTok{{-}\textgreater{}} \VariableTok{None}\NormalTok{:}
    \CommentTok{"""Orchestrate: parse bond range → select factory → scan PES →}
\CommentTok{    benchmark (optional) → orbital render (optional) → print summary.}
\CommentTok{    """}
\NormalTok{    ...}
\end{Highlighting}
\end{Shaded}

\emph{(Full implementation: \texttt{report2\_scripts/main.py})}

\begin{center}\rule{0.5\linewidth}{0.5pt}\end{center}

\section{Part 3 --- Implementation Guidelines (Instructional
Pseudocode)}\label{part-3-implementation-guidelines-instructional-pseudocode}

\subsection{\texorpdfstring{Module 1:
\texttt{molecule.py}}{Module 1: molecule.py}}\label{module-1-molecule.py}

\begin{Shaded}
\begin{Highlighting}[]
\CommentTok{\# {-}{-}{-} molecule.py pseudocode {-}{-}{-}}

\KeywordTok{def}\NormalTok{ create\_h2(bond\_length}\OperatorTok{=}\FloatTok{0.735}\NormalTok{):}
    \CommentTok{\# Place two H on z{-}axis at ±R/2}
\NormalTok{    atoms }\OperatorTok{=}\NormalTok{ [}
\NormalTok{        Atom(}\StringTok{"H"}\NormalTok{, np.array([}\FloatTok{0.0}\NormalTok{, }\FloatTok{0.0}\NormalTok{, }\OperatorTok{{-}}\NormalTok{bond\_length }\OperatorTok{/} \DecValTok{2}\NormalTok{])),}
\NormalTok{        Atom(}\StringTok{"H"}\NormalTok{, np.array([}\FloatTok{0.0}\NormalTok{, }\FloatTok{0.0}\NormalTok{, }\OperatorTok{+}\NormalTok{bond\_length }\OperatorTok{/} \DecValTok{2}\NormalTok{])),}
\NormalTok{    ]}
    \ControlFlowTok{return}\NormalTok{ MoleculeSpec(atoms, basis}\OperatorTok{=}\StringTok{"sto{-}3g"}\NormalTok{, charge}\OperatorTok{=}\DecValTok{0}\NormalTok{,}
\NormalTok{                        multiplicity}\OperatorTok{=}\DecValTok{1}\NormalTok{, name}\OperatorTok{=}\SpecialStringTok{f"H2\_R}\SpecialCharTok{\{}\NormalTok{bond\_length}\SpecialCharTok{:.3f\}}\SpecialStringTok{"}\NormalTok{)}


\KeywordTok{def}\NormalTok{ create\_lih(bond\_length}\OperatorTok{=}\FloatTok{1.595}\NormalTok{):}
    \CommentTok{\# Li at origin, H at (0,0,R)}
\NormalTok{    atoms }\OperatorTok{=}\NormalTok{ [}
\NormalTok{        Atom(}\StringTok{"Li"}\NormalTok{, np.array([}\FloatTok{0.0}\NormalTok{, }\FloatTok{0.0}\NormalTok{, }\FloatTok{0.0}\NormalTok{])),}
\NormalTok{        Atom(}\StringTok{"H"}\NormalTok{,  np.array([}\FloatTok{0.0}\NormalTok{, }\FloatTok{0.0}\NormalTok{, bond\_length])),}
\NormalTok{    ]}
    \ControlFlowTok{return}\NormalTok{ MoleculeSpec(atoms, basis}\OperatorTok{=}\StringTok{"sto{-}3g"}\NormalTok{, name}\OperatorTok{=}\SpecialStringTok{f"LiH\_R}\SpecialCharTok{\{}\NormalTok{bond\_length}\SpecialCharTok{:.3f\}}\SpecialStringTok{"}\NormalTok{)}


\KeywordTok{def}\NormalTok{ create\_h2o(oh\_length}\OperatorTok{=}\FloatTok{0.957}\NormalTok{, angle\_deg}\OperatorTok{=}\FloatTok{104.5}\NormalTok{):}
\NormalTok{    alpha }\OperatorTok{=}\NormalTok{ np.radians(angle\_deg)}
    \CommentTok{\# H positions: (±R*sin(α/2), 0, {-}R*cos(α/2))}
\NormalTok{    h1 }\OperatorTok{=}\NormalTok{ np.array([}\OperatorTok{+}\NormalTok{oh\_length }\OperatorTok{*}\NormalTok{ np.sin(alpha }\OperatorTok{/} \DecValTok{2}\NormalTok{), }\FloatTok{0.0}\NormalTok{,}
                    \OperatorTok{{-}}\NormalTok{oh\_length }\OperatorTok{*}\NormalTok{ np.cos(alpha }\OperatorTok{/} \DecValTok{2}\NormalTok{)])}
\NormalTok{    h2 }\OperatorTok{=}\NormalTok{ np.array([}\OperatorTok{{-}}\NormalTok{oh\_length }\OperatorTok{*}\NormalTok{ np.sin(alpha }\OperatorTok{/} \DecValTok{2}\NormalTok{), }\FloatTok{0.0}\NormalTok{,}
                    \OperatorTok{{-}}\NormalTok{oh\_length }\OperatorTok{*}\NormalTok{ np.cos(alpha }\OperatorTok{/} \DecValTok{2}\NormalTok{)])}
\NormalTok{    atoms }\OperatorTok{=}\NormalTok{ [Atom(}\StringTok{"O"}\NormalTok{, np.zeros(}\DecValTok{3}\NormalTok{)), Atom(}\StringTok{"H"}\NormalTok{, h1), Atom(}\StringTok{"H"}\NormalTok{, h2)]}
    \ControlFlowTok{return}\NormalTok{ MoleculeSpec(atoms, basis}\OperatorTok{=}\StringTok{"sto{-}3g"}\NormalTok{, name}\OperatorTok{=}\SpecialStringTok{f"H2O\_R}\SpecialCharTok{\{}\NormalTok{oh\_length}\SpecialCharTok{:.3f\}}\SpecialStringTok{"}\NormalTok{)}


\KeywordTok{def}\NormalTok{ set\_bond\_length(mol, atom\_idx\_a, atom\_idx\_b, new\_length):}
\NormalTok{    new\_mol }\OperatorTok{=}\NormalTok{ copy.deepcopy(mol)}
\NormalTok{    pa }\OperatorTok{=}\NormalTok{ new\_mol.atoms[atom\_idx\_a].position  }\CommentTok{\# shape: (3,)}
\NormalTok{    pb }\OperatorTok{=}\NormalTok{ new\_mol.atoms[atom\_idx\_b].position  }\CommentTok{\# shape: (3,)}
\NormalTok{    midpoint }\OperatorTok{=}\NormalTok{ (pa }\OperatorTok{+}\NormalTok{ pb) }\OperatorTok{/} \FloatTok{2.0}               \CommentTok{\# shape: (3,)}
\NormalTok{    d\_hat }\OperatorTok{=}\NormalTok{ (pb }\OperatorTok{{-}}\NormalTok{ pa)                        }\CommentTok{\# direction vector}
\NormalTok{    d\_hat }\OperatorTok{=}\NormalTok{ d\_hat }\OperatorTok{/}\NormalTok{ np.linalg.norm(d\_hat)    }\CommentTok{\# unit vector}
    \CommentTok{\# Reposition symmetrically about midpoint}
\NormalTok{    new\_mol.atoms[atom\_idx\_a].position }\OperatorTok{=}\NormalTok{ midpoint }\OperatorTok{{-}}\NormalTok{ d\_hat }\OperatorTok{*}\NormalTok{ new\_length }\OperatorTok{/} \DecValTok{2}
\NormalTok{    new\_mol.atoms[atom\_idx\_b].position }\OperatorTok{=}\NormalTok{ midpoint }\OperatorTok{+}\NormalTok{ d\_hat }\OperatorTok{*}\NormalTok{ new\_length }\OperatorTok{/} \DecValTok{2}
    \ControlFlowTok{return}\NormalTok{ new\_mol}
\end{Highlighting}
\end{Shaded}

\emph{(Full implementation: \texttt{report2\_scripts/molecule.py})}

\subsection{\texorpdfstring{Module 2:
\texttt{classical\_hf.py}}{Module 2: classical\_hf.py}}\label{module-2-classical_hf.py}

\begin{Shaded}
\begin{Highlighting}[]
\CommentTok{\# {-}{-}{-} classical\_hf.py pseudocode {-}{-}{-}}

\KeywordTok{def}\NormalTok{ \_build\_pyscf\_mol(mol\_spec):}
    \CommentTok{"""Shared helper: MoleculeSpec → PySCF Mole."""}
    \CommentTok{\# Format atom string: "H 0.0 0.0 0.0; H 0.0 0.0 0.735"}
\NormalTok{    atom\_str }\OperatorTok{=} \StringTok{"; "}\NormalTok{.join(}
        \SpecialStringTok{f"}\SpecialCharTok{\{}\NormalTok{a}\SpecialCharTok{.}\NormalTok{symbol}\SpecialCharTok{\}}\SpecialStringTok{ }\SpecialCharTok{\{}\NormalTok{a}\SpecialCharTok{.}\NormalTok{position[}\DecValTok{0}\NormalTok{]}\SpecialCharTok{:.6f\}}\SpecialStringTok{ }\SpecialCharTok{\{}\NormalTok{a}\SpecialCharTok{.}\NormalTok{position[}\DecValTok{1}\NormalTok{]}\SpecialCharTok{:.6f\}}\SpecialStringTok{ }\SpecialCharTok{\{}\NormalTok{a}\SpecialCharTok{.}\NormalTok{position[}\DecValTok{2}\NormalTok{]}\SpecialCharTok{:.6f\}}\SpecialStringTok{"}
        \ControlFlowTok{for}\NormalTok{ a }\KeywordTok{in}\NormalTok{ mol\_spec.atoms}
\NormalTok{    )}
\NormalTok{    mol }\OperatorTok{=}\NormalTok{ gto.Mole()}
\NormalTok{    mol.atom }\OperatorTok{=}\NormalTok{ atom\_str}
\NormalTok{    mol.basis }\OperatorTok{=}\NormalTok{ mol\_spec.basis}
\NormalTok{    mol.charge }\OperatorTok{=}\NormalTok{ mol\_spec.charge}
\NormalTok{    mol.spin }\OperatorTok{=}\NormalTok{ mol\_spec.multiplicity }\OperatorTok{{-}} \DecValTok{1}  \CommentTok{\# PySCF uses 2S, not 2S+1}
\NormalTok{    mol.unit }\OperatorTok{=} \StringTok{"Angstrom"}
\NormalTok{    mol.build()}
    \ControlFlowTok{return}\NormalTok{ mol}


\KeywordTok{def}\NormalTok{ run\_hartree\_fock(mol\_spec):}
\NormalTok{    mol }\OperatorTok{=}\NormalTok{ \_build\_pyscf\_mol(mol\_spec)}

    \CommentTok{\# Run restricted Hartree{-}Fock}
\NormalTok{    mf }\OperatorTok{=}\NormalTok{ scf.RHF(mol)}
\NormalTok{    mf.kernel()  }\CommentTok{\# converge SCF}

\NormalTok{    mo\_coeff }\OperatorTok{=}\NormalTok{ mf.mo\_coeff      }\CommentTok{\# shape: (n\_ao, n\_mo) — AO→MO matrix}
\NormalTok{    mo\_energy }\OperatorTok{=}\NormalTok{ mf.mo\_energy    }\CommentTok{\# shape: (n\_mo,)}
\NormalTok{    n\_mo }\OperatorTok{=}\NormalTok{ mo\_coeff.shape[}\DecValTok{1}\NormalTok{]}

    \CommentTok{\# One{-}electron integrals: AO → MO via C\^{}T h\^{}core C}
\NormalTok{    h1\_ao }\OperatorTok{=}\NormalTok{ mf.get\_hcore()                  }\CommentTok{\# shape: (n\_ao, n\_ao)}
\NormalTok{    h1\_mo }\OperatorTok{=}\NormalTok{ mo\_coeff.T }\OperatorTok{@}\NormalTok{ h1\_ao }\OperatorTok{@}\NormalTok{ mo\_coeff   }\CommentTok{\# shape: (n\_mo, n\_mo)}

    \CommentTok{\# Two{-}electron integrals: efficient AO→MO with ao2mo.full}
\NormalTok{    h2\_mo }\OperatorTok{=}\NormalTok{ ao2mo.full(mol, mo\_coeff)       }\CommentTok{\# packed triangular form}
\NormalTok{    h2\_mo }\OperatorTok{=}\NormalTok{ ao2mo.restore(}\DecValTok{1}\NormalTok{, h2\_mo, n\_mo)   }\CommentTok{\# shape: (n\_mo, n\_mo, n\_mo, n\_mo)}

\NormalTok{    overlap }\OperatorTok{=}\NormalTok{ mol.intor(}\StringTok{"int1e\_ovlp"}\NormalTok{)       }\CommentTok{\# shape: (n\_ao, n\_ao)}
\NormalTok{    e\_nuc }\OperatorTok{=}\NormalTok{ mol.energy\_nuc()                }\CommentTok{\# scalar (Hartree)}

    \ControlFlowTok{return}\NormalTok{ HFResult(}
\NormalTok{        energy\_hf}\OperatorTok{=}\NormalTok{mf.e\_tot, energy\_nuc}\OperatorTok{=}\NormalTok{e\_nuc,}
\NormalTok{        mo\_coefficients}\OperatorTok{=}\NormalTok{mo\_coeff, mo\_energies}\OperatorTok{=}\NormalTok{mo\_energy,}
\NormalTok{        one\_electron\_integrals}\OperatorTok{=}\NormalTok{h1\_mo, two\_electron\_integrals}\OperatorTok{=}\NormalTok{h2\_mo,}
\NormalTok{        overlap\_matrix}\OperatorTok{=}\NormalTok{overlap,}
\NormalTok{        n\_electrons}\OperatorTok{=}\NormalTok{mol.nelectron, n\_orbitals}\OperatorTok{=}\NormalTok{n\_mo,}
\NormalTok{        ao\_labels}\OperatorTok{=}\NormalTok{[}\BuiltInTok{str}\NormalTok{(l) }\ControlFlowTok{for}\NormalTok{ l }\KeywordTok{in}\NormalTok{ mol.ao\_labels()],}
\NormalTok{    )}


\KeywordTok{def}\NormalTok{ get\_fci\_energy(mol\_spec):}
\NormalTok{    mol }\OperatorTok{=}\NormalTok{ \_build\_pyscf\_mol(mol\_spec)}
\NormalTok{    mf }\OperatorTok{=}\NormalTok{ scf.RHF(mol).run()       }\CommentTok{\# RHF as starting point}
\NormalTok{    cisolver }\OperatorTok{=}\NormalTok{ fci.FCI(mf)}
\NormalTok{    e\_fci, \_ }\OperatorTok{=}\NormalTok{ cisolver.kernel()   }\CommentTok{\# exact ground{-}state energy}
    \ControlFlowTok{return} \BuiltInTok{float}\NormalTok{(e\_fci)}


\KeywordTok{def}\NormalTok{ get\_cisd\_energy(mol\_spec):}
\NormalTok{    mol }\OperatorTok{=}\NormalTok{ \_build\_pyscf\_mol(mol\_spec)}
\NormalTok{    mf }\OperatorTok{=}\NormalTok{ scf.RHF(mol).run()}
\NormalTok{    myci }\OperatorTok{=}\NormalTok{ ci.CISD(mf).run()}
    \ControlFlowTok{return} \BuiltInTok{float}\NormalTok{(myci.e\_tot)       }\CommentTok{\# E\_CISD (Hartree)}


\KeywordTok{def}\NormalTok{ get\_ccsd\_t\_energy(mol\_spec):}
\NormalTok{    mol }\OperatorTok{=}\NormalTok{ \_build\_pyscf\_mol(mol\_spec)}
\NormalTok{    mf }\OperatorTok{=}\NormalTok{ scf.RHF(mol).run()}
\NormalTok{    mycc }\OperatorTok{=}\NormalTok{ cc.CCSD(mf).run()}
    \CommentTok{\# Perturbative triples correction}
\NormalTok{    e\_t }\OperatorTok{=}\NormalTok{ mycc.ccsd\_t()            }\CommentTok{\# ΔE\_(T)}
    \ControlFlowTok{return} \BuiltInTok{float}\NormalTok{(mycc.e\_tot }\OperatorTok{+}\NormalTok{ e\_t)  }\CommentTok{\# E\_CCSD(T) = E\_CCSD + ΔE\_(T)}
\end{Highlighting}
\end{Shaded}

\emph{(Full implementation: \texttt{report2\_scripts/classical\_hf.py})}

\subsection{\texorpdfstring{Module 3:
\texttt{hamiltonian\_builder.py}}{Module 3: hamiltonian\_builder.py}}\label{module-3-hamiltonian_builder.py}

\begin{Shaded}
\begin{Highlighting}[]
\CommentTok{\# {-}{-}{-} hamiltonian\_builder.py pseudocode {-}{-}{-}}

\KeywordTok{def}\NormalTok{ build\_hamiltonian(h1, h2, n\_electrons, e\_nuc,}
\NormalTok{                      mapping}\OperatorTok{=}\StringTok{"jordan\_wigner"}\NormalTok{,}
\NormalTok{                      active\_electrons}\OperatorTok{=}\VariableTok{None}\NormalTok{, active\_orbitals}\OperatorTok{=}\VariableTok{None}\NormalTok{):}
\NormalTok{    n\_mo }\OperatorTok{=}\NormalTok{ h1.shape[}\DecValTok{0}\NormalTok{]}
    \CommentTok{\# n\_spin\_orbs = 2 * n\_mo  (each spatial orbital → 2 spin{-}orbitals)}

    \CommentTok{\# {-}{-}{-} Stage 1: Build fermionic Hamiltonian {-}{-}{-}}
    \CommentTok{\# One{-}electron terms: h\_ij a+\_\{2i+σ\} a\_\{2j+σ\}  (spin{-}diagonal)}
\NormalTok{    fermi\_h }\OperatorTok{=} \DecValTok{0}
    \ControlFlowTok{for}\NormalTok{ i }\KeywordTok{in} \BuiltInTok{range}\NormalTok{(n\_mo):}
        \ControlFlowTok{for}\NormalTok{ j }\KeywordTok{in} \BuiltInTok{range}\NormalTok{(n\_mo):}
            \ControlFlowTok{if} \BuiltInTok{abs}\NormalTok{(h1[i, j]) }\OperatorTok{\textless{}} \FloatTok{1e{-}12}\NormalTok{:}
                \ControlFlowTok{continue}
            \ControlFlowTok{for}\NormalTok{ sigma }\KeywordTok{in}\NormalTok{ (}\DecValTok{0}\NormalTok{, }\DecValTok{1}\NormalTok{):  }\CommentTok{\# alpha, beta}
\NormalTok{                p, q }\OperatorTok{=} \DecValTok{2} \OperatorTok{*}\NormalTok{ i }\OperatorTok{+}\NormalTok{ sigma, }\DecValTok{2} \OperatorTok{*}\NormalTok{ j }\OperatorTok{+}\NormalTok{ sigma}
\NormalTok{                fermi\_h }\OperatorTok{+=}\NormalTok{ h1[i, j] }\OperatorTok{*}\NormalTok{ FermiWord(\{(}\DecValTok{0}\NormalTok{, p): }\StringTok{"+"}\NormalTok{, (}\DecValTok{1}\NormalTok{, q): }\StringTok{"{-}"}\NormalTok{\})}

    \CommentTok{\# Two{-}electron terms: ½ h\_ijkl a+\_p a+\_q a\_s a\_r (all spin combos, skip p==q)}
    \ControlFlowTok{for}\NormalTok{ i }\KeywordTok{in} \BuiltInTok{range}\NormalTok{(n\_mo):}
        \ControlFlowTok{for}\NormalTok{ j }\KeywordTok{in} \BuiltInTok{range}\NormalTok{(n\_mo):}
            \ControlFlowTok{for}\NormalTok{ k }\KeywordTok{in} \BuiltInTok{range}\NormalTok{(n\_mo):}
                \ControlFlowTok{for}\NormalTok{ l }\KeywordTok{in} \BuiltInTok{range}\NormalTok{(n\_mo):}
                    \ControlFlowTok{if} \BuiltInTok{abs}\NormalTok{(h2[i, j, k, l]) }\OperatorTok{\textless{}} \FloatTok{1e{-}12}\NormalTok{:}
                        \ControlFlowTok{continue}
                    \ControlFlowTok{for}\NormalTok{ s1 }\KeywordTok{in}\NormalTok{ (}\DecValTok{0}\NormalTok{, }\DecValTok{1}\NormalTok{):}
                        \ControlFlowTok{for}\NormalTok{ s2 }\KeywordTok{in}\NormalTok{ (}\DecValTok{0}\NormalTok{, }\DecValTok{1}\NormalTok{):}
\NormalTok{                            p }\OperatorTok{=} \DecValTok{2} \OperatorTok{*}\NormalTok{ i }\OperatorTok{+}\NormalTok{ s1}
\NormalTok{                            q }\OperatorTok{=} \DecValTok{2} \OperatorTok{*}\NormalTok{ k }\OperatorTok{+}\NormalTok{ s2}
\NormalTok{                            r }\OperatorTok{=} \DecValTok{2} \OperatorTok{*}\NormalTok{ j }\OperatorTok{+}\NormalTok{ s1}
\NormalTok{                            s }\OperatorTok{=} \DecValTok{2} \OperatorTok{*}\NormalTok{ l }\OperatorTok{+}\NormalTok{ s2}
                            \ControlFlowTok{if}\NormalTok{ p }\OperatorTok{==}\NormalTok{ q:}
                                \ControlFlowTok{continue}
\NormalTok{                            fermi\_h }\OperatorTok{+=} \FloatTok{0.5} \OperatorTok{*}\NormalTok{ h2[i, j, k, l] }\OperatorTok{*}\NormalTok{ FermiWord(}
\NormalTok{                                \{(}\DecValTok{0}\NormalTok{, p): }\StringTok{"+"}\NormalTok{, (}\DecValTok{1}\NormalTok{, q): }\StringTok{"+"}\NormalTok{, (}\DecValTok{2}\NormalTok{, s): }\StringTok{"{-}"}\NormalTok{, (}\DecValTok{3}\NormalTok{, r): }\StringTok{"{-}"}\NormalTok{\}}
\NormalTok{                            )}

    \CommentTok{\# Nuclear repulsion: constant * Identity}
    \CommentTok{\# fermi\_h += e\_nuc * Identity}

    \CommentTok{\# {-}{-}{-} Stage 2: Fermion{-}to{-}qubit mapping {-}{-}{-}}
    \ControlFlowTok{if}\NormalTok{ mapping }\OperatorTok{==} \StringTok{"jordan\_wigner"}\NormalTok{:}
\NormalTok{        qubit\_h }\OperatorTok{=}\NormalTok{ qml.jordan\_wigner(fermi\_h)   }\CommentTok{\# JW: a+\_p → ½(X{-}iY)∏Z}
    \ControlFlowTok{else}\NormalTok{:}
\NormalTok{        qubit\_h }\OperatorTok{=}\NormalTok{ qml.parity\_transform(fermi\_h)  }\CommentTok{\# parity encoding}
\NormalTok{    qubit\_h }\OperatorTok{=}\NormalTok{ qml.simplify(qubit\_h)  }\CommentTok{\# merge \& cancel redundant Pauli terms}
    \CommentTok{\# ... wrap into QubitHamiltonian dataclass ...}

    \CommentTok{\# Alternative (simpler): one{-}call PennyLane interface}
    \CommentTok{\# H, n\_qubits = qml.qchem.molecular\_hamiltonian(}
    \CommentTok{\#     symbols, coords, basis=basis, mapping=mapping)}


\KeywordTok{def}\NormalTok{ pennylane\_to\_qutip(hamiltonian, n\_qubits):}
    \CommentTok{"""Cross{-}validate by converting to QuTiP and diagonalising."""}
\NormalTok{    pauli\_map }\OperatorTok{=}\NormalTok{ \{}\StringTok{"I"}\NormalTok{: qutip.qeye(}\DecValTok{2}\NormalTok{), }\StringTok{"X"}\NormalTok{: qutip.sigmax(),}
                 \StringTok{"Y"}\NormalTok{: qutip.sigmay(), }\StringTok{"Z"}\NormalTok{: qutip.sigmaz()\}}
\NormalTok{    H\_qt }\OperatorTok{=} \DecValTok{0}
    \ControlFlowTok{for}\NormalTok{ coeff, obs }\KeywordTok{in} \BuiltInTok{zip}\NormalTok{(hamiltonian.coeffs, hamiltonian.ops):}
        \CommentTok{\# Decompose obs into per{-}qubit Pauli labels}
\NormalTok{        factors }\OperatorTok{=}\NormalTok{ [pauli\_map[}\StringTok{"I"}\NormalTok{]] }\OperatorTok{*}\NormalTok{ n\_qubits}
        \ControlFlowTok{for}\NormalTok{ wire, pauli\_label }\KeywordTok{in}\NormalTok{ \_decompose\_pauli(obs):}
\NormalTok{            factors[wire] }\OperatorTok{=}\NormalTok{ pauli\_map[pauli\_label]}
        \CommentTok{\# Build tensor product: ⊗\_q σ\_q}
\NormalTok{        term }\OperatorTok{=}\NormalTok{ qutip.tensor(factors)        }\CommentTok{\# dim: (2\^{}N, 2\^{}N)}
\NormalTok{        H\_qt }\OperatorTok{=}\NormalTok{ H\_qt }\OperatorTok{+}\NormalTok{ coeff }\OperatorTok{*}\NormalTok{ term}
    \CommentTok{\# Exact diagonalisation}
\NormalTok{    eigenvalues }\OperatorTok{=}\NormalTok{ H\_qt.eigenenergies()      }\CommentTok{\# sorted ascending}
\NormalTok{    e0 }\OperatorTok{=} \BuiltInTok{float}\NormalTok{(eigenvalues[}\DecValTok{0}\NormalTok{])              }\CommentTok{\# ground{-}state energy}
    \ControlFlowTok{return}\NormalTok{ e0}
\end{Highlighting}
\end{Shaded}

\emph{(Full implementation:
\texttt{report2\_scripts/hamiltonian\_builder.py})}

\subsection{\texorpdfstring{Module 4:
\texttt{ansatz.py}}{Module 4: ansatz.py}}\label{module-4-ansatz.py}

\begin{Shaded}
\begin{Highlighting}[]
\CommentTok{\# {-}{-}{-} ansatz.py pseudocode {-}{-}{-}}

\KeywordTok{def}\NormalTok{ build\_uccsd\_ansatz(n\_qubits, n\_electrons, hf\_state}\OperatorTok{=}\VariableTok{None}\NormalTok{):}
    \CommentTok{\# Prepare HF reference state: first n\_electrons orbitals occupied}
    \ControlFlowTok{if}\NormalTok{ hf\_state }\KeywordTok{is} \VariableTok{None}\NormalTok{:}
\NormalTok{        hf\_state }\OperatorTok{=}\NormalTok{ np.zeros(n\_qubits, dtype}\OperatorTok{=}\BuiltInTok{int}\NormalTok{)}
\NormalTok{        hf\_state[:n\_electrons] }\OperatorTok{=} \DecValTok{1}   \CommentTok{\# e.g. |1100⟩ for H₂}

    \CommentTok{\# Enumerate spin{-}allowed excitations via PennyLane qchem}
\NormalTok{    singles, doubles }\OperatorTok{=}\NormalTok{ qml.qchem.excitations(n\_electrons, n\_qubits)}
    \CommentTok{\# H₂ example: singles = [[0,2], [1,3]], doubles = [[0,1,2,3]]}
\NormalTok{    excitations }\OperatorTok{=}\NormalTok{ singles }\OperatorTok{+}\NormalTok{ doubles}
\NormalTok{    n\_params }\OperatorTok{=} \BuiltInTok{len}\NormalTok{(excitations)             }\CommentTok{\# = |singles| + |doubles|}

    \KeywordTok{def}\NormalTok{ ansatz\_fn(params):                  }\CommentTok{\# params shape: (n\_params,)}
\NormalTok{        qml.BasisState(hf\_state, wires}\OperatorTok{=}\BuiltInTok{range}\NormalTok{(n\_qubits))}

\NormalTok{        idx }\OperatorTok{=} \DecValTok{0}
        \ControlFlowTok{for}\NormalTok{ exc }\KeywordTok{in}\NormalTok{ singles:}
            \CommentTok{\# exp(iθ (a+\_a a\_i − h.c.)) via first{-}order Trotter}
\NormalTok{            qml.SingleExcitation(params[idx], wires}\OperatorTok{=}\NormalTok{exc)}
\NormalTok{            idx }\OperatorTok{+=} \DecValTok{1}
        \ControlFlowTok{for}\NormalTok{ exc }\KeywordTok{in}\NormalTok{ doubles:}
            \CommentTok{\# exp(iθ (a+\_a a+\_b a\_j a\_i − h.c.))}
\NormalTok{            qml.DoubleExcitation(params[idx], wires}\OperatorTok{=}\NormalTok{exc)}
\NormalTok{            idx }\OperatorTok{+=} \DecValTok{1}

    \CommentTok{\# }\AlertTok{NOTE}\CommentTok{: Excitation ordering determines Trotter decomposition.}
    \CommentTok{\# Optimizer compensates for Trotter error.}
    \CommentTok{\# For H₂: by symmetry, singles are zero at the optimum;}
    \CommentTok{\# only the double excitation matters (1 effective parameter).}
    \ControlFlowTok{return}\NormalTok{ ansatz\_fn, n\_params, excitations}


\KeywordTok{def}\NormalTok{ build\_hardware\_efficient\_ansatz(n\_qubits, n\_layers}\OperatorTok{=}\DecValTok{2}\NormalTok{, entangler}\OperatorTok{=}\StringTok{"CNOT"}\NormalTok{):}
\NormalTok{    n\_params }\OperatorTok{=} \DecValTok{2} \OperatorTok{*}\NormalTok{ n\_qubits }\OperatorTok{*}\NormalTok{ n\_layers      }\CommentTok{\# RY + RZ per qubit per layer}

    \KeywordTok{def}\NormalTok{ ansatz\_fn(params):                   }\CommentTok{\# params shape: (n\_params,)}
\NormalTok{        p }\OperatorTok{=}\NormalTok{ params.reshape(n\_layers, n\_qubits, }\DecValTok{2}\NormalTok{)  }\CommentTok{\# shape: (L, N\_q, 2)}

        \ControlFlowTok{for}\NormalTok{ l }\KeywordTok{in} \BuiltInTok{range}\NormalTok{(n\_layers):}
            \CommentTok{\# Single{-}qubit rotation layer}
            \ControlFlowTok{for}\NormalTok{ q }\KeywordTok{in} \BuiltInTok{range}\NormalTok{(n\_qubits):}
\NormalTok{                qml.RY(p[l, q, }\DecValTok{0}\NormalTok{], wires}\OperatorTok{=}\NormalTok{q)}
\NormalTok{                qml.RZ(p[l, q, }\DecValTok{1}\NormalTok{], wires}\OperatorTok{=}\NormalTok{q)}
            \CommentTok{\# Entangling ladder (nearest{-}neighbour)}
            \ControlFlowTok{for}\NormalTok{ q }\KeywordTok{in} \BuiltInTok{range}\NormalTok{(n\_qubits }\OperatorTok{{-}} \DecValTok{1}\NormalTok{):}
                \ControlFlowTok{if}\NormalTok{ entangler }\OperatorTok{==} \StringTok{"CNOT"}\NormalTok{:}
\NormalTok{                    qml.CNOT(wires}\OperatorTok{=}\NormalTok{[q, q }\OperatorTok{+} \DecValTok{1}\NormalTok{])}
                \ControlFlowTok{elif}\NormalTok{ entangler }\OperatorTok{==} \StringTok{"CZ"}\NormalTok{:}
\NormalTok{                    qml.CZ(wires}\OperatorTok{=}\NormalTok{[q, q }\OperatorTok{+} \DecValTok{1}\NormalTok{])}

    \ControlFlowTok{return}\NormalTok{ ansatz\_fn, n\_params}
\end{Highlighting}
\end{Shaded}

\emph{(Full implementation: \texttt{report2\_scripts/ansatz.py})}

\subsection{\texorpdfstring{Module 5:
\texttt{vqe\_solver.py}}{Module 5: vqe\_solver.py}}\label{module-5-vqe_solver.py}

\begin{Shaded}
\begin{Highlighting}[]
\CommentTok{\# {-}{-}{-} vqe\_solver.py pseudocode {-}{-}{-}}

\KeywordTok{def}\NormalTok{ run\_vqe(hamiltonian, ansatz\_fn, n\_params, n\_qubits,}
\NormalTok{            init\_params}\OperatorTok{=}\VariableTok{None}\NormalTok{, optimizer}\OperatorTok{=}\StringTok{"Adam"}\NormalTok{, learning\_rate}\OperatorTok{=}\FloatTok{0.05}\NormalTok{,}
\NormalTok{            maxiter}\OperatorTok{=}\DecValTok{200}\NormalTok{, tol}\OperatorTok{=}\FloatTok{1e{-}6}\NormalTok{, device\_name}\OperatorTok{=}\StringTok{"default.qubit"}\NormalTok{,}
\NormalTok{            diff\_method}\OperatorTok{=}\StringTok{"backprop"}\NormalTok{, callback}\OperatorTok{=}\VariableTok{None}\NormalTok{):}

\NormalTok{    dev }\OperatorTok{=}\NormalTok{ qml.device(device\_name, wires}\OperatorTok{=}\NormalTok{n\_qubits)}

    \AttributeTok{@qml.qnode}\NormalTok{(dev, diff\_method}\OperatorTok{=}\NormalTok{diff\_method)}
    \KeywordTok{def}\NormalTok{ cost\_fn(params):                    }\CommentTok{\# params shape: (n\_params,)}
\NormalTok{        ansatz\_fn(params)}
        \ControlFlowTok{return}\NormalTok{ qml.expval(hamiltonian)      }\CommentTok{\# \textless{}ψ(θ)|H|ψ(θ)\textgreater{}}

    \CommentTok{\# Initialise with small random perturbation to break symmetry}
    \ControlFlowTok{if}\NormalTok{ init\_params }\KeywordTok{is} \VariableTok{None}\NormalTok{:}
\NormalTok{        init\_params }\OperatorTok{=}\NormalTok{ np.random.normal(}\DecValTok{0}\NormalTok{, }\FloatTok{0.01}\NormalTok{, size}\OperatorTok{=}\NormalTok{n\_params)  }\CommentTok{\# shape: (n\_params,)}
\NormalTok{    params }\OperatorTok{=}\NormalTok{ init\_params.copy()}
\NormalTok{    trajectory }\OperatorTok{=}\NormalTok{ []                         }\CommentTok{\# list of (params\_copy, energy)}
\NormalTok{    gradient\_norms }\OperatorTok{=}\NormalTok{ []}

    \ControlFlowTok{if}\NormalTok{ optimizer }\KeywordTok{in}\NormalTok{ (}\StringTok{"Adam"}\NormalTok{, }\StringTok{"GradientDescent"}\NormalTok{):}
        \CommentTok{\# PennyLane built{-}in optimiser (autograd / backprop)}
\NormalTok{        opt }\OperatorTok{=}\NormalTok{ qml.AdamOptimizer(stepsize}\OperatorTok{=}\NormalTok{learning\_rate)  }\CommentTok{\# or GradientDescentOptimizer}
        \ControlFlowTok{for}\NormalTok{ k }\KeywordTok{in} \BuiltInTok{range}\NormalTok{(maxiter):}
\NormalTok{            params, prev\_e }\OperatorTok{=}\NormalTok{ opt.step\_and\_cost(cost\_fn, params)}
\NormalTok{            energy }\OperatorTok{=} \BuiltInTok{float}\NormalTok{(cost\_fn(params))}
\NormalTok{            grad }\OperatorTok{=}\NormalTok{ qml.grad(cost\_fn)(params)          }\CommentTok{\# shape: (n\_params,)}
\NormalTok{            g\_norm }\OperatorTok{=} \BuiltInTok{float}\NormalTok{(np.linalg.norm(grad))}
\NormalTok{            trajectory.append((params.copy(), energy))}
\NormalTok{            gradient\_norms.append(g\_norm)}
            \ControlFlowTok{if}\NormalTok{ callback:}
\NormalTok{                callback(k, params, energy)}
            \CommentTok{\# Convergence: |ΔE| \textless{} tol}
            \ControlFlowTok{if} \BuiltInTok{len}\NormalTok{(trajectory) }\OperatorTok{\textgreater{}} \DecValTok{1} \KeywordTok{and} \BuiltInTok{abs}\NormalTok{(trajectory[}\OperatorTok{{-}}\DecValTok{1}\NormalTok{][}\DecValTok{1}\NormalTok{] }\OperatorTok{{-}}\NormalTok{ trajectory[}\OperatorTok{{-}}\DecValTok{2}\NormalTok{][}\DecValTok{1}\NormalTok{]) }\OperatorTok{\textless{}}\NormalTok{ tol:}
                \ControlFlowTok{break}

    \ControlFlowTok{elif}\NormalTok{ optimizer }\OperatorTok{==} \StringTok{"COBYLA"}\NormalTok{:}
        \CommentTok{\# SciPy derivative{-}free optimiser}
        \KeywordTok{def}\NormalTok{ objective(p):}
\NormalTok{            e }\OperatorTok{=} \BuiltInTok{float}\NormalTok{(cost\_fn(p))}
\NormalTok{            trajectory.append((p.copy(), e))}
            \ControlFlowTok{return}\NormalTok{ e}
\NormalTok{        scipy.optimize.minimize(objective, params, method}\OperatorTok{=}\StringTok{"COBYLA"}\NormalTok{,}
\NormalTok{                                options}\OperatorTok{=}\NormalTok{\{}\StringTok{"maxiter"}\NormalTok{: maxiter, }\StringTok{"rhobeg"}\NormalTok{: }\FloatTok{0.5}\NormalTok{\})}

    \ControlFlowTok{elif}\NormalTok{ optimizer }\OperatorTok{==} \StringTok{"SPSA"}\NormalTok{:}
        \CommentTok{\# Simultaneous Perturbation Stochastic Approximation}
\NormalTok{        a0, c0, A, alpha, gamma }\OperatorTok{=} \FloatTok{0.1}\NormalTok{, }\FloatTok{0.1}\NormalTok{, }\DecValTok{10}\NormalTok{, }\FloatTok{0.602}\NormalTok{, }\FloatTok{0.101}
        \ControlFlowTok{for}\NormalTok{ k }\KeywordTok{in} \BuiltInTok{range}\NormalTok{(}\DecValTok{1}\NormalTok{, maxiter }\OperatorTok{+} \DecValTok{1}\NormalTok{):}
\NormalTok{            a\_k }\OperatorTok{=}\NormalTok{ a0 }\OperatorTok{/}\NormalTok{ (k }\OperatorTok{+}\NormalTok{ A) }\OperatorTok{**}\NormalTok{ alpha       }\CommentTok{\# step size (decaying)}
\NormalTok{            c\_k }\OperatorTok{=}\NormalTok{ c0 }\OperatorTok{/}\NormalTok{ k }\OperatorTok{**}\NormalTok{ gamma             }\CommentTok{\# perturbation size (decaying)}
\NormalTok{            delta }\OperatorTok{=}\NormalTok{ np.random.choice([}\OperatorTok{{-}}\DecValTok{1}\NormalTok{, }\DecValTok{1}\NormalTok{], size}\OperatorTok{=}\NormalTok{n\_params)  }\CommentTok{\# ±1 vector}
\NormalTok{            f\_plus  }\OperatorTok{=} \BuiltInTok{float}\NormalTok{(cost\_fn(params }\OperatorTok{+}\NormalTok{ c\_k }\OperatorTok{*}\NormalTok{ delta))}
\NormalTok{            f\_minus }\OperatorTok{=} \BuiltInTok{float}\NormalTok{(cost\_fn(params }\OperatorTok{{-}}\NormalTok{ c\_k }\OperatorTok{*}\NormalTok{ delta))}
            \CommentTok{\# Gradient estimate: 2 circuit evals regardless of n\_params}
\NormalTok{            g\_hat }\OperatorTok{=}\NormalTok{ (f\_plus }\OperatorTok{{-}}\NormalTok{ f\_minus) }\OperatorTok{/}\NormalTok{ (}\DecValTok{2} \OperatorTok{*}\NormalTok{ c\_k }\OperatorTok{*}\NormalTok{ delta)  }\CommentTok{\# shape: (n\_params,)}
\NormalTok{            params }\OperatorTok{=}\NormalTok{ params }\OperatorTok{{-}}\NormalTok{ a\_k }\OperatorTok{*}\NormalTok{ g\_hat}
\NormalTok{            trajectory.append((params.copy(), f\_plus))}
\NormalTok{            gradient\_norms.append(}\BuiltInTok{float}\NormalTok{(np.linalg.norm(g\_hat)))}

    \CommentTok{\# Select result with lowest energy from full trajectory}
\NormalTok{    best\_idx }\OperatorTok{=} \BuiltInTok{int}\NormalTok{(np.argmin([e }\ControlFlowTok{for}\NormalTok{ \_, e }\KeywordTok{in}\NormalTok{ trajectory]))}
\NormalTok{    best\_params, best\_energy }\OperatorTok{=}\NormalTok{ trajectory[best\_idx]}
    \CommentTok{\# ... return VQEResult ...}
\end{Highlighting}
\end{Shaded}

\emph{(Full implementation: \texttt{report2\_scripts/vqe\_solver.py})}

\subsection{\texorpdfstring{Module 6:
\texttt{pes\_scanner.py}}{Module 6: pes\_scanner.py}}\label{module-6-pes_scanner.py}

\begin{Shaded}
\begin{Highlighting}[]
\CommentTok{\# {-}{-}{-} pes\_scanner.py pseudocode {-}{-}{-}}

\KeywordTok{def}\NormalTok{ scan\_pes(molecule\_factory, bond\_lengths, basis}\OperatorTok{=}\StringTok{"sto{-}3g"}\NormalTok{,}
\NormalTok{             mapping}\OperatorTok{=}\StringTok{"jordan\_wigner"}\NormalTok{, ansatz\_type}\OperatorTok{=}\StringTok{"uccsd"}\NormalTok{,}
\NormalTok{             optimizer}\OperatorTok{=}\StringTok{"Adam"}\NormalTok{, maxiter}\OperatorTok{=}\DecValTok{200}\NormalTok{,}
\NormalTok{             use\_prev\_params}\OperatorTok{=}\VariableTok{True}\NormalTok{, compute\_classical}\OperatorTok{=}\VariableTok{True}\NormalTok{, callback}\OperatorTok{=}\VariableTok{None}\NormalTok{):}

\NormalTok{    n\_points }\OperatorTok{=} \BuiltInTok{len}\NormalTok{(bond\_lengths)}
\NormalTok{    energies\_vqe }\OperatorTok{=}\NormalTok{ np.zeros(n\_points)       }\CommentTok{\# shape: (N,)}
\NormalTok{    energies\_hf  }\OperatorTok{=}\NormalTok{ np.zeros(n\_points)       }\CommentTok{\# shape: (N,)}
\NormalTok{    energies\_fci }\OperatorTok{=}\NormalTok{ np.zeros(n\_points)       }\CommentTok{\# shape: (N,)}
\NormalTok{    prev\_params }\OperatorTok{=} \VariableTok{None}

    \ControlFlowTok{for}\NormalTok{ idx, R }\KeywordTok{in} \BuiltInTok{enumerate}\NormalTok{(bond\_lengths):}
        \CommentTok{\# 1. Build geometry at bond length R\_k}
\NormalTok{        mol\_spec }\OperatorTok{=}\NormalTok{ molecule\_factory(R)}
\NormalTok{        mol\_spec.basis }\OperatorTok{=}\NormalTok{ basis}

        \CommentTok{\# 2. Run Hartree{-}Fock → MO integrals}
\NormalTok{        hf\_result }\OperatorTok{=}\NormalTok{ run\_hartree\_fock(mol\_spec)}
\NormalTok{        energies\_hf[idx] }\OperatorTok{=}\NormalTok{ hf\_result.energy\_hf}

        \CommentTok{\# 3. Build qubit Hamiltonian (JW or parity)}
\NormalTok{        qham }\OperatorTok{=}\NormalTok{ build\_hamiltonian(}
\NormalTok{            hf\_result.one\_electron\_integrals,    }\CommentTok{\# shape: (M, M)}
\NormalTok{            hf\_result.two\_electron\_integrals,    }\CommentTok{\# shape: (M, M, M, M)}
\NormalTok{            hf\_result.n\_electrons,}
\NormalTok{            hf\_result.energy\_nuc,}
\NormalTok{            mapping}\OperatorTok{=}\NormalTok{mapping,}
\NormalTok{        )}

        \CommentTok{\# 4. Build ansatz (UCCSD or hardware{-}efficient)}
        \ControlFlowTok{if}\NormalTok{ ansatz\_type }\OperatorTok{==} \StringTok{"uccsd"}\NormalTok{:}
\NormalTok{            ansatz\_fn, n\_params, \_ }\OperatorTok{=}\NormalTok{ build\_uccsd\_ansatz(}
\NormalTok{                qham.n\_qubits, hf\_result.n\_electrons)}
        \ControlFlowTok{else}\NormalTok{:}
\NormalTok{            ansatz\_fn, n\_params }\OperatorTok{=}\NormalTok{ build\_hardware\_efficient\_ansatz(}
\NormalTok{                qham.n\_qubits, n\_layers}\OperatorTok{=}\DecValTok{2}\NormalTok{)}

        \CommentTok{\# 5. Warm{-}start: reuse previous geometry\textquotesingle{}s optimal params}
        \CommentTok{\#    (only if param count unchanged; smooth PES → smooth θ*)}
\NormalTok{        init\_p }\OperatorTok{=} \VariableTok{None}
        \ControlFlowTok{if}\NormalTok{ use\_prev\_params }\KeywordTok{and}\NormalTok{ prev\_params }\KeywordTok{is} \KeywordTok{not} \VariableTok{None}\NormalTok{:}
            \ControlFlowTok{if} \BuiltInTok{len}\NormalTok{(prev\_params) }\OperatorTok{==}\NormalTok{ n\_params:}
\NormalTok{                init\_p }\OperatorTok{=}\NormalTok{ prev\_params  }\CommentTok{\# warm{-}start halves iterations}

        \CommentTok{\# 6. Run VQE}
\NormalTok{        vqe\_result }\OperatorTok{=}\NormalTok{ run\_vqe(}
\NormalTok{            qham.hamiltonian, ansatz\_fn, n\_params, qham.n\_qubits,}
\NormalTok{            init\_params}\OperatorTok{=}\NormalTok{init\_p, optimizer}\OperatorTok{=}\NormalTok{optimizer, maxiter}\OperatorTok{=}\NormalTok{maxiter)}
\NormalTok{        energies\_vqe[idx] }\OperatorTok{=}\NormalTok{ vqe\_result.optimal\_energy}
\NormalTok{        prev\_params }\OperatorTok{=}\NormalTok{ vqe\_result.optimal\_params}

        \CommentTok{\# 7. Classical benchmarks (optional)}
        \ControlFlowTok{if}\NormalTok{ compute\_classical:}
\NormalTok{            energies\_fci[idx] }\OperatorTok{=}\NormalTok{ get\_fci\_energy(mol\_spec)}
            \CommentTok{\# ... similarly for CISD, CCSD(T) ...}

    \CommentTok{\# Caveat: at large R, RHF may converge to a broken{-}symmetry solution.}
    \CommentTok{\# Active space must remain fixed across the scan for warm{-}start.}
    \ControlFlowTok{return}\NormalTok{ PESData(bond\_lengths}\OperatorTok{=}\NormalTok{bond\_lengths, energies\_vqe}\OperatorTok{=}\NormalTok{energies\_vqe,}
\NormalTok{                   energies\_hf}\OperatorTok{=}\NormalTok{energies\_hf, energies\_fci}\OperatorTok{=}\NormalTok{energies\_fci, ...)}
\end{Highlighting}
\end{Shaded}

\emph{(Full implementation: \texttt{report2\_scripts/pes\_scanner.py})}

\subsection{\texorpdfstring{Module 7:
\texttt{orbital\_grid.py}}{Module 7: orbital\_grid.py}}\label{module-7-orbital_grid.py}

\begin{Shaded}
\begin{Highlighting}[]
\CommentTok{\# {-}{-}{-} orbital\_grid.py pseudocode {-}{-}{-}}

\KeywordTok{def}\NormalTok{ evaluate\_orbital\_on\_grid(mo\_coefficients, orbital\_index,}
\NormalTok{                              atom\_coords, atom\_symbols, basis\_name,}
\NormalTok{                              grid\_extent}\OperatorTok{=}\FloatTok{5.0}\NormalTok{, grid\_spacing}\OperatorTok{=}\FloatTok{0.1}\NormalTok{):}
    \CommentTok{\# 1. Build PySCF Mole for AO evaluation}
\NormalTok{    mol }\OperatorTok{=}\NormalTok{ gto.Mole()}
\NormalTok{    atom\_str }\OperatorTok{=} \StringTok{"; "}\NormalTok{.join(}
        \SpecialStringTok{f"}\SpecialCharTok{\{}\NormalTok{sym}\SpecialCharTok{\}}\SpecialStringTok{ }\SpecialCharTok{\{}\NormalTok{c[}\DecValTok{0}\NormalTok{]}\SpecialCharTok{:.6f\}}\SpecialStringTok{ }\SpecialCharTok{\{}\NormalTok{c[}\DecValTok{1}\NormalTok{]}\SpecialCharTok{:.6f\}}\SpecialStringTok{ }\SpecialCharTok{\{}\NormalTok{c[}\DecValTok{2}\NormalTok{]}\SpecialCharTok{:.6f\}}\SpecialStringTok{"}
        \ControlFlowTok{for}\NormalTok{ sym, c }\KeywordTok{in} \BuiltInTok{zip}\NormalTok{(atom\_symbols, atom\_coords)}
\NormalTok{    )}
\NormalTok{    mol.atom }\OperatorTok{=}\NormalTok{ atom\_str}
\NormalTok{    mol.basis }\OperatorTok{=}\NormalTok{ basis\_name}
\NormalTok{    mol.unit }\OperatorTok{=} \StringTok{"Bohr"}           \CommentTok{\# PySCF uses Bohr internally}
\NormalTok{    mol.build()}

    \CommentTok{\# 2. Construct uniform 3D grid centred on molecular centroid}
\NormalTok{    center }\OperatorTok{=}\NormalTok{ np.mean(atom\_coords, axis}\OperatorTok{=}\DecValTok{0}\NormalTok{)    }\CommentTok{\# shape: (3,)}
\NormalTok{    x }\OperatorTok{=}\NormalTok{ np.arange(center[}\DecValTok{0}\NormalTok{] }\OperatorTok{{-}}\NormalTok{ grid\_extent,}
\NormalTok{                   center[}\DecValTok{0}\NormalTok{] }\OperatorTok{+}\NormalTok{ grid\_extent, grid\_spacing)}
\NormalTok{    y }\OperatorTok{=}\NormalTok{ np.arange(...)                        }\CommentTok{\# analogous for y, z}
\NormalTok{    z }\OperatorTok{=}\NormalTok{ np.arange(...)}
\NormalTok{    nx, ny, nz }\OperatorTok{=} \BuiltInTok{len}\NormalTok{(x), }\BuiltInTok{len}\NormalTok{(y), }\BuiltInTok{len}\NormalTok{(z)      }\CommentTok{\# typically \textasciitilde{}100 each}

    \CommentTok{\# Flatten to (N\_grid, 3) for batch evaluation}
\NormalTok{    xx, yy, zz }\OperatorTok{=}\NormalTok{ np.meshgrid(x, y, z, indexing}\OperatorTok{=}\StringTok{"ij"}\NormalTok{)}
\NormalTok{    coords }\OperatorTok{=}\NormalTok{ np.stack([xx.ravel(), yy.ravel(), zz.ravel()], axis}\OperatorTok{=}\DecValTok{1}\NormalTok{)}
    \CommentTok{\# coords shape: (nx*ny*nz, 3)}

    \CommentTok{\# 3. Evaluate all K AO basis functions at all grid points}
\NormalTok{    ao\_values }\OperatorTok{=}\NormalTok{ mol.eval\_gto(}\StringTok{"GTOval\_cart"}\NormalTok{, coords)}
    \CommentTok{\# ao\_values shape: (nx*ny*nz, K) — exploits Gaussian sparsity internally}

    \CommentTok{\# 4. Contract with MO coefficient column p: φ\_p(r) = Σ\_μ C\_\{μp\} χ\_μ(r)}
\NormalTok{    mo\_col }\OperatorTok{=}\NormalTok{ mo\_coefficients[:, orbital\_index]   }\CommentTok{\# shape: (K,)}
\NormalTok{    mo\_values }\OperatorTok{=}\NormalTok{ ao\_values }\OperatorTok{@}\NormalTok{ mo\_col               }\CommentTok{\# shape: (nx*ny*nz,)}

    \CommentTok{\# 5. Reshape to 3D volume}
\NormalTok{    volume }\OperatorTok{=}\NormalTok{ mo\_values.reshape(nx, ny, nz)       }\CommentTok{\# shape: (N\_x, N\_y, N\_z)}

    \ControlFlowTok{return}\NormalTok{ OrbitalVolumeData(}
\NormalTok{        values}\OperatorTok{=}\NormalTok{volume,}
\NormalTok{        grid\_origin}\OperatorTok{=}\NormalTok{(}\BuiltInTok{float}\NormalTok{(x[}\DecValTok{0}\NormalTok{]), }\BuiltInTok{float}\NormalTok{(y[}\DecValTok{0}\NormalTok{]), }\BuiltInTok{float}\NormalTok{(z[}\DecValTok{0}\NormalTok{])),}
\NormalTok{        grid\_spacing}\OperatorTok{=}\NormalTok{(grid\_spacing, grid\_spacing, grid\_spacing),}
\NormalTok{        grid\_shape}\OperatorTok{=}\NormalTok{(nx, ny, nz),}
\NormalTok{        orbital\_index}\OperatorTok{=}\NormalTok{orbital\_index,}
\NormalTok{    )}


\KeywordTok{def}\NormalTok{ evaluate\_electron\_density(mo\_coefficients, n\_occupied, ...):}
    \CommentTok{\# Electron density: ρ(r) = Σ\_\{i=0\}\^{}\{n\_occ{-}1\} 2 * |φ\_i(r)|²}
    \CommentTok{\# Factor of 2: RHF double occupation}
\NormalTok{    density }\OperatorTok{=}\NormalTok{ np.zeros((nx, ny, nz))             }\CommentTok{\# shape: (N\_x, N\_y, N\_z)}
    \ControlFlowTok{for}\NormalTok{ i }\KeywordTok{in} \BuiltInTok{range}\NormalTok{(n\_occupied):}
\NormalTok{        orb\_data }\OperatorTok{=}\NormalTok{ evaluate\_orbital\_on\_grid(mo\_coefficients, i, ...)}
\NormalTok{        density }\OperatorTok{+=} \FloatTok{2.0} \OperatorTok{*}\NormalTok{ orb\_data.values }\OperatorTok{**} \DecValTok{2}
    \ControlFlowTok{return}\NormalTok{ OrbitalVolumeData(values}\OperatorTok{=}\NormalTok{density, ..., orbital\_index}\OperatorTok{={-}}\DecValTok{1}\NormalTok{)}

    \CommentTok{\# UNIT CAVEAT: If atom\_coords are in Ångström, convert first:}
    \CommentTok{\# atom\_coords\_bohr = atom\_coords * 1.8897259886}
\end{Highlighting}
\end{Shaded}

\emph{(Full implementation: \texttt{report2\_scripts/orbital\_grid.py})}

\subsection{\texorpdfstring{Module 8:
\texttt{isosurface.py}}{Module 8: isosurface.py}}\label{module-8-isosurface.py}

\begin{Shaded}
\begin{Highlighting}[]
\CommentTok{\# {-}{-}{-} isosurface.py pseudocode {-}{-}{-}}

\KeywordTok{def}\NormalTok{ marching\_cubes(volume, isovalue, grid\_origin, grid\_spacing):}
    \CommentTok{"""}
\CommentTok{    Lorensen \& Cline (1987) via scikit{-}image.}
\CommentTok{    1. Each voxel (2×2×2 corners) classified inside/outside → 8{-}bit index}
\CommentTok{    2. Index selects triangulation from 256{-}entry lookup table (15 unique)}
\CommentTok{    3. Edge crossings linearly interpolated between corner values}
\CommentTok{    4. Vertex normals from scalar field gradient}
\CommentTok{    """}
    \CommentTok{\# scikit{-}image returns (verts, faces, normals, values)}
\NormalTok{    verts, faces, normals, \_ }\OperatorTok{=}\NormalTok{ sk\_marching\_cubes(}
\NormalTok{        volume,                          }\CommentTok{\# shape: (N\_x, N\_y, N\_z)}
\NormalTok{        level}\OperatorTok{=}\NormalTok{isovalue,}
\NormalTok{        spacing}\OperatorTok{=}\NormalTok{grid\_spacing,            }\CommentTok{\# (dx, dy, dz) in Bohr}
\NormalTok{    )}
    \CommentTok{\# verts shape: (N\_v, 3) in grid{-}spacing units}
    \CommentTok{\# faces shape: (N\_f, 3) triangle vertex indices}
    \CommentTok{\# normals shape: (N\_v, 3)}

    \CommentTok{\# Translate to world coordinates}
\NormalTok{    origin }\OperatorTok{=}\NormalTok{ np.array(grid\_origin)           }\CommentTok{\# shape: (3,)}
\NormalTok{    verts }\OperatorTok{=}\NormalTok{ verts }\OperatorTok{+}\NormalTok{ origin[}\VariableTok{None}\NormalTok{, :]          }\CommentTok{\# broadcast add}

\NormalTok{    signs }\OperatorTok{=}\NormalTok{ np.ones(}\BuiltInTok{len}\NormalTok{(verts), dtype}\OperatorTok{=}\NormalTok{np.float32)  }\CommentTok{\# +1 for positive lobe}
    \ControlFlowTok{return}\NormalTok{ IsosurfaceMesh(}
\NormalTok{        vertices}\OperatorTok{=}\NormalTok{verts.astype(np.float32),}
\NormalTok{        normals}\OperatorTok{=}\NormalTok{normals.astype(np.float32),}
\NormalTok{        faces}\OperatorTok{=}\NormalTok{faces.astype(np.uint32),}
\NormalTok{        signs}\OperatorTok{=}\NormalTok{signs,}
\NormalTok{    )}


\KeywordTok{def}\NormalTok{ extract\_dual\_isosurface(volume, isovalue, grid\_origin, grid\_spacing):}
    \CommentTok{"""Extract both +c and {-}c surfaces for positive/negative lobes."""}
    \CommentTok{\# Positive lobe (sign = +1)}
\NormalTok{    mesh\_pos }\OperatorTok{=}\NormalTok{ marching\_cubes(volume, }\OperatorTok{+}\NormalTok{isovalue, grid\_origin, grid\_spacing)}
    \CommentTok{\# Negative lobe (sign = {-}1)}
\NormalTok{    mesh\_neg }\OperatorTok{=}\NormalTok{ marching\_cubes(volume, }\OperatorTok{{-}}\NormalTok{isovalue, grid\_origin, grid\_spacing)}
\NormalTok{    mesh\_neg.signs[:] }\OperatorTok{=} \OperatorTok{{-}}\FloatTok{1.0}                 \CommentTok{\# tag as negative lobe}

\NormalTok{    n\_verts\_pos }\OperatorTok{=} \BuiltInTok{len}\NormalTok{(mesh\_pos.vertices)}

    \CommentTok{\# Concatenate meshes; offset negative face indices by |V\_pos|}
\NormalTok{    combined\_verts   }\OperatorTok{=}\NormalTok{ np.concatenate([mesh\_pos.vertices, mesh\_neg.vertices])}
\NormalTok{    combined\_normals }\OperatorTok{=}\NormalTok{ np.concatenate([mesh\_pos.normals, mesh\_neg.normals])}
\NormalTok{    combined\_signs   }\OperatorTok{=}\NormalTok{ np.concatenate([mesh\_pos.signs, mesh\_neg.signs])}
\NormalTok{    neg\_faces\_offset }\OperatorTok{=}\NormalTok{ mesh\_neg.faces }\OperatorTok{+}\NormalTok{ n\_verts\_pos   }\CommentTok{\# index offset}
\NormalTok{    combined\_faces   }\OperatorTok{=}\NormalTok{ np.concatenate([mesh\_pos.faces, neg\_faces\_offset])}

    \ControlFlowTok{return}\NormalTok{ IsosurfaceMesh(vertices}\OperatorTok{=}\NormalTok{combined\_verts, normals}\OperatorTok{=}\NormalTok{combined\_normals,}
\NormalTok{                          faces}\OperatorTok{=}\NormalTok{combined\_faces, signs}\OperatorTok{=}\NormalTok{combined\_signs)}

    \CommentTok{\# EDGE CASE: marching\_cubes fails if isovalue outside data range.}
    \CommentTok{\# For negative surface, need volume.min() \textless{} {-}c AND volume.max() \textgreater{} {-}c.}
    \CommentTok{\# Very small isovalues near 0 can produce degenerate triangles.}
\end{Highlighting}
\end{Shaded}

\emph{(Full implementation: \texttt{report2\_scripts/isosurface.py})}

\subsection{\texorpdfstring{Module 9:
\texttt{renderer.py}}{Module 9: renderer.py}}\label{module-9-renderer.py}

\begin{Shaded}
\begin{Highlighting}[]
\CommentTok{\# {-}{-}{-} renderer.py pseudocode {-}{-}{-}}

\KeywordTok{def}\NormalTok{ init\_opengl\_context(width}\OperatorTok{=}\DecValTok{1400}\NormalTok{, height}\OperatorTok{=}\DecValTok{900}\NormalTok{, title}\OperatorTok{=}\StringTok{"VQE Molecular Explorer"}\NormalTok{):}
\NormalTok{    glfw.init()}
\NormalTok{    glfw.window\_hint(glfw.CONTEXT\_VERSION\_MAJOR, }\DecValTok{3}\NormalTok{)}
\NormalTok{    glfw.window\_hint(glfw.CONTEXT\_VERSION\_MINOR, }\DecValTok{3}\NormalTok{)}
\NormalTok{    glfw.window\_hint(glfw.OPENGL\_PROFILE, glfw.OPENGL\_CORE\_PROFILE)}
\NormalTok{    glfw.window\_hint(glfw.SAMPLES, }\DecValTok{4}\NormalTok{)        }\CommentTok{\# 4× MSAA}
\NormalTok{    window }\OperatorTok{=}\NormalTok{ glfw.create\_window(width, height, title, }\VariableTok{None}\NormalTok{, }\VariableTok{None}\NormalTok{)}
\NormalTok{    glfw.make\_context\_current(window)}
\NormalTok{    glEnable(GL\_DEPTH\_TEST)}
\NormalTok{    glEnable(GL\_MULTISAMPLE)}
\NormalTok{    glEnable(GL\_BLEND)}
\NormalTok{    glBlendFunc(GL\_SRC\_ALPHA, GL\_ONE\_MINUS\_SRC\_ALPHA)}
    \ControlFlowTok{return}\NormalTok{ window}


\KeywordTok{def}\NormalTok{ create\_orbital\_vao(mesh):}
    \CommentTok{"""Upload IsosurfaceMesh to GPU.}
\CommentTok{    Vertex layout: [x, y, z, nx, ny, nz, sign] — 7 float32, 28 bytes/vertex}
\CommentTok{    Position: location 0, Normal: location 1, Sign: location 2.}
\CommentTok{    Uses GL\_DYNAMIC\_DRAW (mesh changes with bond length).}
\CommentTok{    """}
\NormalTok{    n\_verts }\OperatorTok{=} \BuiltInTok{len}\NormalTok{(mesh.vertices)}
\NormalTok{    vertices }\OperatorTok{=}\NormalTok{ np.zeros((n\_verts, }\DecValTok{7}\NormalTok{), dtype}\OperatorTok{=}\NormalTok{np.float32)}
\NormalTok{    vertices[:, }\DecValTok{0}\NormalTok{:}\DecValTok{3}\NormalTok{] }\OperatorTok{=}\NormalTok{ mesh.vertices         }\CommentTok{\# position (x, y, z)}
\NormalTok{    vertices[:, }\DecValTok{3}\NormalTok{:}\DecValTok{6}\NormalTok{] }\OperatorTok{=}\NormalTok{ mesh.normals          }\CommentTok{\# normal (nx, ny, nz)}
\NormalTok{    vertices[:, }\DecValTok{6}\NormalTok{]   }\OperatorTok{=}\NormalTok{ mesh.signs            }\CommentTok{\# lobe sign (±1)}

\NormalTok{    vao }\OperatorTok{=}\NormalTok{ glGenVertexArrays(}\DecValTok{1}\NormalTok{)}
\NormalTok{    vbo }\OperatorTok{=}\NormalTok{ glGenBuffers(}\DecValTok{1}\NormalTok{)}
\NormalTok{    ebo }\OperatorTok{=}\NormalTok{ glGenBuffers(}\DecValTok{1}\NormalTok{)}
\NormalTok{    glBindVertexArray(vao)}

\NormalTok{    glBindBuffer(GL\_ARRAY\_BUFFER, vbo)}
\NormalTok{    glBufferData(GL\_ARRAY\_BUFFER, vertices.nbytes, vertices, GL\_DYNAMIC\_DRAW)}
\NormalTok{    glBindBuffer(GL\_ELEMENT\_ARRAY\_BUFFER, ebo)}
\NormalTok{    glBufferData(GL\_ELEMENT\_ARRAY\_BUFFER, mesh.faces.nbytes, mesh.faces, GL\_DYNAMIC\_DRAW)}

\NormalTok{    stride }\OperatorTok{=} \DecValTok{7} \OperatorTok{*} \DecValTok{4}                           \CommentTok{\# 28 bytes}
    \CommentTok{\# Attribute 0: position (3 floats, offset 0)}
\NormalTok{    glVertexAttribPointer(}\DecValTok{0}\NormalTok{, }\DecValTok{3}\NormalTok{, GL\_FLOAT, GL\_FALSE, stride, ctypes.c\_void\_p(}\DecValTok{0}\NormalTok{))}
\NormalTok{    glEnableVertexAttribArray(}\DecValTok{0}\NormalTok{)}
    \CommentTok{\# Attribute 1: normal (3 floats, offset 12)}
\NormalTok{    glVertexAttribPointer(}\DecValTok{1}\NormalTok{, }\DecValTok{3}\NormalTok{, GL\_FLOAT, GL\_FALSE, stride, ctypes.c\_void\_p(}\DecValTok{12}\NormalTok{))}
\NormalTok{    glEnableVertexAttribArray(}\DecValTok{1}\NormalTok{)}
    \CommentTok{\# Attribute 2: sign (1 float, offset 24)}
\NormalTok{    glVertexAttribPointer(}\DecValTok{2}\NormalTok{, }\DecValTok{1}\NormalTok{, GL\_FLOAT, GL\_FALSE, stride, ctypes.c\_void\_p(}\DecValTok{24}\NormalTok{))}
\NormalTok{    glEnableVertexAttribArray(}\DecValTok{2}\NormalTok{)}

\NormalTok{    glBindVertexArray(}\DecValTok{0}\NormalTok{)}
    \ControlFlowTok{return}\NormalTok{ vao, vbo, ebo, }\BuiltInTok{len}\NormalTok{(mesh.faces) }\OperatorTok{*} \DecValTok{3}


\KeywordTok{def}\NormalTok{ update\_orbital\_vao(vbo, ebo, mesh):}
    \CommentTok{"""Orphan old buffer → upload new mesh data (avoids GPU sync stalls)."""}
    \CommentTok{\# ... rebuild vertex array same layout as create\_orbital\_vao ...}
\NormalTok{    glBindBuffer(GL\_ARRAY\_BUFFER, vbo)}
\NormalTok{    glBufferData(GL\_ARRAY\_BUFFER, vertices.nbytes, vertices, GL\_DYNAMIC\_DRAW)}
\NormalTok{    glBindBuffer(GL\_ELEMENT\_ARRAY\_BUFFER, ebo)}
\NormalTok{    glBufferData(GL\_ELEMENT\_ARRAY\_BUFFER, mesh.faces.nbytes, mesh.faces, GL\_DYNAMIC\_DRAW)}
    \ControlFlowTok{return} \BuiltInTok{len}\NormalTok{(mesh.faces) }\OperatorTok{*} \DecValTok{3}


\KeywordTok{def}\NormalTok{ create\_pes\_curve\_vao(bond\_lengths, energies\_dict):}
    \CommentTok{"""Each method (HF, CISD, VQE, FCI) → coloured GL\_LINE\_STRIP.}
\CommentTok{    Vertex layout: [R, E, 0, r, g, b, a] — 7 floats per vertex.}
\CommentTok{    """}
\NormalTok{    METHOD\_COLORS }\OperatorTok{=}\NormalTok{ \{}\StringTok{"HF"}\NormalTok{: (}\FloatTok{0.5}\NormalTok{, }\FloatTok{0.5}\NormalTok{, }\FloatTok{0.5}\NormalTok{, }\DecValTok{1}\NormalTok{), }\StringTok{"CISD"}\NormalTok{: (}\DecValTok{0}\NormalTok{, }\FloatTok{0.8}\NormalTok{, }\DecValTok{0}\NormalTok{, }\DecValTok{1}\NormalTok{),}
                     \StringTok{"VQE"}\NormalTok{: (}\FloatTok{0.2}\NormalTok{, }\FloatTok{0.4}\NormalTok{, }\FloatTok{0.9}\NormalTok{, }\DecValTok{1}\NormalTok{), }\StringTok{"FCI"}\NormalTok{: (}\DecValTok{1}\NormalTok{, }\FloatTok{0.84}\NormalTok{, }\DecValTok{0}\NormalTok{, }\DecValTok{1}\NormalTok{),}
                     \StringTok{"CCSD(T)"}\NormalTok{: (}\FloatTok{0.6}\NormalTok{, }\FloatTok{0.2}\NormalTok{, }\FloatTok{0.8}\NormalTok{, }\DecValTok{1}\NormalTok{)\}}
\NormalTok{    all\_verts }\OperatorTok{=}\NormalTok{ []}
    \ControlFlowTok{for}\NormalTok{ method, energies }\KeywordTok{in}\NormalTok{ energies\_dict.items():}
\NormalTok{        rgba }\OperatorTok{=}\NormalTok{ METHOD\_COLORS.get(method, (}\FloatTok{0.7}\NormalTok{, }\FloatTok{0.7}\NormalTok{, }\FloatTok{0.7}\NormalTok{, }\DecValTok{1}\NormalTok{))}
        \ControlFlowTok{for}\NormalTok{ R, E }\KeywordTok{in} \BuiltInTok{zip}\NormalTok{(bond\_lengths, energies):}
\NormalTok{            all\_verts.append([R, E, }\FloatTok{0.0}\NormalTok{, }\OperatorTok{*}\NormalTok{rgba])  }\CommentTok{\# shape per vert: (7,)}
    \CommentTok{\# ... upload to VBO ...}


\KeywordTok{def}\NormalTok{ draw\_molecule(atom\_positions, atom\_symbols, bonds, shader\_program):}
    \CommentTok{\# CPK colours and scaled vdW radii}
\NormalTok{    ATOM\_COLORS }\OperatorTok{=}\NormalTok{ \{}\StringTok{"H"}\NormalTok{: (}\DecValTok{1}\NormalTok{,}\DecValTok{1}\NormalTok{,}\DecValTok{1}\NormalTok{), }\StringTok{"C"}\NormalTok{: (}\FloatTok{0.5}\NormalTok{,}\FloatTok{0.5}\NormalTok{,}\FloatTok{0.5}\NormalTok{), }\StringTok{"O"}\NormalTok{: (}\DecValTok{1}\NormalTok{,}\DecValTok{0}\NormalTok{,}\DecValTok{0}\NormalTok{), }\StringTok{"Li"}\NormalTok{: (}\FloatTok{0.6}\NormalTok{,}\FloatTok{0.2}\NormalTok{,}\FloatTok{0.8}\NormalTok{)\}}
\NormalTok{    VDW\_RADII   }\OperatorTok{=}\NormalTok{ \{}\StringTok{"H"}\NormalTok{: }\FloatTok{1.20}\NormalTok{, }\StringTok{"C"}\NormalTok{: }\FloatTok{1.70}\NormalTok{, }\StringTok{"O"}\NormalTok{: }\FloatTok{1.52}\NormalTok{, }\StringTok{"Li"}\NormalTok{: }\FloatTok{1.82}\NormalTok{\}}
    \ControlFlowTok{for}\NormalTok{ pos, sym }\KeywordTok{in} \BuiltInTok{zip}\NormalTok{(atom\_positions, atom\_symbols):}
\NormalTok{        radius }\OperatorTok{=}\NormalTok{ VDW\_RADII.get(sym, }\FloatTok{1.5}\NormalTok{) }\OperatorTok{*} \FloatTok{0.3}   \CommentTok{\# scaled ×0.3}
\NormalTok{        model }\OperatorTok{=}\NormalTok{ \_make\_sphere\_model\_matrix(pos, radius)}
        \CommentTok{\# ... set uniforms, draw instanced icosphere ...}
    \ControlFlowTok{for}\NormalTok{ i, j }\KeywordTok{in}\NormalTok{ bonds:}
\NormalTok{        model }\OperatorTok{=}\NormalTok{ \_make\_cylinder\_model\_matrix(atom\_positions[i], atom\_positions[j], }\FloatTok{0.1}\NormalTok{)}
        \CommentTok{\# ... draw unit cylinder ...}


\KeywordTok{class}\NormalTok{ ArcballCamera:}
    \CommentTok{"""Arcball: yaw/pitch from mouse drag, zoom from scroll."""}
    \KeywordTok{def} \FunctionTok{\_\_init\_\_}\NormalTok{(}\VariableTok{self}\NormalTok{, position}\OperatorTok{=}\VariableTok{None}\NormalTok{, target}\OperatorTok{=}\VariableTok{None}\NormalTok{):}
        \VariableTok{self}\NormalTok{.yaw, }\VariableTok{self}\NormalTok{.pitch }\OperatorTok{=} \OperatorTok{{-}}\FloatTok{90.0}\NormalTok{, }\FloatTok{20.0}
        \VariableTok{self}\NormalTok{.zoom }\OperatorTok{=} \FloatTok{10.0}
        \CommentTok{\# ...}

    \KeywordTok{def}\NormalTok{ on\_mouse\_drag(}\VariableTok{self}\NormalTok{, dx, dy):}
        \VariableTok{self}\NormalTok{.yaw   }\OperatorTok{+=}\NormalTok{ dx }\OperatorTok{*} \FloatTok{0.3}               \CommentTok{\# Δyaw = 0.3 * dx}
        \VariableTok{self}\NormalTok{.pitch }\OperatorTok{{-}=}\NormalTok{ dy }\OperatorTok{*} \FloatTok{0.3}               \CommentTok{\# Δpitch = {-}0.3 * dy}
        \VariableTok{self}\NormalTok{.pitch  }\OperatorTok{=}\NormalTok{ np.clip(}\VariableTok{self}\NormalTok{.pitch, }\OperatorTok{{-}}\DecValTok{89}\NormalTok{, }\DecValTok{89}\NormalTok{)}
        \VariableTok{self}\NormalTok{.\_update()}

    \KeywordTok{def}\NormalTok{ on\_scroll(}\VariableTok{self}\NormalTok{, delta):}
        \VariableTok{self}\NormalTok{.zoom }\OperatorTok{{-}=}\NormalTok{ delta }\OperatorTok{*} \FloatTok{0.5}
        \VariableTok{self}\NormalTok{.zoom }\OperatorTok{=} \BuiltInTok{max}\NormalTok{(}\FloatTok{1.0}\NormalTok{, }\VariableTok{self}\NormalTok{.zoom)}
        \VariableTok{self}\NormalTok{.\_update()}

    \KeywordTok{def}\NormalTok{ \_update(}\VariableTok{self}\NormalTok{):}
        \CommentTok{\# Recompute position in spherical coords relative to target}
\NormalTok{        yr, pr }\OperatorTok{=}\NormalTok{ np.radians(}\VariableTok{self}\NormalTok{.yaw), np.radians(}\VariableTok{self}\NormalTok{.pitch)}
        \VariableTok{self}\NormalTok{.position }\OperatorTok{=} \VariableTok{self}\NormalTok{.target }\OperatorTok{+} \VariableTok{self}\NormalTok{.zoom }\OperatorTok{*}\NormalTok{ np.array([}
\NormalTok{            np.cos(pr) }\OperatorTok{*}\NormalTok{ np.cos(yr), np.sin(pr), np.cos(pr) }\OperatorTok{*}\NormalTok{ np.sin(yr)])}

    \KeywordTok{def}\NormalTok{ get\_view\_matrix(}\VariableTok{self}\NormalTok{):}
        \ControlFlowTok{return}\NormalTok{ \_look\_at(}\VariableTok{self}\NormalTok{.position, }\VariableTok{self}\NormalTok{.target, up}\OperatorTok{=}\NormalTok{[}\DecValTok{0}\NormalTok{, }\DecValTok{1}\NormalTok{, }\DecValTok{0}\NormalTok{])}

    \KeywordTok{def}\NormalTok{ get\_projection\_matrix(}\VariableTok{self}\NormalTok{, aspect, fov}\OperatorTok{=}\FloatTok{45.0}\NormalTok{):}
        \ControlFlowTok{return}\NormalTok{ \_perspective(fov, aspect, near}\OperatorTok{=}\FloatTok{0.1}\NormalTok{, far}\OperatorTok{=}\FloatTok{200.0}\NormalTok{)}
\end{Highlighting}
\end{Shaded}

\emph{(Full implementation: \texttt{report2\_scripts/renderer.py})}

\subsection{\texorpdfstring{Module 10:
\texttt{hardware\_runner.py}}{Module 10: hardware\_runner.py}}\label{module-10-hardware_runner.py}

\begin{Shaded}
\begin{Highlighting}[]
\CommentTok{\# {-}{-}{-} hardware\_runner.py pseudocode {-}{-}{-}}

\KeywordTok{def}\NormalTok{ create\_ibm\_device(n\_qubits, backend\_name}\OperatorTok{=}\StringTok{"ibm\_brisbane"}\NormalTok{,}
\NormalTok{                      shots}\OperatorTok{=}\DecValTok{8192}\NormalTok{, ibm\_token}\OperatorTok{=}\VariableTok{None}\NormalTok{):}
\NormalTok{    token }\OperatorTok{=}\NormalTok{ ibm\_token }\KeywordTok{or}\NormalTok{ os.environ.get(}\StringTok{"IBMQ\_TOKEN"}\NormalTok{)}
\NormalTok{    dev }\OperatorTok{=}\NormalTok{ qml.device(}
        \StringTok{"qiskit.ibmq"}\NormalTok{,}
\NormalTok{        wires}\OperatorTok{=}\NormalTok{n\_qubits,}
\NormalTok{        backend}\OperatorTok{=}\NormalTok{backend\_name,}
\NormalTok{        shots}\OperatorTok{=}\NormalTok{shots,}
\NormalTok{        ibmqx\_token}\OperatorTok{=}\NormalTok{token,}
\NormalTok{    )}
    \ControlFlowTok{return}\NormalTok{ dev}


\KeywordTok{def}\NormalTok{ create\_ionq\_device(n\_qubits, backend}\OperatorTok{=}\StringTok{"ionq\_harmony"}\NormalTok{,}
\NormalTok{                       shots}\OperatorTok{=}\DecValTok{1024}\NormalTok{, api\_key}\OperatorTok{=}\VariableTok{None}\NormalTok{):}
\NormalTok{    key }\OperatorTok{=}\NormalTok{ api\_key }\KeywordTok{or}\NormalTok{ os.environ.get(}\StringTok{"IONQ\_API\_KEY"}\NormalTok{)}
\NormalTok{    dev }\OperatorTok{=}\NormalTok{ qml.device(}
        \StringTok{"ionq.simulator"}\NormalTok{,           }\CommentTok{\# or "ionq.qpu" for real hardware}
\NormalTok{        wires}\OperatorTok{=}\NormalTok{n\_qubits,}
\NormalTok{        shots}\OperatorTok{=}\NormalTok{shots,}
\NormalTok{        api\_key}\OperatorTok{=}\NormalTok{key,}
\NormalTok{    )}
    \ControlFlowTok{return}\NormalTok{ dev}


\KeywordTok{def}\NormalTok{ estimate\_shot\_budget(hamiltonian, target\_precision}\OperatorTok{=}\FloatTok{1.6e{-}3}\NormalTok{):}
    \CommentTok{\# Variance bound: Var[E] ≤ Σ|c\_k|² / S  (independent measurements)}
\NormalTok{    coeffs }\OperatorTok{=}\NormalTok{ np.array(hamiltonian.coeffs)    }\CommentTok{\# shape: (n\_paulis,)}
\NormalTok{    variance\_bound }\OperatorTok{=}\NormalTok{ np.}\BuiltInTok{sum}\NormalTok{(coeffs }\OperatorTok{**} \DecValTok{2}\NormalTok{)}
    \CommentTok{\# Need: √Var ≤ ε ⇒ S ≥ Σ|c\_k|² / ε²}
\NormalTok{    shots\_needed }\OperatorTok{=} \BuiltInTok{int}\NormalTok{(np.ceil(variance\_bound }\OperatorTok{/}\NormalTok{ target\_precision }\OperatorTok{**} \DecValTok{2}\NormalTok{))}
\NormalTok{    shots\_needed }\OperatorTok{=} \BuiltInTok{max}\NormalTok{(}\DecValTok{1000}\NormalTok{, }\BuiltInTok{min}\NormalTok{(shots\_needed, }\DecValTok{10\_000\_000}\NormalTok{))  }\CommentTok{\# clamp}
    \ControlFlowTok{return}\NormalTok{ shots\_needed}
    \CommentTok{\# H₂ (STO{-}3G): Σ|c\_k|² ≈ 0.3, ε = 1.6e{-}3 ⇒ S ≈ 117,000}
    \CommentTok{\# Pauli grouping into commuting cliques reduces actual requirement.}
\end{Highlighting}
\end{Shaded}

\emph{(Full implementation:
\texttt{report2\_scripts/hardware\_runner.py})}

\subsection{\texorpdfstring{Module 11:
\texttt{benchmarker.py}}{Module 11: benchmarker.py}}\label{module-11-benchmarker.py}

\begin{Shaded}
\begin{Highlighting}[]
\CommentTok{\# {-}{-}{-} benchmarker.py pseudocode {-}{-}{-}}

\KeywordTok{def}\NormalTok{ run\_benchmark(pes\_data):}
\NormalTok{    e\_hf  }\OperatorTok{=}\NormalTok{ pes\_data.energies\_hf       }\CommentTok{\# shape: (N,)}
\NormalTok{    e\_fci }\OperatorTok{=}\NormalTok{ pes\_data.energies\_fci      }\CommentTok{\# shape: (N,)}
\NormalTok{    denom }\OperatorTok{=}\NormalTok{ e\_fci }\OperatorTok{{-}}\NormalTok{ e\_hf               }\CommentTok{\# shape: (N,) — negative values}

\NormalTok{    methods }\OperatorTok{=}\NormalTok{ [}\StringTok{"HF"}\NormalTok{, }\StringTok{"VQE"}\NormalTok{]}
\NormalTok{    energies }\OperatorTok{=}\NormalTok{ \{}\StringTok{"HF"}\NormalTok{: e\_hf, }\StringTok{"VQE"}\NormalTok{: pes\_data.energies\_vqe, }\StringTok{"FCI"}\NormalTok{: e\_fci\}}
    \ControlFlowTok{if}\NormalTok{ pes\_data.energies\_cisd }\KeywordTok{is} \KeywordTok{not} \VariableTok{None}\NormalTok{:}
\NormalTok{        methods.append(}\StringTok{"CISD"}\NormalTok{)}
\NormalTok{        energies[}\StringTok{"CISD"}\NormalTok{] }\OperatorTok{=}\NormalTok{ pes\_data.energies\_cisd}
    \ControlFlowTok{if}\NormalTok{ pes\_data.energies\_ccsd\_t }\KeywordTok{is} \KeywordTok{not} \VariableTok{None}\NormalTok{:}
\NormalTok{        methods.append(}\StringTok{"CCSD(T)"}\NormalTok{)}
\NormalTok{        energies[}\StringTok{"CCSD(T)"}\NormalTok{] }\OperatorTok{=}\NormalTok{ pes\_data.energies\_ccsd\_t}

\NormalTok{    correlation\_recovery }\OperatorTok{=}\NormalTok{ \{\}}
\NormalTok{    errors\_vs\_fci }\OperatorTok{=}\NormalTok{ \{\}}
    \ControlFlowTok{for}\NormalTok{ method }\KeywordTok{in}\NormalTok{ methods:}
\NormalTok{        e\_m }\OperatorTok{=}\NormalTok{ energies[method]                     }\CommentTok{\# shape: (N,)}
        \CommentTok{\# Correlation energy recovery: r = (E\_method {-} E\_HF) / (E\_FCI {-} E\_HF)}
        \CommentTok{\# Guard division{-}by{-}zero where |E\_FCI {-} E\_HF| \textless{} 1e{-}10}
\NormalTok{        recovery }\OperatorTok{=}\NormalTok{ np.where(}
\NormalTok{            np.}\BuiltInTok{abs}\NormalTok{(denom) }\OperatorTok{\textgreater{}} \FloatTok{1e{-}10}\NormalTok{,}
\NormalTok{            (e\_m }\OperatorTok{{-}}\NormalTok{ e\_hf) }\OperatorTok{/}\NormalTok{ denom,}
            \FloatTok{0.0}\NormalTok{,}
\NormalTok{        )                                          }\CommentTok{\# shape: (N,)}
\NormalTok{        correlation\_recovery[method] }\OperatorTok{=}\NormalTok{ recovery}
        \CommentTok{\# Absolute error vs FCI: |E\_method {-} E\_FCI|}
\NormalTok{        errors\_vs\_fci[method] }\OperatorTok{=}\NormalTok{ np.}\BuiltInTok{abs}\NormalTok{(e\_m }\OperatorTok{{-}}\NormalTok{ e\_fci)  }\CommentTok{\# shape: (N,)}

    \ControlFlowTok{return}\NormalTok{ BenchmarkResult(}
\NormalTok{        bond\_lengths}\OperatorTok{=}\NormalTok{pes\_data.bond\_lengths,}
\NormalTok{        methods}\OperatorTok{=}\NormalTok{methods, energies}\OperatorTok{=}\NormalTok{energies,}
\NormalTok{        correlation\_recovery}\OperatorTok{=}\NormalTok{correlation\_recovery,}
\NormalTok{        errors\_vs\_fci}\OperatorTok{=}\NormalTok{errors\_vs\_fci,}
\NormalTok{    )}


\KeywordTok{def}\NormalTok{ print\_benchmark\_table(result, bond\_lengths\_subset}\OperatorTok{=}\VariableTok{None}\NormalTok{):}
    \CommentTok{\# Select 5 evenly spaced bond lengths if no subset given}
    \ControlFlowTok{if}\NormalTok{ bond\_lengths\_subset }\KeywordTok{is} \VariableTok{None}\NormalTok{:}
\NormalTok{        indices }\OperatorTok{=}\NormalTok{ np.linspace(}\DecValTok{0}\NormalTok{, }\BuiltInTok{len}\NormalTok{(result.bond\_lengths) }\OperatorTok{{-}} \DecValTok{1}\NormalTok{, }\DecValTok{5}\NormalTok{, dtype}\OperatorTok{=}\BuiltInTok{int}\NormalTok{)}
    \ControlFlowTok{else}\NormalTok{:}
\NormalTok{        indices }\OperatorTok{=}\NormalTok{ [np.argmin(np.}\BuiltInTok{abs}\NormalTok{(result.bond\_lengths }\OperatorTok{{-}}\NormalTok{ bl))}
                   \ControlFlowTok{for}\NormalTok{ bl }\KeywordTok{in}\NormalTok{ bond\_lengths\_subset]}

    \CommentTok{\# Print header: R (Å) | HF | CISD | VQE | FCI | |VQE{-}FCI|}
\NormalTok{    header }\OperatorTok{=} \SpecialStringTok{f"}\SpecialCharTok{\{}\StringTok{\textquotesingle{}R (Å)\textquotesingle{}}\SpecialCharTok{:\textgreater{}8\}}\SpecialStringTok{"}
    \ControlFlowTok{for}\NormalTok{ m }\KeywordTok{in}\NormalTok{ result.methods }\OperatorTok{+}\NormalTok{ [}\StringTok{"FCI"}\NormalTok{]:}
\NormalTok{        header }\OperatorTok{+=} \SpecialStringTok{f" | }\SpecialCharTok{\{}\NormalTok{m}\SpecialCharTok{:\textgreater{}12\}}\SpecialStringTok{"}
\NormalTok{    header }\OperatorTok{+=} \SpecialStringTok{f" | }\SpecialCharTok{\{}\StringTok{\textquotesingle{}|VQE{-}FCI|\textquotesingle{}}\SpecialCharTok{:\textgreater{}12\}}\SpecialStringTok{"}
    \BuiltInTok{print}\NormalTok{(header)}
    \CommentTok{\# ... print rows for each bond length in indices ...}
\end{Highlighting}
\end{Shaded}

\emph{(Full implementation: \texttt{report2\_scripts/benchmarker.py})}

\subsection{\texorpdfstring{Module 12:
\texttt{main.py}}{Module 12: main.py}}\label{module-12-main.py}

\begin{Shaded}
\begin{Highlighting}[]
\CommentTok{\# {-}{-}{-} main.py pseudocode {-}{-}{-}}

\KeywordTok{def}\NormalTok{ parse\_args():}
\NormalTok{    parser }\OperatorTok{=}\NormalTok{ argparse.ArgumentParser(description}\OperatorTok{=}\StringTok{"VQE PES Explorer"}\NormalTok{)}
\NormalTok{    parser.add\_argument(}\StringTok{"{-}{-}molecule"}\NormalTok{, default}\OperatorTok{=}\StringTok{"h2"}\NormalTok{, choices}\OperatorTok{=}\NormalTok{[}\StringTok{"h2"}\NormalTok{, }\StringTok{"lih"}\NormalTok{, }\StringTok{"h2o"}\NormalTok{])}
\NormalTok{    parser.add\_argument(}\StringTok{"{-}{-}basis"}\NormalTok{, default}\OperatorTok{=}\StringTok{"sto{-}3g"}\NormalTok{)}
\NormalTok{    parser.add\_argument(}\StringTok{"{-}{-}bond{-}range"}\NormalTok{, default}\OperatorTok{=}\StringTok{"0.3,3.0,20"}\NormalTok{)}
\NormalTok{    parser.add\_argument(}\StringTok{"{-}{-}ansatz"}\NormalTok{, default}\OperatorTok{=}\StringTok{"uccsd"}\NormalTok{, choices}\OperatorTok{=}\NormalTok{[}\StringTok{"uccsd"}\NormalTok{, }\StringTok{"hardware\_efficient"}\NormalTok{])}
\NormalTok{    parser.add\_argument(}\StringTok{"{-}{-}optimizer"}\NormalTok{, default}\OperatorTok{=}\StringTok{"Adam"}\NormalTok{,}
\NormalTok{                        choices}\OperatorTok{=}\NormalTok{[}\StringTok{"Adam"}\NormalTok{, }\StringTok{"GradientDescent"}\NormalTok{, }\StringTok{"COBYLA"}\NormalTok{, }\StringTok{"SPSA"}\NormalTok{])}
\NormalTok{    parser.add\_argument(}\StringTok{"{-}{-}maxiter"}\NormalTok{, }\BuiltInTok{type}\OperatorTok{=}\BuiltInTok{int}\NormalTok{, default}\OperatorTok{=}\DecValTok{200}\NormalTok{)}
\NormalTok{    parser.add\_argument(}\StringTok{"{-}{-}hardware"}\NormalTok{, action}\OperatorTok{=}\StringTok{"store\_true"}\NormalTok{)}
\NormalTok{    parser.add\_argument(}\StringTok{"{-}{-}backend"}\NormalTok{, default}\OperatorTok{=}\StringTok{"ibm\_brisbane"}\NormalTok{)}
\NormalTok{    parser.add\_argument(}\StringTok{"{-}{-}render"}\NormalTok{, action}\OperatorTok{=}\StringTok{"store\_true"}\NormalTok{)}
\NormalTok{    parser.add\_argument(}\StringTok{"{-}{-}benchmark"}\NormalTok{, action}\OperatorTok{=}\StringTok{"store\_true"}\NormalTok{)}
\NormalTok{    parser.add\_argument(}\StringTok{"{-}{-}orbital"}\NormalTok{, }\BuiltInTok{type}\OperatorTok{=}\BuiltInTok{int}\NormalTok{, default}\OperatorTok{=}\DecValTok{0}\NormalTok{)    }\CommentTok{\# MO index to visualise}
\NormalTok{    parser.add\_argument(}\StringTok{"{-}{-}config"}\NormalTok{, default}\OperatorTok{=}\StringTok{"config.yaml"}\NormalTok{)}
    \ControlFlowTok{return}\NormalTok{ parser.parse\_args()}


\KeywordTok{def}\NormalTok{ run\_pipeline(args):}
    \CommentTok{\# 1. Parse bond range → linspace array}
\NormalTok{    start, stop, n\_pts }\OperatorTok{=}\NormalTok{ args.bond\_range.split(}\StringTok{","}\NormalTok{)}
\NormalTok{    bond\_lengths }\OperatorTok{=}\NormalTok{ np.linspace(}\BuiltInTok{float}\NormalTok{(start), }\BuiltInTok{float}\NormalTok{(stop), }\BuiltInTok{int}\NormalTok{(n\_pts))}

    \CommentTok{\# 2. Select molecule factory}
\NormalTok{    factories }\OperatorTok{=}\NormalTok{ \{}\StringTok{"h2"}\NormalTok{: create\_h2, }\StringTok{"lih"}\NormalTok{: create\_lih, }\StringTok{"h2o"}\NormalTok{: create\_h2o\}}
\NormalTok{    mol\_factory }\OperatorTok{=}\NormalTok{ factories[args.molecule]}

    \CommentTok{\# 3. PES scan}
\NormalTok{    pes\_data }\OperatorTok{=}\NormalTok{ scan\_pes(}
\NormalTok{        mol\_factory, bond\_lengths,}
\NormalTok{        basis}\OperatorTok{=}\NormalTok{args.basis, ansatz\_type}\OperatorTok{=}\NormalTok{args.ansatz,}
\NormalTok{        optimizer}\OperatorTok{=}\NormalTok{args.optimizer, maxiter}\OperatorTok{=}\NormalTok{args.maxiter,}
\NormalTok{        compute\_classical}\OperatorTok{=}\NormalTok{args.benchmark,}
\NormalTok{    )}

    \CommentTok{\# 4. Benchmark (optional)}
    \ControlFlowTok{if}\NormalTok{ args.benchmark:}
\NormalTok{        bench }\OperatorTok{=}\NormalTok{ run\_benchmark(pes\_data)}
\NormalTok{        print\_benchmark\_table(bench)}

    \CommentTok{\# 5. 3D orbital render (optional)}
    \ControlFlowTok{if}\NormalTok{ args.render:}
\NormalTok{        mol\_eq }\OperatorTok{=}\NormalTok{ mol\_factory()               }\CommentTok{\# equilibrium geometry}
\NormalTok{        hf\_eq }\OperatorTok{=}\NormalTok{ run\_hartree\_fock(mol\_eq)}
        \CommentTok{\# Convert atom coords: Å → Bohr}
\NormalTok{        atom\_coords\_bohr }\OperatorTok{=}\NormalTok{ np.array([a.position }\OperatorTok{*} \FloatTok{1.8897} \ControlFlowTok{for}\NormalTok{ a }\KeywordTok{in}\NormalTok{ mol\_eq.atoms])}
\NormalTok{        orb\_data }\OperatorTok{=}\NormalTok{ evaluate\_orbital\_on\_grid(}
\NormalTok{            hf\_eq.mo\_coefficients, args.orbital,}
\NormalTok{            atom\_coords\_bohr, [a.symbol }\ControlFlowTok{for}\NormalTok{ a }\KeywordTok{in}\NormalTok{ mol\_eq.atoms],}
\NormalTok{            mol\_eq.basis,}
\NormalTok{        )}
\NormalTok{        mesh }\OperatorTok{=}\NormalTok{ extract\_dual\_isosurface(}
\NormalTok{            orb\_data.values, isovalue}\OperatorTok{=}\FloatTok{0.05}\NormalTok{,}
\NormalTok{            orb\_data.grid\_origin, orb\_data.grid\_spacing,}
\NormalTok{        )}
\NormalTok{        window }\OperatorTok{=}\NormalTok{ init\_opengl\_context()}
\NormalTok{        vao, vbo, ebo, n\_idx }\OperatorTok{=}\NormalTok{ create\_orbital\_vao(mesh)}
\NormalTok{        camera }\OperatorTok{=}\NormalTok{ ArcballCamera()}
        \CommentTok{\# ... enter render loop with arcball interaction ...}

    \CommentTok{\# 6. Print summary}
\NormalTok{    eq\_idx }\OperatorTok{=}\NormalTok{ np.argmin(pes\_data.energies\_vqe)}
    \BuiltInTok{print}\NormalTok{(}\SpecialStringTok{f"VQE min energy: }\SpecialCharTok{\{}\NormalTok{pes\_data}\SpecialCharTok{.}\NormalTok{energies\_vqe[eq\_idx]}\SpecialCharTok{:.6f\}}\SpecialStringTok{ Ha "}
          \SpecialStringTok{f"at R = }\SpecialCharTok{\{}\NormalTok{bond\_lengths[eq\_idx]}\SpecialCharTok{:.3f\}}\SpecialStringTok{ Å"}\NormalTok{)}
\end{Highlighting}
\end{Shaded}

\emph{(Full implementation: \texttt{report2\_scripts/main.py})}

\begin{center}\rule{0.5\linewidth}{0.5pt}\end{center}

\section{Part 4 --- OpenGL Visualization
Design}\label{part-4-opengl-visualization-design}

\subsection{1. Molecular Orbital
Isosurface}\label{molecular-orbital-isosurface}

\subsubsection{VAO/VBO Layout}\label{vaovbo-layout}

Each vertex stores 7 float32 values (28 bytes):

{\def\LTcaptype{none} % do not increment counter
\begin{longtable}[]{@{}cccl@{}}
\toprule\noalign{}
Offset & Attribute & Size & Description \\
\midrule\noalign{}
\endhead
\bottomrule\noalign{}
\endlastfoot
0 & Position & 3 floats & \((x, y, z)\) in Bohr \\
12 & Normal & 3 floats & Surface normal for lighting \\
24 & Sign & 1 float & \(+1\) (positive lobe) or \(-1\) (negative
lobe) \\
\end{longtable}
}

The EBO stores triangle indices as uint32.

\subsubsection{GLSL Shaders for the
Isosurface}\label{glsl-shaders-for-the-isosurface}

\textbf{Vertex Shader (\texttt{orbital.vert}):}

The vertex shader transforms each mesh vertex from object space to clip
space via the standard \(\mathbf{P} \cdot \mathbf{V} \cdot \mathbf{M}\)
pipeline. It accepts three vertex attributes --- position \((x, y, z)\),
normal \((n_x, n_y, n_z)\), and lobe sign (\(\pm 1\)) --- bound at
locations 0, 1, and 2 respectively. The world-space position and the
transformed normal (via the inverse-transpose of the model matrix,
\(\mathbf{n}' = (\mathbf{M}^{-T})\,\mathbf{n}\)) are passed as varyings
to the fragment stage.

\emph{(Full shader: \texttt{report2\_scripts/shaders/orbital.vert})}

\textbf{Fragment Shader (\texttt{orbital.frag}):}

The fragment shader implements Phong shading with two-sided lighting
(using \(|\hat{\mathbf{n}} \cdot \hat{\mathbf{l}}|\) instead of
\(\max(0, \hat{\mathbf{n}} \cdot \hat{\mathbf{l}})\), so that back-faces
of orbital lobes are visible). The lobe sign interpolated from the
vertex stage selects the base colour: positive lobes are rendered in
blue \((0.2, 0.4, 0.9)\) and negative lobes in red \((0.9, 0.2, 0.2)\),
blended via \texttt{mix(negColor,\ posColor,\ (sign\ +\ 1)/2)}. A
specular highlight with shininess exponent 32 and strength 0.4 is added.
The final fragment is output with \(\alpha = 0.7\) (semi-transparent) to
allow visual inspection of internal structure.

\subsubsection{Mesh Update on Bond Length
Change}\label{mesh-update-on-bond-length-change}

When the user moves a bond-length slider:

\begin{enumerate}
\def\labelenumi{\arabic{enumi}.}
\tightlist
\item
  Recompute the MO at the new geometry:
  \texttt{evaluate\_orbital\_on\_grid(...)} with updated atom positions.
\item
  Re-extract the isosurface: \texttt{extract\_dual\_isosurface(...)}.
\item
  Upload new mesh data via
  \texttt{update\_orbital\_vao(vbo,\ ebo,\ new\_mesh)}.
\end{enumerate}

To avoid stalls, use \textbf{double-buffered VBOs}: allocate two VBOs
and alternate between them. When a bond-length change triggers
recomputation (which can run on a background thread), the new mesh is
uploaded to the inactive VBO. Once the upload completes, the active and
inactive VBOs are swapped, giving a seamless transition with no pipeline
stall.

\emph{(See \texttt{update\_orbital\_vao} in
\texttt{report2\_scripts/renderer.py})}

\subsection{2. Potential Energy Surface
Curve}\label{potential-energy-surface-curve}

The PES is rendered as a 2D curve in a 3D viewport, positioned to the
side of the molecular orbital. For each method (HF, CISD, VQE, FCI,
etc.), bond lengths are mapped to the \(x\)-axis and energies to the
\(y\)-axis, with \(z = 0\). Each method is drawn as a coloured
\texttt{GL\_LINE\_STRIP} with a distinct RGBA colour (grey for HF, green
for CISD, blue for VQE, gold for FCI, purple for CCSD(T)). Vertices
store \texttt{{[}x,\ y,\ z,\ r,\ g,\ b,\ a{]}} (7 floats). The VBO is
uploaded once per PES scan.

\emph{(See \texttt{create\_pes\_curve\_vao} in
\texttt{report2\_scripts/renderer.py})}

A \textbf{moving marker} (small sphere) at the current bond length
indicates which point on the PES corresponds to the displayed orbital.

\subsection{3. Molecular Geometry}\label{molecular-geometry}

Atoms are rendered as spheres (using a tessellated icosphere VAO) and
bonds as cylinders (using a tessellated tube VAO):

\begin{itemize}
\tightlist
\item
  Atom colours follow CPK convention (H=white, O=red, Li=purple,
  C=grey).
\item
  Radii are scaled van der Waals radii (\(\times 0.3\)).
\item
  Bond cylinders connect bonded atom pairs with radius \(\sim 0.1\)
  Bohr.
\item
  As the PES scan progresses, atom positions update and the molecular
  geometry animates smoothly.
\end{itemize}

\subsection{4. Camera, Lighting, and UI
Layout}\label{camera-lighting-and-ui-layout}

\textbf{Layout:} The window is divided conceptually: - Left 2/3: 3D
viewport with orbital isosurface and molecular geometry. - Right 1/3:
PES curve plot and info panel (via Dear ImGui).

\textbf{Lighting:} A single directional light from above-right, plus
ambient. The light position is fixed in camera space so it rotates with
the view.

\textbf{Camera:} Arcball camera (see \texttt{ArcballCamera} class) with:
- Left-drag: rotate - Scroll: zoom - Middle-drag: pan

\textbf{ImGui Panel:} The Dear ImGui overlay renders a side panel
displaying the molecule name, basis set, current bond length, VQE and
FCI energies, VQE--FCI error, and orbital index. A
\texttt{slider\_float} widget allows the user to drag the bond length
between 0.3 and 3.0 Å, triggering
\texttt{update\_visualization(current\_R)} whenever the value changes.

\emph{(See \texttt{report2\_scripts/renderer.py} for the full render
loop and ImGui integration.)}

\begin{center}\rule{0.5\linewidth}{0.5pt}\end{center}

\section{Part 5 --- Hardware Execution
Plan}\label{part-5-hardware-execution-plan}

\subsection{1. IBM Quantum Submission via
PennyLane}\label{ibm-quantum-submission-via-pennylane}

\subsubsection{1.1 Circuit Transpilation and Qubit
Layout}\label{circuit-transpilation-and-qubit-layout}

The PennyLane-Qiskit device handles transpilation automatically. To
control the process, pass a \texttt{transpile\_options} dictionary with
\texttt{optimization\_level} (0--3, where 3 performs the most aggressive
gate cancellation and routing), and an \texttt{initial\_layout}
specifying which physical qubits to use (e.g., 4 qubits forming a line
or star in the coupling map with minimal CNOT error). Use the
\texttt{select\_optimal\_qubits} strategy from the QAOA report.

\subsubsection{1.2 Pauli Grouping and Measurement
Circuits}\label{pauli-grouping-and-measurement-circuits}

The H₂ Hamiltonian has 15 Pauli terms. Not all can be measured
simultaneously. Commuting Pauli terms can be grouped:

\textbf{Commuting clique partitioning:} Group Paulis that mutually
commute (can be simultaneously diagonalised). For H₂:

\begin{itemize}
\tightlist
\item
  Group 1 (all Z-type):
  \(\{I, Z_0, Z_1, Z_2, Z_3, Z_0Z_1, Z_0Z_2, Z_0Z_3, Z_1Z_2, Z_1Z_3, Z_2Z_3\}\)
  --- all commute, measured in the computational basis.
\item
  Group 2: \(\{X_0X_1Y_2Y_3, Y_0Y_1X_2X_3\}\) --- commute.
\item
  Group 3: \(\{X_0Y_1Y_2X_3, Y_0X_1X_2Y_3\}\) --- commute.
\end{itemize}

This requires only 3 distinct measurement circuits instead of 15
separate ones.

PennyLane performs this grouping automatically when using
\texttt{qml.expval(H)} with a hardware device. The
\texttt{grouping\_type} parameter (e.g., \texttt{"qwc"} for qubit-wise
commuting) on \texttt{qml.Hamiltonian} controls the strategy.

\subsubsection{1.3 Shot Budget for Chemical
Accuracy}\label{shot-budget-for-chemical-accuracy}

For the H₂ Hamiltonian with \(\sum_k |c_k|^2 \approx 0.3\)
Hartree\(^2\):

\[S \geq \frac{\sum_k |c_k|^2}{\epsilon^2} = \frac{0.3}{(1.6 \times 10^{-3})^2} \approx 117{,}000 \text{ shots}\]

With Pauli grouping (3 groups), distribute shots proportionally to the
variance contribution of each group. In practice, 100,000--200,000 total
shots (across all groups) suffices for chemical accuracy on H₂.

\subsection{2. Error Mitigation for
Chemistry}\label{error-mitigation-for-chemistry}

\subsubsection{2.1 ZNE for VQE Energy}\label{zne-for-vqe-energy}

ZNE is particularly well-suited for VQE because: - The output is a
scalar (energy), and ZNE extrapolates scalar values. - The noise in VQE
circuits is dominated by incoherent errors (after Pauli twirling), which
scale approximately linearly or polynomially with noise factor.

Apply ZNE exactly as described in the QAOA report (global unitary
folding at scale factors 1, 3, 5 with Richardson extrapolation). For the
H₂ UCCSD circuit (depth \textasciitilde16 CNOTs): - \(\lambda = 1\): 16
CNOTs - \(\lambda = 3\): 48 CNOTs - \(\lambda = 5\): 80 CNOTs

At 80 CNOTs on ibm\_brisbane (\(\sim 0.5\%\) CNOT error), the total
error probability is \(\sim 33\%\), which is marginal but still within
the regime where ZNE works. For deeper circuits (LiH, H₂O), ZNE may not
suffice.

\subsubsection{2.2 Symmetry Verification}\label{symmetry-verification}

Chemistry problems have symmetries that the exact ground state obeys: -
\textbf{Particle number:} \(\langle\hat{N}\rangle = N_e\) -
\textbf{Spin:} \(\langle\hat{S}_z\rangle = 0\) (for singlet states),
\(\langle\hat{S}^2\rangle = 0\)

Post-selection: discard measurement outcomes that violate these
symmetries. This reduces the effective shot count but removes unphysical
contributions. For H₂, filter the raw bitstring samples to keep only
those whose Hamming weight equals \(N_e\) (the electron count), then
recompute the energy expectation from the filtered sample set only.

\subsubsection{2.3 Expected Performance}\label{expected-performance}

{\def\LTcaptype{none} % do not increment counter
\begin{longtable}[]{@{}
  >{\raggedright\arraybackslash}p{(\linewidth - 4\tabcolsep) * \real{0.3333}}
  >{\centering\arraybackslash}p{(\linewidth - 4\tabcolsep) * \real{0.3333}}
  >{\centering\arraybackslash}p{(\linewidth - 4\tabcolsep) * \real{0.3333}}@{}}
\toprule\noalign{}
\begin{minipage}[b]{\linewidth}\raggedright
Configuration
\end{minipage} & \begin{minipage}[b]{\linewidth}\centering
\(|E_{\text{VQE}} - E_{\text{FCI}}|\)
\end{minipage} & \begin{minipage}[b]{\linewidth}\centering
Chemical accuracy?
\end{minipage} \\
\midrule\noalign{}
\endhead
\bottomrule\noalign{}
\endlastfoot
H₂, noiseless sim & \(< 10^{-6}\) Ha & Yes \\
H₂, noisy sim (1\% depol) & \(\sim 5\) mHa & No \\
H₂, ibm\_brisbane (raw) & \(\sim 10\)--\(30\) mHa & No \\
H₂, ibm\_brisbane + ZNE & \(\sim 2\)--\(5\) mHa & Marginal \\
H₂, ibm\_brisbane + ZNE + symmetry & \(\sim 1\)--\(3\) mHa & Possible \\
\end{longtable}
}

\subsection{3. Comparison Checklist}\label{comparison-checklist}

{\def\LTcaptype{none} % do not increment counter
\begin{longtable}[]{@{}lcccc@{}}
\toprule\noalign{}
Metric & Sim (noiseless) & Sim (noisy) & Hardware & HW + ZNE \\
\midrule\noalign{}
\endhead
\bottomrule\noalign{}
\endlastfoot
\(E_{\text{VQE}}\) at \(R_e\) & ✓ & ✓ & ✓ & ✓ \\
\(|E_{\text{VQE}} - E_{\text{FCI}}|\) & ✓ & ✓ & ✓ & ✓ \\
Correlation recovery ratio & ✓ & ✓ & ✓ & ✓ \\
PES shape (qualitative) & ✓ & ✓ & ✓ & ✓ \\
Dissociation energy \(D_e\) & ✓ & ✓ & ✓ & ✓ \\
Equilibrium \(R_e\) & ✓ & ✓ & ✓ & ✓ \\
VQE iterations to converge & ✓ & ✓ & --- & --- \\
Total circuit evaluations & ✓ & ✓ & ✓ & ✓ \\
\end{longtable}
}

\textbf{Plots to produce:} 1. PES: \(E(R)\) for HF, CISD, CCSD(T), VQE,
FCI on same axes. 2. Error: \(|E_{\text{method}} - E_{\text{FCI}}|\) vs
\(R\) for each method. 3. Correlation recovery ratio vs \(R\). 4. VQE
convergence: energy vs iteration at \(R_e\) and at \(R = 3.0\) Å.

\begin{center}\rule{0.5\linewidth}{0.5pt}\end{center}

\section{Part 6 --- Benchmarking \&
Validation}\label{part-6-benchmarking-validation}

\subsection{1. Validation Against PySCF
FCI}\label{validation-against-pyscf-fci}

To validate VQE energies, a helper function builds a PySCF \texttt{Mole}
from the \texttt{MoleculeSpec}, runs RHF + FCI to obtain the exact
ground-state energy \(E_{\text{FCI}}\), and computes the absolute error
\(\Delta E = |E_{\text{VQE}} - E_{\text{FCI}}|\). The error is reported
in Hartree, milliHartree (\(\times 1000\)), and kcal/mol
(\(\times 627.509\)). Chemical accuracy
(\(\Delta E < 1.6 \times 10^{-3}\) Ha) is flagged with a boolean.

\emph{(Full implementation: \texttt{report2\_scripts/benchmarker.py} and
\texttt{report2\_scripts/classical\_hf.py})}

\subsection{2. Metrics to Measure and
Plot}\label{metrics-to-measure-and-plot}

\textbf{Correlation energy recovered:}

\[r_{\text{corr}} = \frac{E_{\text{VQE}} - E_{\text{HF}}}{E_{\text{FCI}} - E_{\text{HF}}}\]

A value of \(r_{\text{corr}} = 1.0\) means VQE recovers all correlation
energy.

\textbf{State fidelity:} If the VQE state \(|\psi_{\text{VQE}}\rangle\)
and exact ground state \(|\psi_0\rangle\) are both available (in
simulation):

\[F = |\langle\psi_{\text{VQE}}|\psi_0\rangle|^2\]

In PennyLane, extract the statevector via \texttt{qml.state()} and
compute overlap with the FCI eigenvector.

\textbf{Circuit depth vs accuracy:}

{\def\LTcaptype{none} % do not increment counter
\begin{longtable}[]{@{}ccc@{}}
\toprule\noalign{}
UCCSD layers & CNOT depth (H₂) & \(|E - E_{\text{FCI}}|\) \\
\midrule\noalign{}
\endhead
\bottomrule\noalign{}
\endlastfoot
1 Trotter step & \textasciitilde16 & \(< 10^{-6}\) \\
2 Trotter steps & \textasciitilde32 & \(< 10^{-8}\) \\
\end{longtable}
}

For H₂, a single Trotter step of UCCSD is exact (the ansatz is
expressible). For LiH and H₂O, multiple Trotter steps improve accuracy.

\subsection{3. Results Table Template}\label{results-table-template}

{\def\LTcaptype{none} % do not increment counter
\begin{longtable}[]{@{}
  >{\centering\arraybackslash}p{(\linewidth - 12\tabcolsep) * \real{0.0889}}
  >{\centering\arraybackslash}p{(\linewidth - 12\tabcolsep) * \real{0.1333}}
  >{\centering\arraybackslash}p{(\linewidth - 12\tabcolsep) * \real{0.1444}}
  >{\centering\arraybackslash}p{(\linewidth - 12\tabcolsep) * \real{0.1667}}
  >{\centering\arraybackslash}p{(\linewidth - 12\tabcolsep) * \real{0.1222}}
  >{\centering\arraybackslash}p{(\linewidth - 12\tabcolsep) * \real{0.1222}}
  >{\centering\arraybackslash}p{(\linewidth - 12\tabcolsep) * \real{0.2222}}@{}}
\toprule\noalign{}
\begin{minipage}[b]{\linewidth}\centering
\(R\) (Å)
\end{minipage} & \begin{minipage}[b]{\linewidth}\centering
\(E_{\text{HF}}\)
\end{minipage} & \begin{minipage}[b]{\linewidth}\centering
\(E_{\text{CISD}}\)
\end{minipage} & \begin{minipage}[b]{\linewidth}\centering
\(E_{\text{CCSD(T)}}\)
\end{minipage} & \begin{minipage}[b]{\linewidth}\centering
\(E_{\text{VQE}}\)
\end{minipage} & \begin{minipage}[b]{\linewidth}\centering
\(E_{\text{FCI}}\)
\end{minipage} & \begin{minipage}[b]{\linewidth}\centering
\(|E_{\text{VQE}} - E_{\text{FCI}}|\)
\end{minipage} \\
\midrule\noalign{}
\endhead
\bottomrule\noalign{}
\endlastfoot
0.50 & & & & & & \\
0.75 & & & & & & \\
1.00 & & & & & & \\
1.50 & & & & & & \\
2.50 & & & & & & \\
\end{longtable}
}

\begin{center}\rule{0.5\linewidth}{0.5pt}\end{center}

\section{Part 7 --- Research Contribution
Angle}\label{part-7-research-contribution-angle}

\subsection{1. Three Directions Toward a Publishable
Result}\label{three-directions-toward-a-publishable-result}

\textbf{Direction 1: Adaptive Ansatz PES --- ADAPT-VQE vs UCCSD Across
Dissociation}

Implement ADAPT-VQE (Grimsley et al., 2019), which builds the ansatz one
operator at a time by selecting the operator with the largest gradient
at each step. Compare the PES from ADAPT-VQE vs fixed UCCSD across the
full dissociation curve for H₂, LiH, and BeH₂. Key questions: Does
ADAPT-VQE use fewer parameters? Is the PES smoother? Does it handle the
multi-reference region (near dissociation) better? Benchmark circuit
depth and parameter count at each bond length.

\textbf{Direction 2: Noise-Aware VQE for Robust PES}

Study how noisy VQE distorts the PES and propose corrections. At each
bond length, run VQE under varying noise strengths and measure: (a)
equilibrium bond length shift, (b) dissociation energy error, (c) PES
curvature distortion (affects vibrational frequencies). Propose a
noise-aware cost function
\(\tilde{E}(\theta; \lambda) = E_{\text{noisy}}(\theta) - f(\lambda)\)
where \(f(\lambda)\) is a noise correction learned from ZNE. Show that
this gives a more accurate PES than raw noisy VQE.

\textbf{Direction 3: SHARC-VQE Partitioned Hamiltonian Benchmarking}

Implement the SHARC (Singlet-Triplet Hamiltonian Arithmetic)
partitioning from Ralli et al.~(2023), which splits the molecular
Hamiltonian into fragments that can be solved independently on smaller
quantum devices. Benchmark the partitioned approach vs monolithic UCCSD
for H₂O and small organic molecules (methane, ethylene). Measure the
accuracy-vs-resource tradeoff and identify the optimal partitioning
strategy for NISQ hardware.

\subsection{2. Most Relevant 2024--2025
Papers}\label{most-relevant-20242025-papers}

\begin{enumerate}
\def\labelenumi{\arabic{enumi}.}
\item
  \textbf{Grimsley et al.~(2024).} ``ADAPT-VQE is insensitive to rough
  parameter landscapes and barren plateaus.'' \emph{npj Quantum
  Information}. --- Robustness of adaptive ansätze.
\item
  \textbf{Motta et al.~(2024).} ``Bridging quantum computing and quantum
  chemistry: Applications of VQE to molecular systems.'' \emph{WIREs
  Computational Molecular Science}. --- Comprehensive VQE review.
\item
  \textbf{Ralli et al.~(2023--2024).} ``SHARC-VQE: Simplified
  Hamiltonian Approach with Resource-efficient Circuits.''
  \emph{Physical Review Research}. --- Hamiltonian partitioning for VQE.
\item
  \textbf{Gonthier et al.~(2024).} ``Identifying challenges towards
  practical quantum advantage through resource estimation.''
  \emph{Communications Physics}. --- What it takes for quantum advantage
  in chemistry.
\item
  \textbf{Lee et al.~(2024).} ``Is there evidence for exponential
  quantum advantage in quantum chemistry?'' \emph{Nature
  Communications}. --- Critical analysis of quantum chemistry advantage
  claims.
\item
  \textbf{Cerezo et al.~(2024).} ``Variational quantum algorithms:
  fundamental limits, robustness, and practical implementations.''
  \emph{Reviews of Modern Physics}. --- Definitive VQA review.
\item
  \textbf{Tilly et al.~(2022, updated 2024).} ``The Variational Quantum
  Eigensolver: A review of methods and best practices.'' \emph{Physics
  Reports}. --- Canonical VQE reference.
\end{enumerate}

\subsection{3. Tractable Molecules on Today's
Hardware}\label{tractable-molecules-on-todays-hardware}

{\def\LTcaptype{none} % do not increment counter
\begin{longtable}[]{@{}
  >{\raggedright\arraybackslash}p{(\linewidth - 8\tabcolsep) * \real{0.3333}}
  >{\centering\arraybackslash}p{(\linewidth - 8\tabcolsep) * \real{0.1667}}
  >{\centering\arraybackslash}p{(\linewidth - 8\tabcolsep) * \real{0.1667}}
  >{\centering\arraybackslash}p{(\linewidth - 8\tabcolsep) * \real{0.1667}}
  >{\centering\arraybackslash}p{(\linewidth - 8\tabcolsep) * \real{0.1667}}@{}}
\toprule\noalign{}
\begin{minipage}[b]{\linewidth}\raggedright
Molecule
\end{minipage} & \begin{minipage}[b]{\linewidth}\centering
Qubits (STO-3G)
\end{minipage} & \begin{minipage}[b]{\linewidth}\centering
Qubits (6-31G*)
\end{minipage} & \begin{minipage}[b]{\linewidth}\centering
Active Space
\end{minipage} & \begin{minipage}[b]{\linewidth}\centering
Tractable?
\end{minipage} \\
\midrule\noalign{}
\endhead
\bottomrule\noalign{}
\endlastfoot
H₂ & 4 (2 tapered) & 8 & Full & Yes (demonstrated) \\
LiH & 12 (10 tapered) & 22 & (2e, 5o) → 10 qubits & Yes with active
space \\
BeH₂ & 14 & 26 & (4e, 6o) → 12 qubits & Marginal \\
H₂O & 14 (12 tapered) & 38 & (4e, 4o) → 8 qubits & Yes with active
space \\
NH₃ & 16 & 40 & (6e, 6o) → 12 qubits & Marginal \\
N₂ & 20 & 48 & (6e, 6o) → 12 qubits & Marginal (strong corr.) \\
\end{longtable}
}

\textbf{Why active space is essential:} The full orbital space is far
too large for current hardware. Active space methods freeze core
electrons and discard high-energy virtual orbitals, reducing the qubit
count to a tractable number. The active space must be carefully chosen
to include the orbitals involved in bond-breaking (the valence space).

\textbf{N₂ is the holy grail:} The triple bond in N₂ involves strong
static correlation across 6 electrons in 6 orbitals. Accurately
computing the N₂ PES is a long-standing challenge for both classical and
quantum methods. Demonstrating VQE on N₂ with an active space of (6e,
6o) on 12 qubits would be a significant milestone.

\end{document}
